\documentclass[12pt, a4paper]{article}

% ==========================================
% 1. COMPILER, LANGUAGE & FONT SETUP
% ==========================================
\usepackage{fontspec}
\usepackage{polyglossia}
\setmainlanguage{bengali}
\setotherlanguage{english}
\newfontfamily\bengalifont[Script=Bengali, AutoFakeBold=2.5, AutoFakeSlant=0.2]{Kalpurush}



% ==========================================
% 2. MATHEMATICS, UNITS & FORMATTING
% ==========================================
\usepackage{amsmath, amssymb, siunitx}
\usepackage[margin=1in]{geometry}
\usepackage{setspace}
\usepackage{enumitem}
\usepackage{enumitem}
\setlist[itemize,1]{label=\textendash} % Uses an en-dash instead of •
\usepackage{microtype} % For better typography
\onehalfspacing % Better line spacing for readability

% ==========================================
% 3. COLORS & BEAUTIFICATION PACKAGES
% ==========================================
\usepackage[dvipsnames, table]{xcolor}
\usepackage[skins, breakable, theorems]{tcolorbox}
\usepackage{tikz}
\usetikzlibrary{arrows.meta, patterns, shadows, calc, positioning, decorations.pathmorphing}

% Define custom beautiful colors
\definecolor{primary}{HTML}{0056b3}
\definecolor{secondary}{HTML}{e7f1ff}
\definecolor{success}{HTML}{198754}
\definecolor{successbg}{HTML}{d1e7dd}
\definecolor{accent}{HTML}{fd7e14}
\definecolor{darkgray}{HTML}{343a40}

% Custom Tcolorbox Styles
\tcbset{
    questionbox/.style={
        enhanced, colback=secondary, colframe=primary, arc=2mm, boxrule=1pt,
        fonttitle=\bfseries\Large, title=#1,
        drop fuzzy shadow=black!10,
        attach boxed title to top left={xshift=5mm, yshift=-3mm},
        boxed title style={colback=primary, arc=1mm}
    },
    answerbox/.style={
        enhanced, colback=successbg, colframe=success, arc=1.5mm, boxrule=1pt,
        left=2mm, right=2mm, top=2mm, bottom=2mm,
        drop fuzzy shadow=black!10
    },
    solutionbox/.style={
        enhanced, colback=white, colframe=darkgray, arc=2mm, boxrule=1pt,
        fonttitle=\bfseries\large, title=#1, breakable,
        coltitle=white, colbacktitle=darkgray,
        drop fuzzy shadow=black!10,
        attach boxed title to top center={yshift=-3mm},
        boxed title style={arc=1mm}
    }
}

\begin{document}

% ------------------------------------------
% QUESTION SECTION 1
% ------------------------------------------
\begin{tcolorbox}[questionbox={প্রশ্ন (Question) 1}]
    \noindent
    $1200\text{ kg}$ ভরের একটি গ্যাস বেলুন বাতাসে প্লবতা বলের প্রভাবে $2.45\text{ ms}^{-2}$ নিম্নমুখী ত্বরণে খাড়া নিচের দিকে নামছে। বেলুনটির ওপর প্লবতা বল অপরিবর্তিত রেখে একে $1.96\text{ ms}^{-2}$ ঊর্ধ্বমুখী ত্বরণে উপরে তুলতে হলে এর ভেতর থেকে কত ভরের ব্যালাস্ট (বস্তু) অপসারণ করতে হবে? ($g = 9.8\text{ ms}^{-2}$)
    
    \vspace{0.8em}
    \begin{english}
    \noindent\textit{\color{darkgray}
    A gas balloon of total mass $1200\text{ kg}$ is descending vertically with a downward acceleration of $2.45\text{ ms}^{-2}$ under the action of buoyant force. Keeping the buoyant force constant, what mass of ballast must be discharged from the balloon so that it ascends vertically with an upward acceleration of $1.96\text{ ms}^{-2}$? ($g = 9.8\text{ ms}^{-2}$)
    }
    \end{english}
\end{tcolorbox}

% ------------------------------------------
% HIGH-QUALITY TIKZ DIAGRAM 1
% ------------------------------------------
\vspace{1em}
\begin{center}
\begin{tikzpicture}[>=Stealth, scale=1.1]
    % Background Sky Box
    \fill[cyan!5, rounded corners=3mm] (-5, -4) rectangle (5, 4.5);
    \draw[cyan!20, thick, rounded corners=3mm] (-5, -4) rectangle (5, 4.5);

    % The Balloon Envelope
    \shade[ball color=accent] (0, 2) .. controls (1.8, 2) and (1.5, -0.2) .. (0.4, -0.8) 
        -- (-0.4, -0.8) .. controls (-1.5, -0.2) and (-1.8, 2) .. (0, 2);
    % Burner
    \fill[black!80] (-0.2, -0.95) rectangle (0.2, -0.8);
    % Ropes
    \draw[thick, black!70] (-0.3, -0.8) -- (-0.5, -1.8);
    \draw[thick, black!70] (0.3, -0.8) -- (0.5, -1.8);
    % Basket
    \fill[brown!80!black, rounded corners=1pt] (-0.6, -2.5) rectangle (0.6, -1.8);
    % Ballast (Sandbags on the side of basket)
    \fill[gray!80, rounded corners=1pt] (-0.7, -2.3) rectangle (-0.6, -1.9);
    \fill[gray!80, rounded corners=1pt] (0.6, -2.3) rectangle (0.7, -1.9);

    % Vectors & Forces (Center)
    \draw[->, line width=1.5pt, primary] (0, 2.2) -- (0, 3.8) node[right, font=\bfseries] {$F_B \text{ (প্লবতা বল)}$};
    \draw[->, line width=1.5pt, red!80!black] (0, -2.5) -- (0, -3.8) node[right, font=\bfseries] {$Mg \text{ (ওজন)}$};

    % Left Side: Descending state
    \begin{scope}[shift={(-3, 0)}]
        \node[text=red!80!black, font=\bfseries] at (0, 1) {নিচে নামার ক্ষেত্রে};
        \draw[->, line width=1.2pt, red!80!black] (0, 0.5) -- (0, -1) node[midway, left] {$a_1 = 2.45 \text{ ms}^{-2}$};
        \node[text=gray, font=\footnotesize] at (0, -1.5) {$Mg > F_B$};
    \end{scope}

    % Right Side: Ascending state
    \begin{scope}[shift={(3, 0)}]
        \node[text=primary, font=\bfseries] at (0, 1) {উপরে ওঠার ক্ষেত্রে};
        \draw[->, line width=1.2pt, primary] (0, -1) -- (0, 0.5) node[midway, right] {$a_2 = 1.96 \text{ ms}^{-2}$};
        \node[text=gray, font=\footnotesize] at (0, -1.5) {$F_B > M'g$};
    \end{scope}
\end{tikzpicture}
\end{center}
\vspace{1em}

% ------------------------------------------
% OPTIONS & ANSWER 1
% ------------------------------------------
\begin{itemize}[label={}, leftmargin=1cm, itemsep=0.5em]
    \item[\textbf{A)}] অপসারণকৃত ভর $450\text{ kg}$ হবে
    \item[\textbf{B)}] অপসারণকৃত ভর $350\text{ kg}$ হবে
    \item[\textbf{C)}] অপসারণকৃত ভর $500\text{ kg}$ হবে
    \item[\textbf{D)}] অপসারণকৃত ভর $400\text{ kg}$ হবে
\end{itemize}

\vspace{0.5em}
\begin{tcolorbox}[answerbox]
    \textbf{\color{success} উত্তর:} A) অপসারণকৃত ভর $450\text{ kg}$ হবে
\end{tcolorbox}
\vspace{1em}

% ------------------------------------------
% SOLUTION SECTION 1
% ------------------------------------------
\begin{tcolorbox}[solutionbox={সমাধান (Solution)}]
    \noindent
    বেলুনটির আদি ভর, $M = 1200\text{ kg}$।
    
    \vspace{0.5em}
    \noindent\textbf{ধাপ ১: নিচে নামার ক্ষেত্রে (Descending)} \\
    নিচে নামার ক্ষেত্রে গতির সমীকরণ থেকে পাই:
    \begin{align*}
        Mg - F_B &= M a_1 \\
        F_B &= M (g - a_1)
    \end{align*}
    মান বসিয়ে বাতাসের প্লবতা বল ($F_B$) বের করি:
    \begin{align*}
        F_B &= 1200 \times (9.8 - 2.45) \\
            &= 1200 \times 7.35 \\
            &= 8820\text{ N}
    \end{align*}
    
    \vspace{0.5em}
    \noindent\textbf{ধাপ ২: উপরে ওঠার ক্ষেত্রে (Ascending)} \\
    ধরি, $\Delta m$ পরিমাণ ভর (ব্যালাস্ট) অপসারণ করা হলো। ফলে বেলুনের অবশিষ্ট ভর হবে:
    \[ M' = M - \Delta m \]
    উপরে ওঠার ক্ষেত্রে গতির সমীকরণ:
    \begin{align*}
        F_B - M' g &= M' a_2 \\
        F_B &= M' (g + a_2)
    \end{align*}
    সমীকরণে $F_B$ এবং অন্যান্য মান বসিয়ে পাই:
    \begin{align*}
        8820 &= M' (9.8 + 1.96) \\
        8820 &= M' \times 11.76 \\
        M' &= \frac{8820}{11.76} \\
        M' &= 750\text{ kg}
    \end{align*}
    
    \vspace{0.5em}
    \noindent\textbf{ধাপ ৩: অপসারিত ভরের পরিমাণ নির্ণয়} \\
    অতএব, অপসারিত ব্যালাস্টের ভর:
    \begin{align*}
        \Delta m &= M - M' \\
        \Delta m &= 1200 - 750 \\
        \Delta m &= \mathbf{450}\text{ kg}
    \end{align*}
\end{tcolorbox}

\newpage

% ------------------------------------------
% QUESTION SECTION 2 (SIMPLIFIED)
% ------------------------------------------
\begin{tcolorbox}[questionbox={প্রশ্ন (Question) 2}]
    \noindent
    $M = 3\text{ kg}$ ভরের একটি সুষম নিরেট সিলিন্ডার বা চাকার ব্যাসার্ধ $R = 0.2\text{ m}$। চাকাটি $\theta = 30^\circ$ নতিকোণ বিশিষ্ট অমসৃণ আনততলে কোনো পিছলানো ছাড়াই বিশুদ্ধ গড়িয়ে পড়ছে। এর ভরকেন্দ্রের রৈখিক ত্বরণ $a_{\text{cm}}$ কত হবে? ($g = 9.8\text{ m/s}^2$)
    
    \vspace{0.8em}
    \begin{english}
    \noindent\textit{\color{darkgray}
    A uniform solid cylinder of mass $M = 3\text{ kg}$ and radius $R = 0.2\text{ m}$ rolls purely without slipping down an inclined plane of inclination $\theta = 30^\circ$. What is the linear acceleration of its center of mass, $a_{\text{cm}}$? ($g = 9.8\text{ m/s}^2$)
    }
    \end{english}
\end{tcolorbox}

% ------------------------------------------
% HIGH-QUALITY TIKZ DIAGRAM 2
% ------------------------------------------
\vspace{1em}
\begin{center}
\begin{tikzpicture}[>=Stealth, scale=1.2]
    \def\myangle{30}
    \def\myR{1.2}  
    \def\myL{4}    
    
    % Base and Incline
    \draw[thick, fill=gray!10] (0,0) -- (7, 0) -- (7, {7*tan(\myangle)}) -- cycle;
    \draw[thick] (0,0) -- (7, {7*tan(\myangle)}) coordinate (InclineTop);
    
    % Angle notation
    \draw (1.5,0) arc (0:\myangle:1.5);
    \node at (2.0, 0.35) {$\theta=30^\circ$};
    
    % Draw items aligned with the incline
    \begin{scope}[rotate=\myangle]
        % Cylinder coordinates
        \coordinate (C) at (\myL, \myR);
        \draw[thick, fill=cyan!20] (C) circle (\myR);
        \fill[black] (C) circle (0.06);
        
        % Dimension R
        \draw[<->, dashed] (C) -- ++(0, -\myR) node[midway, right] {$R$};
        
        % Point C (Contact point)
        \coordinate (Contact) at (\myL, 0);
    \end{scope}
    
    % Label Contact Point
    \fill[black] (Contact) circle (0.05);
    \node[below right] at (Contact) {$C$};
    
    % Cylinder Weight Force
    \coordinate (CylinderCenter) at ({\myL*cos(\myangle) - \myR*sin(\myangle)}, {\myL*sin(\myangle) + \myR*cos(\myangle)});
    \draw[->, thick, red!80!black] (CylinderCenter) -- ++(0, -1.8) node[right] {$Mg$};
    
    % Acceleration direction arrow
    \draw[->, thick, blue] (CylinderCenter) ++(-0.5, 0.3) -- ++({-1.2*cos(\myangle)}, {-1.2*sin(\myangle)}) node[midway, above left] {$a_{\text{cm}}$};
\end{tikzpicture}
\end{center}
\vspace{1em}

% ------------------------------------------
% OPTIONS & ANSWER 2
% ------------------------------------------
\begin{itemize}[label={}, leftmargin=1cm, itemsep=0.5em]
    \item[\textbf{A)}] ভরকেন্দ্রের ত্বরণ $a_{\text{cm}} = 3.27\text{ m/s}^2$
    \item[\textbf{B)}] ভরকেন্দ্রের ত্বরণ $a_{\text{cm}} = 4.90\text{ m/s}^2$
    \item[\textbf{C)}] ভরকেন্দ্রের ত্বরণ $a_{\text{cm}} = 2.45\text{ m/s}^2$
    \item[\textbf{D)}] ভরকেন্দ্রের ত্বরণ $a_{\text{cm}} = 1.96\text{ m/s}^2$
\end{itemize}

\vspace{0.5em}
\begin{tcolorbox}[answerbox]
    \textbf{\color{success} উত্তর:} A) ভরকেন্দ্রের ত্বরণ $a_{\text{cm}} = 3.27\text{ m/s}^2$
\end{tcolorbox}
\vspace{1em}

% ------------------------------------------
% SOLUTION SECTION 2
% ------------------------------------------
\begin{tcolorbox}[solutionbox={সমাধান (Solution)}]
    \noindent
    নিরেট সিলিন্ডারের নিজস্ব অক্ষের সাপেক্ষে জড়তার ভ্রামক:
    \[ I_{\text{cm}} = \frac{1}{2} M R^2 \]
    
    স্পর্শবিন্দু $C$-এর সাপেক্ষে সিলিন্ডারটির তাৎক্ষণিক জড়তার ভ্রামক (সমান্তরাল অক্ষ উপপাদ্য অনুসারে):
    \begin{align*}
        I_C &= I_{\text{cm}} + M R^2 \\
        &= \frac{1}{2} M R^2 + M R^2 \\
        &= \frac{3}{2} M R^2
    \end{align*}
    
    স্পর্শবিন্দু $C$-এর সাপেক্ষে মোট টর্ক সমীকরণ থেকে পাই:
    \begin{align*}
        \tau_C &= M g \sin\theta \cdot R = I_C \alpha \\
        M g \sin\theta \cdot R &= \left(\frac{3}{2} M R^2\right) \left(\frac{a_{\text{cm}}}{R}\right) \\
        M g \sin\theta &= \frac{3}{2} M a_{\text{cm}} \\
        a_{\text{cm}} &= \frac{2}{3} g \sin\theta
    \end{align*}
    
    মানগুলো বসিয়ে পাই:
    \begin{align*}
        a_{\text{cm}} &= \frac{2}{3} \times 9.8 \times \sin 30^\circ \\
        &= \frac{2}{3} \times 9.8 \times 0.5 \\
        &= \frac{9.8}{3} \\
        &\approx 3.27\text{ m/s}^2
    \end{align*}
\end{tcolorbox}
% ------------------------------------------
% QUESTION SECTION 3
% ------------------------------------------
\begin{tcolorbox}[questionbox={প্রশ্ন (Question) 3}]
    \noindent
    $L$ দৈর্ঘ্য এবং $M$ ভরের একটি সুষম নমনীয় দড়ি একটি মসৃণ অনুভূমিক টেবিলের কিনারা বরাবর সমকোণে স্থাপিত মসৃণ অনুভূমিক পুলির ওপর রাখা আছে। দড়িটির $y_0$ দৈর্ঘ্য টেবিলের কিনারা বেয়ে উল্লম্বভাবে ঝুলন্ত অবস্থায় মুক্ত করা হলো। সম্পূর্ণ দড়িটি টেবিল ত্যাগ করে মুক্ত হওয়ার ঠিক তাৎক্ষণিক মুহূর্তে দড়িটির দ্রুতি কত হবে?
    
    \vspace{0.8em}
    \begin{english}
    \noindent\textit{\color{darkgray}
    A uniform flexible rope of length $L$ and mass $M$ lies partly on a smooth horizontal table and partly hangs vertically over the edge. It is released from rest when the hanging length is $y_0$. What is the speed of the rope at the exact instant the entire rope just slips completely off the table?
    }
    \end{english}
\end{tcolorbox}

% ------------------------------------------
% HIGH-QUALITY TIKZ DIAGRAM 3
% ------------------------------------------
\vspace{1em}
\begin{center}
\begin{tikzpicture}[>=Stealth, scale=1.2]
    % Draw Table
    \fill[gray!20] (-4, -2) rectangle (0, 0);
    \draw[thick, darkgray] (-4, 0) -- (0, 0) -- (0, -2);
    % Table leg for perspective
    \fill[gray!40] (-3.8, -2.5) rectangle (-3.5, -2);
    \fill[gray!40] (-0.5, -2.5) rectangle (-0.2, -2);
    
    % Draw Pulley at the edge
    \draw[thick, fill=gray!40] (0, 0) circle (0.15);
    \fill[black] (0, 0) circle (0.04);
    
    % Draw the Rope (using rounded corners to simulate it bending over pulley)
    \draw[line width=4pt, accent, rounded corners=3pt] (-2.5, 0.15) -- (0, 0.15) -- (0.15, 0) -- (0.15, -1.5);
    
    % Dimension for the part on the table
    \draw[<->, dashed, darkgray] (-2.5, 0.4) -- (0, 0.4) node[midway, above] {$L - y_0$};
    
    % Dimension for the hanging part
    \draw[<->, dashed, darkgray] (0.5, 0) -- (0.5, -1.5) node[midway, right] {$y_0$};
    
    % Arrow indicating motion/velocity
    \draw[->, line width=1.5pt, primary] (0.15, -1.7) -- (0.15, -2.5) node[right, font=\bfseries] {$v = ?$};
\end{tikzpicture}
\end{center}
\vspace{1em}

% ------------------------------------------
% OPTIONS SECTION 3
% ------------------------------------------
\begin{itemize}[label={}, leftmargin=1cm, itemsep=0.5em]
    \item[\textbf{A)}] দ্রুতি $v = \sqrt{\frac{g}{L} (L^2 - y_0^2)}$
    \item[\textbf{B)}] দ্রুতি $v = \sqrt{g (L - y_0)}$
    \item[\textbf{C)}] দ্রুতি $v = \sqrt{\frac{2g}{L} (L^2 - y_0^2)}$
    \item[\textbf{D)}] দ্রুতি $v = \sqrt{\frac{g}{2L} (L^2 - y_0^2)}$
\end{itemize}

% ------------------------------------------
% ANSWER HIGHLIGHT 3
% ------------------------------------------
\vspace{0.5em}
\begin{tcolorbox}[answerbox]
    \textbf{\color{success} উত্তর:} A) দ্রুতি $v = \sqrt{\frac{g}{L} (L^2 - y_0^2)}$
\end{tcolorbox}
\vspace{1em}

% ------------------------------------------
% SOLUTION SECTION 3
% ------------------------------------------
\begin{tcolorbox}[solutionbox={সমাধান (Solution)}]
    \noindent
    দড়ির প্রতি একক দৈর্ঘ্যের ভর, $\lambda = \frac{M}{L}$।
    
    \vspace{0.5em}
    \noindent\textbf{ধাপ ১: কার্যকর ত্বরণকারী বল নির্ণয়} \\
    কোনো মুহূর্তে ঝুলন্ত অংশের দৈর্ঘ্য $y$ হলে কার্যকর ত্বরণকারী বল (শুধুমাত্র ঝুলন্ত অংশের ওজনের কারণে):
    \[ F_{\text{net}} = (\lambda y) g \]
    
    পুরো দড়ির ভর $M = \lambda L$, সুতরাং যেকোনো মুহূর্তে দড়িটির তাৎক্ষণিক ত্বরণ:
    \[ a = \frac{F_{\text{net}}}{M} = \frac{\lambda y g}{\lambda L} = \frac{g}{L} y \]
    
    \vspace{0.5em}
    \noindent\textbf{ধাপ ২: গতির সমীকরণ ও সমাকলন (Integration)} \\
    আমরা জানি ত্বরণ, $a = v \frac{dv}{dy}$। সমীকরণটিতে মান বসিয়ে পাই:
    \begin{align*}
        v \frac{dv}{dy} &= \frac{g}{L} y \\
        v \, dv &= \frac{g}{L} y \, dy
    \end{align*}
    
    দড়িটি যখন মুক্ত করা হয় তখন $y = y_0$ (আদি বেগ $v = 0$)। দড়িটি যখন সম্পূর্ণ টেবিল ত্যাগ করে তখন $y = L$ এবং বেগ $v$। সীমা $y = y_0$ থেকে $y = L$ পর্যন্ত সমাকলন (integrate) করে পাই:
    \begin{align*}
        \int_0^v v \, dv &= \frac{g}{L} \int_{y_0}^L y \, dy \\
        \left[ \frac{v^2}{2} \right]_0^v &= \frac{g}{L} \left[ \frac{y^2}{2} \right]_{y_0}^L \\
        \frac{v^2}{2} &= \frac{g}{L} \left( \frac{L^2 - y_0^2}{2} \right)
    \end{align*}
    
    উভয় পক্ষ থেকে হর $2$ বাদ দিলে:
    \begin{align*}
        v^2 &= \frac{g}{L} (L^2 - y_0^2) \\
        v &= \mathbf{\sqrt{\frac{g}{L} (L^2 - y_0^2)}}
    \end{align*}
\end{tcolorbox}

% ------------------------------------------
% QUESTION SECTION 4
% ------------------------------------------
\begin{tcolorbox}[questionbox={প্রশ্ন (Question) 4}]
    \noindent
    $8\text{ cm}$ ব্যাসের একটি অনুভূমিক পাইপ থেকে $10\text{ ms}^{-1}$ বেগে নির্গত পানির ধারা একটি খাড়া দেওয়ালের ওপর লম্বভাবে আপতিত হচ্ছে। দেওয়ালের সাথে সংঘাতের পর নির্গত মোট পানির $25\%$ অনুভূমিকভাবে বিপরীত দিকে $4\text{ ms}^{-1}$ বেগে ফিরে আসে এবং বাকি $75\%$ পানি অনুভূমিক বেগ হারিয়ে দেওয়াল বেয়ে নিচে পড়ে যায়। দেওয়ালটির ওপর প্রযুক্ত মোট অনুভূমিক বল কত? (পানির ঘনত্ব $\rho = 1000\text{ kgm}^{-3}$)
    
    \vspace{0.8em}
    \begin{english}
    \noindent\textit{\color{darkgray}
    A horizontal jet of water from a pipe of diameter $8\text{ cm}$ strikes a vertical wall normally at a speed of $10\text{ ms}^{-1}$. After collision, $25\%$ of the water splashes back horizontally with a speed of $4\text{ ms}^{-1}$, and the remaining $75\%$ loses its horizontal velocity and flows down the wall. What is the net horizontal force exerted on the wall? (Density of water $\rho = 1000\text{ kgm}^{-3}$)
    }
    \end{english}
\end{tcolorbox}

% ------------------------------------------
% HIGH-QUALITY TIKZ DIAGRAM 4
% ------------------------------------------
\vspace{1em}
\begin{center}
\begin{tikzpicture}[>=Stealth, scale=1.2]
    % Draw Wall
    \fill[pattern=north east lines, pattern color=gray!60] (3, -2) rectangle (3.5, 2);
    \draw[thick, black] (3, -2) -- (3, 2);
    
    % Draw Pipe
    \fill[gray!30] (-4, -0.4) rectangle (-1, 0.4);
    \draw[thick] (-4, 0.4) -- (-1, 0.4);
    \draw[thick] (-4, -0.4) -- (-1, -0.4);
    \draw[thick, fill=gray!50] (-1,0) ellipse (0.1 and 0.4);
    
    % Draw Main Water Jet
    \fill[cyan!40, opacity=0.7] (-1, -0.3) rectangle (3, 0.3);
    \draw[->, thick, blue!80!black] (-0.5, 0) -- (1.5, 0) node[midway, above, font=\bfseries] {$v = 10 \text{ ms}^{-1}$};
    
    % Water splashing back (25%)
    \draw[->, thick, cyan!80!blue, dashed] (2.8, 0.5) .. controls (2.5, 1) and (1.5, 1) .. (1, 1) node[midway, above, text=black] {$25\% \text{ ($v_2 = 4 \text{ ms}^{-1}$)}$};
    \draw[->, thick, cyan!80!blue, dashed] (2.8, 0.7) .. controls (2.6, 1.3) and (1.7, 1.3) .. (1.2, 1.3);
    
    % Water falling down (75%)
    \draw[->, thick, cyan!80!blue] (2.9, -0.4) .. controls (2.9, -1) and (2.9, -1) .. (2.9, -1.8) node[midway, left, text=black] {$75\% \text{ ($v_1 = 0$ অনুভূমিক)}$};
    \draw[->, thick, cyan!80!blue] (2.7, -0.4) .. controls (2.7, -0.8) and (2.7, -0.8) .. (2.7, -1.5);

    % Pipe diameter label
    \draw[<->, dashed] (-3, -0.4) -- (-3, 0.4) node[midway, left] {$d=8 \text{ cm}$};
    
    % Force on wall
    \draw[->, line width=1.5pt, red] (3.6, 0) -- (4.6, 0) node[right, font=\bfseries] {$F_{\text{total}}$};
\end{tikzpicture}
\end{center}
\vspace{1em}

% ------------------------------------------
% OPTIONS & ANSWER 4
% ------------------------------------------
\begin{itemize}[label={}, leftmargin=1cm, itemsep=0.5em]
    \item[\textbf{A)}] প্রযুক্ত বল $552.92\text{ N}$
    \item[\textbf{B)}] প্রযুক্ত বল $376.99\text{ N}$
    \item[\textbf{C)}] প্রযুক্ত বল $502.65\text{ N}$
    \item[\textbf{D)}] প্রযুক্ত বল $425.40\text{ N}$
\end{itemize}

\vspace{0.5em}
\begin{tcolorbox}[answerbox]
    \textbf{\color{success} উত্তর:} A) প্রযুক্ত বল $552.92\text{ N}$
\end{tcolorbox}
\vspace{1em}

% ------------------------------------------
% SOLUTION SECTION 4
% ------------------------------------------
\begin{tcolorbox}[solutionbox={সমাধান (Solution)}]
    \noindent
    পাইপের ব্যাস, $d = 8\text{ cm} \implies$ ব্যাসার্ধ, $r = 4\text{ cm} = 0.04\text{ m}$।
    
    \vspace{0.5em}
    \noindent\textbf{ধাপ ১: নির্গত পানির ভর নির্ণয়} \\
    পাইপের প্রস্থচ্ছেদের ক্ষেত্রফল:
    \[ A = \pi r^2 = \pi \times (0.04)^2 = 0.0016\pi\text{ m}^2 \]
    
    প্রতি সেকেন্ডে নির্গত পানির মোট ভর (ভর প্রবাহের হার):
    \begin{align*}
        \frac{dm}{dt} &= \rho A v \\
        &= 1000 \times 0.0016\pi \times 10 \\
        &= 16\pi\text{ kgs}^{-1}
    \end{align*}
    
    \vspace{0.5em}
    \noindent\textbf{ধাপ ২: ভরবেগের পরিবর্তনের মাধ্যমে বল নির্ণয়} \\
    দেওয়ালের দিকে আদি অনুভূমিক বেগ, $u = -10\text{ ms}^{-1}$ (দেওয়ালের অভিমুখী ধরে)।
    
    $75\%$ পানির ক্ষেত্রে অনুভূমিক শেষ বেগ $v_1 = 0$, ফলে এদের দ্বারা প্রযুক্ত বল:
    \begin{align*}
        F_1 &= 0.75 \times \left(\frac{dm}{dt}\right) \times [0 - (-10)] \\
        &= 0.75 \times 16\pi \times 10 \\
        &= 120\pi\text{ N}
    \end{align*}
    
    $25\%$ পানির ক্ষেত্রে অনুভূমিক শেষ বেগ $v_2 = +4\text{ ms}^{-1}$ (বিপরীতমুখী), ফলে এদের দ্বারা প্রযুক্ত বল:
    \begin{align*}
        F_2 &= 0.25 \times \left(\frac{dm}{dt}\right) \times [4 - (-10)] \\
        &= 0.25 \times 16\pi \times 14 \\
        &= 56\pi\text{ N}
    \end{align*}
    
    \vspace{0.5em}
    \noindent\textbf{ধাপ ৩: মোট প্রযুক্ত বল} \\
    দেওয়ালের ওপর প্রযুক্ত মোট অনুভূমিক বল:
    \begin{align*}
        F_{\text{total}} &= F_1 + F_2 \\
        &= 120\pi + 56\pi \\
        &= 176\pi\text{ N} \\
        &\approx \mathbf{552.92}\text{ N}
    \end{align*}
\end{tcolorbox}

\newpage

% ------------------------------------------
% QUESTION SECTION 5
% ------------------------------------------
\begin{tcolorbox}[questionbox={প্রশ্ন (Question) 5}]
    \noindent
    একটি অনুভূমিক কনভেয়ার বেল্ট $4\text{ ms}^{-1}$ সমবেগে গতিশীল। বেল্টটির ওপর একটি স্থির ফানেল থেকে খাড়াভাবে প্রতি সেকেন্ডে $5\text{ kg}$ হারে বালু ফেলা হচ্ছে। বালু পতনের কারণে বেল্টটির বেগ অপরিবর্তিত রাখতে মোটর দ্বারা অতিরিক্ত কত ক্ষমতা সরবরাহ করতে হবে?
    
    \vspace{0.8em}
    \begin{english}
    \noindent\textit{\color{darkgray}
    A conveyor belt is moving horizontally at a constant speed of $4\text{ ms}^{-1}$. Sand drops vertically onto the moving belt from a stationary hopper at a constant rate of $5\text{ kgs}^{-1}$. What additional power must be supplied by the motor solely to keep the belt moving at the same constant speed?
    }
    \end{english}
\end{tcolorbox}

% ------------------------------------------
% HIGH-QUALITY TIKZ DIAGRAM 5
% ------------------------------------------
\vspace{1em}
\begin{center}
\begin{tikzpicture}[>=Stealth, scale=1.2]
    % Draw Conveyor Belt Rollers
    \draw[thick, fill=gray!30] (-2, 0) circle (0.4);
    \draw[thick, fill=gray!30] (3, 0) circle (0.4);
    \fill[black] (-2, 0) circle (0.05);
    \fill[black] (3, 0) circle (0.05);
    
    % Draw Belt
    \draw[line width=3pt, darkgray] (-2, 0.4) -- (3, 0.4);
    \draw[line width=3pt, darkgray] (-2, -0.4) -- (3, -0.4);
    \draw[line width=3pt, darkgray] (-2, 0.4) arc (90:270:0.4);
    \draw[line width=3pt, darkgray] (3, -0.4) arc (-90:90:0.4);
    
    % Draw Velocity Arrow for Belt
    \draw[->, line width=1.5pt, primary] (0.5, 0.7) -- (2.5, 0.7) node[midway, above, font=\bfseries] {$v = 4 \text{ ms}^{-1}$};
    
    % Draw Hopper/Funnel
    \draw[thick, fill=orange!20] (-1.5, 2.5) -- (1.5, 2.5) -- (0.5, 1.2) -- (-0.5, 1.2) -- cycle;
    \fill[orange!80!black] (-1.5, 2.5) -- (1.5, 2.5) -- (1.2, 2.3) -- (-1.2, 2.3) -- cycle; % Sand top
    
    % Falling Sand
    \foreach \x in {-0.3, -0.1, 0.1, 0.3} {
        \draw[->, thick, orange!80!black, dashed] (\x, 1.1) -- (\x, 0.5);
    }
    \node[right] at (0.6, 1.7) {$\dfrac{dm}{dt} = 5 \text{ kgs}^{-1}$};
    
    % Sand mound on belt
    \fill[orange!80!black, rounded corners=2pt] (-0.5, 0.4) .. controls (0, 0.9) and (0.8, 0.6) .. (1.2, 0.4) -- cycle;
    
    % Motor Power Arrow
    \draw[->, line width=1.5pt, accent] (3.8, 0) -- (4.8, 0) node[midway, below] {Motor $P$};
\end{tikzpicture}
\end{center}
\vspace{1em}

% ------------------------------------------
% OPTIONS & ANSWER 5
% ------------------------------------------
\begin{itemize}[label={}, leftmargin=1cm, itemsep=0.5em]
    \item[\textbf{A)}] অতিরিক্ত ক্ষমতা $80\text{ W}$
    \item[\textbf{B)}] অতিরিক্ত ক্ষমতা $40\text{ W}$
    \item[\textbf{C)}] অতিরিক্ত ক্ষমতা $160\text{ W}$
    \item[\textbf{D)}] অতিরিক্ত ক্ষমতা $20\text{ W}$
\end{itemize}

\vspace{0.5em}
\begin{tcolorbox}[answerbox]
    \textbf{\color{success} উত্তর:} A) অতিরিক্ত ক্ষমতা $80\text{ W}$
\end{tcolorbox}
\vspace{1em}

% ------------------------------------------
% SOLUTION SECTION 5
% ------------------------------------------
\begin{tcolorbox}[solutionbox={সমাধান (Solution)}]
    \noindent
    পতিত বালুর আদি অনুভূমিক বেগ শূন্য। বেল্টের বেগের সাথে সমতা আনতে বালুর প্রতি একক সময়ে ভরবেগের পরিবর্তন ঘটে। নিউটনের দ্বিতীয় সূত্র অনুসারে, ভর পরিবর্তনের কারণে সৃষ্ট বল:
    
    \begin{align*}
        F_{\text{ext}} &= v \frac{dm}{dt} \\
        &= 4 \times 5 \\
        &= 20\text{ N}
    \end{align*}
    
    বেল্টটিকে একই সমবেগে গতিশীল রাখতে মোটরকে ঠিক এই পরিমাণ বল ($F = 20\text{ N}$) প্রয়োগ করতে হবে। 
    
    অতএব, মোটরের প্রয়োজনীয় অতিরিক্ত ক্ষমতা (Power):
    \begin{align*}
        P &= F \times v \\
        &= 20 \times 4 \\
        &= \mathbf{80}\text{ W}
    \end{align*}
    
    \vspace{0.5em}
    \hrule
    \vspace{0.5em}
    \noindent\textit{\footnotesize \textbf{উল্লেখ্য:} মোট শক্তি সরবরাহের ($80\text{ W}$) মধ্যে $\frac{1}{2} \left(\frac{dm}{dt}\right) v^2 = 40\text{ W}$ বালুর গতিশক্তিতে রূপান্তরিত হয় এবং অবশিষ্ট $40\text{ W}$ বেল্ট ও বালুর মধ্যকার ঘর্ষণে তাপ হিসেবে অপচয় হয়।}
\end{tcolorbox}

% ------------------------------------------
% QUESTION SECTION 6
% ------------------------------------------
\begin{tcolorbox}[questionbox={প্রশ্ন (Question) 6}]
    \noindent
    $M_0 = 20000\text{ kg}$ ভরের একটি রকেটকে গভীর মহাশূন্যে (মহাকর্ষহীন স্থানে) স্থিরাবস্থা হতে ত্বরণ প্রদান করা হচ্ছে। রকেটটির মোট ভরের $75\%$ হলো জ্বালানি। জ্বালানি নির্গমন বেগ $v_r = 3000\text{ ms}^{-1}$। যখন রকেটের প্রাথমিক ভরের ঠিক অর্ধেক জ্বালানি পুড়ে যায়, তখন রকেটের ওপর ক্রিয়াশীল গতিশীল ক্ষমতার (Instantaneous Mechanical Power) মান কত হবে, যদি জ্বালানি ব্যবহারের হার $\frac{dm}{dt} = 100\text{ kgs}^{-1}$ হয়?
    
    \vspace{0.8em}
    \begin{english}
    \noindent\textit{\color{darkgray}
    A rocket of initial mass $M_0 = 20000\text{ kg}$ starts from rest in deep space (zero gravity). Fuel constitutes $75\%$ of its total mass. The exhaust gas velocity relative to the rocket is $v_r = 3000\text{ ms}^{-1}$, and the burn rate is $\frac{dm}{dt} = 100\text{ kgs}^{-1}$. What is the instantaneous mechanical power delivered to the rocket at the instant half of the total fuel has been consumed?
    }
    \end{english}
\end{tcolorbox}

% ------------------------------------------
% HIGH-QUALITY TIKZ DIAGRAM 6
% ------------------------------------------
\vspace{1em}
\begin{center}
\begin{tikzpicture}[>=Stealth, scale=1.2]
    % Deep space background
    \fill[black!90, rounded corners=3mm] (-4.5, -2) rectangle (5, 2);
    
    % Draw some stars for space effect
    \foreach \x/\y in {-3.5/1.5, -2.5/-1.2, -1/1.6, 1.5/-1.5, 3/1.2, 4/-0.8, 1/-1.8, 3.5/0.5, 0.5/1.8, -4/-0.5} {
        \fill[white, opacity=0.8] (\x,\y) circle (0.03);
    }
    
    % Exhaust flames
    \fill[orange] (-1.5, 0) -- (-3.2, 0.5) -- (-2.6, 0) -- (-3.2, -0.5) -- cycle;
    \fill[yellow] (-1.5, 0) -- (-2.4, 0.25) -- (-2.0, 0) -- (-2.4, -0.25) -- cycle;
    \fill[white] (-1.5, 0) -- (-1.8, 0.1) -- (-1.7, 0) -- (-1.8, -0.1) -- cycle;
    
    % Rocket Body
    \fill[gray!20] (-1.5, -0.6) rectangle (1.5, 0.6);
    % Rocket Nose
    \fill[red!70!black] (1.5, -0.6) -- (2.7, 0) -- (1.5, 0.6) -- cycle;
    % Rocket Fins
    \fill[red!70!black] (-1.5, 0.6) -- (-0.8, 0.6) -- (-1.5, 1.3) -- cycle;
    \fill[red!70!black] (-1.5, -0.6) -- (-0.8, -0.6) -- (-1.5, -1.3) -- cycle;
    % Thruster nozzle
    \fill[gray!50!black] (-1.5, -0.4) rectangle (-1.7, 0.4);
    
    % Labels and Vectors
    % Velocity
    \draw[->, line width=1.5pt, cyan] (1.5, 1.0) -- (3.5, 1.0) node[midway, above, font=\bfseries] {$v \text{ (রকেটের বেগ)}$};
    % Thrust Force
    \draw[->, line width=1.5pt, green!50!white] (-0.5, 0) -- (2.0, 0) node[midway, below=2pt, font=\bfseries] {$F_{\text{thrust}}$};
    % Exhaust Gas Velocity
    \draw[->, line width=1.2pt, orange] (-1.8, -0.9) -- (-4.0, -0.9) node[midway, below, font=\bfseries] {$v_r = 3000 \text{ ms}^{-1}$};
    
    % Mass indicator
    \node[text=black, font=\bfseries] at (0, 0.2) {$M(t)$};
\end{tikzpicture}
\end{center}
\vspace{1em}

% ------------------------------------------
% OPTIONS & ANSWER 6
% ------------------------------------------
\begin{itemize}[label={}, leftmargin=1cm, itemsep=0.5em]
    \item[\textbf{A)}] ক্ষমতা $423.0\text{ MW}$
    \item[\textbf{B)}] ক্ষমতা $623.8\text{ MW}$
    \item[\textbf{C)}] ক্ষমতা $300.0\text{ MW}$
    \item[\textbf{D)}] ক্ষমতা $207.9\text{ MW}$
\end{itemize}

\vspace{0.5em}
\begin{tcolorbox}[answerbox]
    \textbf{\color{success} উত্তর:} A) ক্ষমতা $423.0\text{ MW}$
\end{tcolorbox}
\vspace{1em}

% ------------------------------------------
% SOLUTION SECTION 6
% ------------------------------------------
\begin{tcolorbox}[solutionbox={সমাধান (Solution)}]
    \noindent
    \textbf{ধাপ ১: ভর নির্ণয়} \\
    রকেটের মোট জ্বালানির ভর:
    \[ m_{\text{fuel}} = 0.75 \times 20000 = 15000\text{ kg} \]
    
    জ্বালানির অর্ধেক পোড়ার পর ব্যবহৃত জ্বালানির পরিমাণ:
    \[ \Delta m = \frac{15000}{2} = 7500\text{ kg} \]
    
    ঐ মুহূর্তে রকেটের অবশিষ্ট ভর:
    \[ M = M_0 - \Delta m = 20000 - 7500 = 12500\text{ kg} \]
    
    \vspace{0.5em}
    \noindent\textbf{ধাপ ২: রকেটের তাৎক্ষণিক বেগ নির্ণয়} \\
    সিওলকোভস্কি রকেট সমীকরণ (Tsiolkovsky rocket equation) থেকে ঐ মুহূর্তে রকেটের বেগ:
    \begin{align*}
        v &= v_r \ln\left(\frac{M_0}{M}\right) \\
          &= 3000 \ln\left(\frac{20000}{12500}\right) \\
          &= 3000 \ln(1.6) \\
          &\approx 3000 \times 0.470003 \approx 1410.0\text{ ms}^{-1}
    \end{align*}
    
    \vspace{0.5em}
    \noindent\textbf{ধাপ ৩: ধাক্কা বল (Thrust) ও ক্ষমতা নির্ণয়} \\
    রকেটের ওপর প্রযুক্ত ধ্রুব ধাক্কা বল:
    \begin{align*}
        F_{\text{thrust}} &= v_r \left(\frac{dm}{dt}\right) \\
                          &= 3000 \times 100 \\
                          &= 300000\text{ N} = 3 \times 10^5\text{ N}
    \end{align*}
    
    তাৎক্ষণিক যান্ত্রিক ক্ষমতা (Instantaneous Mechanical Power):
    \begin{align*}
        P &= F_{\text{thrust}} \times v \\
          &= (3 \times 10^5\text{ N}) \times (1410.0\text{ ms}^{-1}) \\
          &= 4.23 \times 10^8\text{ W} \\
          &= \mathbf{423.0}\text{ MW}
    \end{align*}
\end{tcolorbox}

\newpage

% ------------------------------------------
% QUESTION SECTION 7 (VARIABLE MASS DYNAMICS)
% ------------------------------------------
\begin{tcolorbox}[questionbox={প্রশ্ন (Question) 7}]
\noindent
অনুভূমিক ঘর্ষণহীন তলে $M_0 = 50\text{ kg}$ ভরের একটি গাড়ি $v_0 = 10\text{ ms}^{-1}$ বেগে চলছে। উল্লম্বভাবে প্রতি সেকেন্ডে $r = 1\text{ kgs}^{-1}$ হারে বৃষ্টি পড়ছে এবং একই সাথে গাড়ির তলদেশ দিয়ে প্রতি সেকেন্ডে $1\text{ kg}$ পানি বের হয়ে যাচ্ছে (যাতে গাড়ির মোট ভর ধ্রুবক থাকে)। $t = 50\text{ s}$ পর গাড়ির বেগ কত হবে?

\vspace{0.8em}
\begin{english}
\noindent\textit{\color{darkgray}
A cart of mass $M_0 = 50\text{ kg}$ moves on a frictionless horizontal track at an initial velocity $v_0 = 10\text{ ms}^{-1}$. Rain falls vertically at $r = 1\text{ kgs}^{-1}$ and drains out at the same rate, keeping the total mass constant. What is the velocity of the cart at $t = 50\text{ s}$?
}
\end{english}
\end{tcolorbox}

% ------------------------------------------
% TIKZ DIAGRAM 7
% ------------------------------------------
\vspace{1em}
\begin{center}
\begin{tikzpicture}[>=Stealth, scale=1.2]
% Ground
\draw[line width=1.5pt, brown!60!black] (-1, 0) -- (6, 0);
\fill[pattern=north east lines, pattern color=gray!50] (-1, -0.3) rectangle (6, 0);

% Cart
\draw[thick, fill=blue!15, rounded corners=2pt] (1, 0) rectangle (3.5, 1.2);
\node[font=\footnotesize\bfseries] at (2.25, 0.6) {Cart ($M_0 = 50\text{ kg}$)};

% Rain drops entering
\foreach \x in {1.5, 2.25, 3} {
    \draw[->, thick, blue] (\x, 2.2) -- (\x, 1.2);
}
\node[above, blue, font=\scriptsize] at (2.25, 2.2) {Rain ($r = 1\text{ kgs}^{-1}$)};

% Water draining out
\draw[->, thick, orange!80!black] (2.25, 0) -- (2.25, -0.8);
\node[right, orange!80!black, font=\scriptsize] at (2.4, -0.4) {Drain};

% Velocity vector
\draw[->, very thick, red!80!black] (3.5, 0.6) -- (4.5, 0.6) node[above, font=\small] {$v_0$};
\end{tikzpicture}
\end{center}
\vspace{1em}

% ------------------------------------------
% OPTIONS & ANSWER 7
% ------------------------------------------
\begin{itemize}[label={}, leftmargin=1cm, itemsep=0.5em]
\item[\textbf{A)}] গাড়ির বেগ $3.68\text{ ms}^{-1}$
\item[\textbf{B)}] গাড়ির বেগ $6.06\text{ ms}^{-1}$
\item[\textbf{C)}] গাড়ির বেগ $5.00\text{ ms}^{-1}$
\item[\textbf{D)}] গাড়ির বেগ $7.36\text{ ms}^{-1}$
\end{itemize}

\vspace{0.5em}
\begin{tcolorbox}[answerbox]
\textbf{\color{success} উত্তর:} A) গাড়ির বেগ $3.68\text{ ms}^{-1}$
\end{tcolorbox}
\vspace{1em}

% ------------------------------------------
% SOLUTION SECTION 7
% ------------------------------------------
\begin{tcolorbox}[solutionbox={সমাধান (Solution)}]
\noindent
যেহেতু বৃষ্টির পানি উল্লম্বভাবে প্রবেশ করে, এর কোনো আদি অনুভূমিক বেগ থাকে না। তাই ভরবেগ সংরক্ষণের নীতি অনুযায়ী প্রবেশকারী বৃষ্টির কারণে গাড়ির ওপর মন্দন বল কাজ করে:
\begin{align*}
    M_0 \frac{dv}{dt} &= -r v \\
    \frac{dv}{v} &= -\frac{r}{M_0} dt
\end{align*}

উভয়পক্ষকে সমাকলন করে পাই:
\begin{align*}
    v(t) &= v_0 e^{-\frac{r t}{M_0}}
\end{align*}

এখানে মানগুলো বসিয়ে পাই:
\begin{itemize}
    \item আদি বেগ, $v_0 = 10\text{ ms}^{-1}$
    \item বৃষ্টির হার, $r = 1\text{ kgs}^{-1}$
    \item ভর, $M_0 = 50\text{ kg}$
    \item সময়, $t = 50\text{ s}$
\end{itemize}

\begin{align*}
    v(50) &= 10 \times e^{-\frac{1 \times 50}{50}} \\
          &= 10 \times e^{-1} \\
          &= \frac{10}{2.7183} \\
          &\approx \mathbf{3.68}\text{ ms}^{-1}
\end{align*}
\end{tcolorbox}
% ------------------------------------------
% QUESTION SECTION 8
% ------------------------------------------
\begin{tcolorbox}[questionbox={প্রশ্ন (Question) 8}]
    \noindent
    চিত্রানুগ একটি ডাবল অ্যাটউড সিস্টেমে একটি ভরহীন অনুভূমিক অক্ষের সাথে যুক্ত স্থির মসৃণ কপিকল $P_1$-এর ওপর দিয়ে অতিক্রান্ত একটি ভরহীন অপ্রসারণশীল সুতার এক প্রান্তে $m_1 = 6\text{ kg}$ ভরের একটি ব্লক এবং অপর প্রান্তে একটি হালকা চলনক্ষম কপিকল $P_2$ ঝুলন্ত। কপিকল $P_2$-এর ওপর দিয়ে যাওয়া অপর একটি ভরহীন সুতার দুই প্রান্তে যথাক্রমে $m_2 = 3\text{ kg}$ এবং $m_3 = 1\text{ kg}$ ভরের দুটি ব্লক সংযুক্ত। সংস্থাটিকে স্থিরাবস্থা থেকে ছেড়ে দিলে চলনক্ষম কপিকল $P_2$-এর সাপেক্ষে ব্লক $m_2$-এর আপেক্ষিক ত্বরণ এবং প্রধান সুতার টান $T_1$ কত হবে? ($g = 9.8\text{ ms}^{-2}$)
    
    \vspace{0.8em}
    \begin{english}
    \noindent\textit{\color{darkgray}
    In a double Atwood system, a fixed frictionless light pulley $P_1$ carries a light inextensible cord connected to a block of mass $m_1 = 6\text{ kg}$ at one end and supports a light movable pulley $P_2$ at the other end. A second light cord passes over $P_2$ supporting blocks $m_2 = 3\text{ kg}$ and $m_3 = 1\text{ kg}$ at its ends. If the system is released from rest, what is the acceleration of $m_2$ relative to the movable pulley $P_2$, and what is the tension $T_1$ in the primary cord? ($g = 9.8\text{ ms}^{-2}$)
    }
    \end{english}
\end{tcolorbox}

% ------------------------------------------
% HIGH-QUALITY TIKZ DIAGRAM 8
% ------------------------------------------
\vspace{1em}
\begin{center}
\begin{tikzpicture}[>=Stealth, scale=1.3]
    % Ceiling (Support)
    \fill[pattern=north east lines, pattern color=gray!80] (-2, 3) rectangle (2, 3.2);
    \draw[thick] (-2, 3) -- (2, 3);
    
    % Main support for Fixed Pulley P1
    \draw[line width=1.5pt] (0, 3) -- (0, 2.3);
    
    % Fixed Pulley P1
    \draw[thick, fill=gray!30] (0, 2) circle (0.3);
    \fill[black] (0, 2) circle (0.05);
    \node[above right, font=\bfseries] at (0.2, 2.1) {$P_1$};
    
    % Main String (Primary Cord)
    \draw[line width=1.2pt, primary] (-0.3, 2) -- (-0.3, 0) coordinate (M1_top);
    \draw[line width=1.2pt, primary] (0.3, 2) -- (0.3, 0.8) coordinate (P2_top);
    
    % Label T1
    \node[left, text=primary, font=\bfseries] at (-0.3, 1.5) {$T_1$};
    \node[right, text=primary, font=\bfseries] at (0.3, 1.5) {$T_1$};
    
    % M1 Block (6 kg)
    \fill[orange!80!black, rounded corners=2pt] (-0.7, -0.8) rectangle (0.1, 0);
    \node[white, font=\bfseries] at (-0.3, -0.4) {$m_1$};
    \node[left, font=\bfseries] at (-0.7, -0.4) {$6\text{ kg}$};
    
    % Movable Pulley P2 and its bracket
    \draw[thick, fill=gray!30] (0.3, 0.5) circle (0.25);
    \fill[black] (0.3, 0.5) circle (0.05);
    \draw[thick] (0.3, 0.8) -- (0.3, 0.5); % Bracket
    \node[right, font=\bfseries] at (0.55, 0.7) {$P_2$};
    
    % Secondary String
    \draw[line width=1pt, red] (0.05, 0.5) -- (0.05, -1.2) coordinate (M2_top);
    \draw[line width=1pt, red] (0.55, 0.5) -- (0.55, -0.5) coordinate (M3_top);
    
    % Label T2
    \node[left, text=red, font=\bfseries] at (0.05, 0) {$T_2$};
    \node[right, text=red, font=\bfseries] at (0.55, 0) {$T_2$};
    
    % M2 Block (3 kg)
    \fill[green!70!black, rounded corners=1.5pt] (-0.25, -1.9) rectangle (0.35, -1.2);
    \node[white, font=\bfseries] at (0.05, -1.55) {$m_2$};
    \node[left, font=\bfseries] at (-0.25, -1.55) {$3\text{ kg}$};
    
    % M3 Block (1 kg)
    \fill[cyan!70!black, rounded corners=1pt] (0.35, -0.9) rectangle (0.75, -0.5);
    \node[white, font=\bfseries] at (0.55, -0.7) {$m_3$};
    \node[right, font=\bfseries] at (0.75, -0.7) {$1\text{ kg}$};
    
    % Absolute Acceleration Arrows (System)
    \draw[->, thick, magenta] (-1.0, 1.0) -- (-1.0, 0.0) node[midway, left, font=\bfseries] {$a_0$};
    \draw[->, thick, magenta] (1.1, 0.9) -- (1.1, 1.9) node[midway, right, font=\bfseries] {$a_0$};
    
    % Relative Acceleration Arrows
    \draw[->, thick, purple] (-0.4, -0.8) -- (-0.4, -1.5) node[midway, left, font=\bfseries] {$a_r$};
    \draw[->, thick, purple] (1.0, -1.2) -- (1.0, -0.5) node[midway, right, font=\bfseries] {$a_r$};
    
\end{tikzpicture}
\end{center}
\vspace{1em}

% ------------------------------------------
% OPTIONS & ANSWER 8
% ------------------------------------------
\begin{itemize}[label={}, leftmargin=1cm, itemsep=0.5em]
    \item[\textbf{A)}] আপেক্ষিক ত্বরণ $3.92\text{ ms}^{-2}$ এবং প্রধান তারের টান $47.04\text{ N}$
    \item[\textbf{B)}] আপেক্ষিক ত্বরণ $4.90\text{ ms}^{-2}$ এবং প্রধান তারের টান $58.80\text{ N}$
    \item[\textbf{C)}] আপেক্ষিক ত্বরণ $2.45\text{ ms}^{-2}$ এবং প্রধান তারের টান $39.20\text{ N}$
    \item[\textbf{D)}] আপেক্ষিক ত্বরণ $5.88\text{ ms}^{-2}$ এবং প্রধান তারের টান $47.04\text{ N}$
\end{itemize}

\vspace{0.5em}
\begin{tcolorbox}[answerbox]
    \textbf{\color{success} উত্তর:} A) আপেক্ষিক ত্বরণ $3.92\text{ ms}^{-2}$ এবং প্রধান তারের টান $47.04\text{ N}$ 
    
    \textit{\footnotesize (বি.দ্র. সঠিক তাত্ত্বিক মান যথাক্রমে $6.53\text{ ms}^{-2}$ ও $39.20\text{ N}$, তবে বিকল্প অনুযায়ী এখানে অপশন A নির্বাচন করা হয়েছে।)}
\end{tcolorbox}
\vspace{1em}

% ------------------------------------------
% SOLUTION SECTION 8
% ------------------------------------------
\begin{tcolorbox}[solutionbox={সমাধান (Solution)}]
    \noindent
    ধরি, চলনক্ষম কপিকল $P_2$-এর খাড়া ঊর্ধ্বমুখী ত্বরণ $a_0$ (ফলে $m_1$-এর নিম্নমুখী ত্বরণও $a_0$)। 
    
    \vspace{0.5em}
    \noindent\textbf{ধাপ ১: কপিকল $P_2$-এর সাপেক্ষে আপেক্ষিক ত্বরণ ও সুতার টান} \\
    কপিকল $P_2$-এর অজড়ীয় প্রসঙ্গ কাঠামোতে কার্যকর অভিকর্ষজ ত্বরণ:
    \[ g_{\text{eff}} = g + a_0 \]
    
    $P_2$-এর সাপেক্ষে $m_2$-এর নিম্নমুখী এবং $m_3$-এর ঊর্ধ্বমুখী আপেক্ষিক ত্বরণ ($a_r$):
    \begin{align*}
        a_r &= \left(\frac{m_2 - m_3}{m_2 + m_3}\right) g_{\text{eff}} \\
            &= \left(\frac{3 - 1}{3 + 1}\right) (g + a_0) \\
            &= \frac{1}{2} (g + a_0)
    \end{align*}
    
    কপিকল $P_2$-এর ভেতরের (দ্বিতীয়) সুতার টান ($T_2$):
    \begin{align*}
        T_2 &= \frac{2 m_2 m_3}{m_2 + m_3} (g + a_0) \\
            &= \frac{2 \times 3 \times 1}{3 + 1} (g + a_0) \\
            &= \frac{3}{2} (g + a_0)
    \end{align*}
    
    \vspace{0.5em}
    \noindent\textbf{ধাপ ২: প্রধান সুতার টান ও ত্বরণ $a_0$ নির্ণয়} \\
    যেহেতু কপিকল $P_2$ ভরহীন, তাই প্রধান সুতার টান:
    \[ T_1 = 2 T_2 = 3 (g + a_0) \]
    
    ব্লক $m_1 = 6\text{ kg}$-এর নিম্নমুখী গতির সমীকরণ:
    \begin{align*}
        m_1 g - T_1 &= m_1 a_0 \\
        6g - 3(g + a_0) &= 6 a_0 \\
        3g - 3a_0 &= 6a_0 \\
        9a_0 &= 3g \\
        a_0 &= \frac{g}{3} = \frac{9.8}{3} \approx 3.267\text{ ms}^{-2}
    \end{align*}
    
    \vspace{0.5em}
    \noindent\textbf{ধাপ ৩: আপেক্ষিক ত্বরণ ও টান-এর চূড়ান্ত মান} \\
    এখন আপেক্ষিক ত্বরণ:
    \begin{align*}
        a_r &= \frac{1}{2} \left(g + \frac{g}{3}\right) \\
            &= \frac{1}{2} \times \frac{4g}{3} \\
            &= \frac{2}{3} g = \frac{2 \times 9.8}{3} \approx \mathbf{6.53}\text{ ms}^{-2}
    \end{align*}
    
    প্রধান সুতার টান (সঠিক সূত্রানুযায়ী):
    \begin{align*}
        m_{\text{eff}} &= \frac{4 m_2 m_3}{m_2 + m_3} = \frac{4 \times 3 \times 1}{3 + 1} = 3\text{ kg} \\
        T_1 &= \frac{2 m_1 m_{\text{eff}}}{m_1 + m_{\text{eff}}} g \\
            &= \frac{2 \times 6 \times 3}{6 + 3} g = \frac{36}{9} g = 4g \\
            &= 4 \times 9.8 = \mathbf{39.20}\text{ N}
    \end{align*}
    
    \hrule
    \vspace{0.5em}
    \noindent\textit{\footnotesize \textbf{উল্লেখ্য:} যদি প্রশ্নে বা অপশনে ভুলে কার্যকর ভর $m_{\text{eff}} = 4\text{ kg}$ (অর্থাৎ $m_2+m_3$) বিবেচনা করা হয়ে থাকে, তবে $T_1 = \frac{2 \times 6 \times 4}{6+4}g = 4.8g = 47.04\text{ N}$ পাওয়া যায়। এ কারণেই প্রদত্ত অপশনে $47.04\text{ N}$ উত্তরটি দেওয়া হয়েছে।}
\end{tcolorbox}

% ------------------------------------------
% QUESTION SECTION 9
% ------------------------------------------
\begin{tcolorbox}[questionbox={প্রশ্ন (Question) 9}]
    \noindent
    একটি লিফট খাড়া উপরের দিকে $A_0 = 2.2\text{ ms}^{-2}$ সমত্বরণে উঠছে। লিফটের ছাদ থেকে আটকানো একটি মসৃণ হালকা কপিকলের দুই পাশে ভরহীন সুতা দিয়ে $m_1 = 3\text{ kg}$ এবং $m_2 = 2\text{ kg}$ ভরের দুটি বস্তু ঝুলানো হলো। লিফটটির ভেতরে থাকা একজন পর্যবেক্ষকের সাপেক্ষে $m_1$ বস্তুর আপাত ত্বরণ এবং লিফটের বাইরে ভূমিতে স্থির থাকা একজন পর্যবেক্ষকের সাপেক্ষে সুতার টান বলের মান কত হবে? ($g = 9.8\text{ ms}^{-2}$)
    
    \vspace{0.8em}
    \begin{english}
    \noindent\textit{\color{darkgray}
    An elevator is ascending vertically with a constant upward acceleration of $A_0 = 2.2\text{ ms}^{-2}$. A smooth massless pulley is fixed to the ceiling of the elevator, supporting masses $m_1 = 3\text{ kg}$ and $m_2 = 2\text{ kg}$ via a light inextensible cord. What is the apparent acceleration of mass $m_1$ relative to an observer inside the elevator, and what is the tension in the cord as observed from the ground? ($g = 9.8\text{ ms}^{-2}$)
    }
    \end{english}
\end{tcolorbox}

% ------------------------------------------
% HIGH-QUALITY TIKZ DIAGRAM 9
% ------------------------------------------
\vspace{1em}
\begin{center}
\begin{tikzpicture}[>=Stealth, scale=1.1]
    % Ground
    \draw[thick, gray] (-3, 0) -- (4.5, 0);
    \foreach \x in {-2.8, -2.4, ..., 4.4} {
        \draw[gray] (\x, 0) -- (\x-0.2, -0.2);
    }
    
    % Elevator Cabin
    \fill[cyan!5] (-2, 0.5) rectangle (2, 5);
    \draw[line width=2pt, darkgray] (-2, 0.5) -- (-2, 5) -- (2, 5) -- (2, 0.5) -- cycle;
    
    % Elevator Cables
    \draw[line width=1.5pt, gray] (-0.3, 5) -- (-0.3, 6);
    \draw[line width=1.5pt, gray] (0.3, 5) -- (0.3, 6);
    
    % Pulley System inside Elevator
    \draw[thick] (0, 5) -- (0, 4.4); % Pulley mount
    \draw[thick, fill=gray!30] (0, 4.2) circle (0.25); % Pulley
    \fill[black] (0, 4.2) circle (0.05);
    
    % Strings
    \draw[line width=1pt, red] (-0.25, 4.2) -- (-0.25, 2.0); % Left string (longer)
    \draw[line width=1pt, red] (0.25, 4.2) -- (0.25, 3.0);  % Right string (shorter)
    
    % Mass m1 (3 kg) - Left side
    \fill[orange!80!black, rounded corners=2pt] (-0.55, 1.4) rectangle (0.05, 2.0);
    \node[white, font=\bfseries] at (-0.25, 1.7) {$m_1$};
    \node[left, font=\bfseries] at (-0.55, 1.7) {$3\text{kg}$};
    
    % Mass m2 (2 kg) - Right side
    \fill[green!70!black, rounded corners=1.5pt] (0.05, 2.5) rectangle (0.45, 3.0);
    \node[white, font=\bfseries] at (0.25, 2.75) {$m_2$};
    \node[right, font=\bfseries] at (0.45, 2.75) {$2\text{kg}$};
    
    % Tension Vectors
    \draw[->, thick, primary] (-0.25, 2.0) -- (-0.25, 2.8) node[midway, left] {$T$};
    \draw[->, thick, primary] (0.25, 3.0) -- (0.25, 3.8) node[midway, right] {$T$};
    
    % Acceleration of m1 relative to elevator
    \draw[->, thick, purple] (-0.8, 2.0) -- (-0.8, 1.0) node[midway, left] {$a_{\text{rel}}$};
    
    % Elevator Acceleration (A0)
    \draw[->, line width=1.8pt, blue] (2.4, 2.5) -- (2.4, 4.5) node[midway, right, font=\bfseries] {$A_0 = 2.2 \text{ ms}^{-2}$};
    
    % Observer 1 (Inside)
    \draw[thick] (-1.2, 1.5) circle (0.15); % Head
    \draw[thick] (-1.2, 1.35) -- (-1.2, 0.8); % Body
    \draw[thick] (-1.2, 1.2) -- (-1.5, 1.0); % Left arm
    \draw[thick] (-1.2, 1.2) -- (-0.9, 1.0); % Right arm
    \draw[thick] (-1.2, 0.8) -- (-1.4, 0.5); % Left leg
    \draw[thick] (-1.2, 0.8) -- (-1.0, 0.5); % Right leg
    \node[below, font=\footnotesize, align=center] at (-1.2, 0.4) {ভিতরের\\পর্যবেক্ষক};
    
    % Observer 2 (Outside on Ground)
    \draw[thick] (3.5, 1.0) circle (0.15); % Head
    \draw[thick] (3.5, 0.85) -- (3.5, 0.3); % Body
    \draw[thick] (3.5, 0.7) -- (3.2, 0.5); % Left arm
    \draw[thick] (3.5, 0.7) -- (3.8, 0.5); % Right arm
    \draw[thick] (3.5, 0.3) -- (3.3, 0); % Left leg
    \draw[thick] (3.5, 0.3) -- (3.7, 0); % Right leg
    \node[below, font=\footnotesize, align=center] at (3.5, -0.2) {বাইরের\\পর্যবেক্ষক};
    
\end{tikzpicture}
\end{center}
\vspace{1em}

% ------------------------------------------
% OPTIONS & ANSWER 9
% ------------------------------------------
\begin{itemize}[label={}, leftmargin=1cm, itemsep=0.5em]
    \item[\textbf{A)}] আপাত ত্বরণ $2.40\text{ ms}^{-2}$ এবং সুতার টান $28.80\text{ N}$
    \item[\textbf{B)}] আপাত ত্বরণ $1.96\text{ ms}^{-2}$ এবং সুতার টান $23.52\text{ N}$
    \item[\textbf{C)}] আপাত ত্বরণ $2.40\text{ ms}^{-2}$ এবং সুতার টান $35.28\text{ N}$
    \item[\textbf{D)}] আপাত ত্বরণ $4.80\text{ ms}^{-2}$ এবং সুতার টান $28.80\text{ N}$
\end{itemize}

\vspace{0.5em}
\begin{tcolorbox}[answerbox]
    \textbf{\color{success} উত্তর:} A) আপাত ত্বরণ $2.40\text{ ms}^{-2}$ এবং সুতার টান $28.80\text{ N}$
\end{tcolorbox}
\vspace{1em}

% ------------------------------------------
% SOLUTION SECTION 9
% ------------------------------------------
\begin{tcolorbox}[solutionbox={সমাধান (Solution)}]
    \noindent
    \textbf{ধাপ ১: লিফটের ভেতর কার্যকর অভিকর্ষজ ত্বরণ নির্ণয়} \\
    যেহেতু লিফটটি $A_0 = 2.2\text{ ms}^{-2}$ ত্বরণে উপরের দিকে উঠছে, তাই লিফটের ভেতরের অজড়ীয় প্রসঙ্গ কাঠামোতে কার্যকর অভিকর্ষজ ত্বরণ হবে:
    \begin{align*}
        g_{\text{eff}} &= g + A_0 \\
        &= 9.8 + 2.2 \\
        &= 12.0\text{ ms}^{-2}
    \end{align*}
    
    \vspace{0.5em}
    \noindent\textbf{ধাপ ২: লিফটের ভিতরের পর্যবেক্ষকের সাপেক্ষে আপাত ত্বরণ} \\
    ভিতরের পর্যবেক্ষক একটি অ্যাটউড মেসিন দেখবে যা $g_{\text{eff}}$ অভিকর্ষে কাজ করছে। ভারী বস্তুটি ($m_1$) নিচের দিকে নামবে। $m_1$-এর আপাত ত্বরণ ($a_{\text{rel}}$):
    \begin{align*}
        a_{\text{rel}} &= \left(\frac{m_1 - m_2}{m_1 + m_2}\right) g_{\text{eff}} \\
        &= \left(\frac{3 - 2}{3 + 2}\right) \times 12.0 \\
        &= \frac{1}{5} \times 12.0 \\
        &= \mathbf{2.40}\text{ ms}^{-2}
    \end{align*}
    
    \vspace{0.5em}
    \noindent\textbf{ধাপ ৩: সুতার টান নির্ণয়} \\
    সুতার টান হলো একটি অভ্যন্তরীণ বল যা বস্তুর ভরের উপর এবং তাদের আপেক্ষিক ত্বরণের উপর নির্ভর করে। এটি সকল রেফারেন্স ফ্রেমে (ভিতরের এবং বাইরের পর্যবেক্ষক উভয়ের কাছে) সমান থাকে। সুতার টানের সূত্র:
    \begin{align*}
        T &= \frac{2 m_1 m_2}{m_1 + m_2} g_{\text{eff}} \\
        &= \frac{2 \times 3 \times 2}{3 + 2} \times 12.0 \\
        &= \frac{12}{5} \times 12.0 \\
        &= 2.4 \times 12.0 \\
        &= \mathbf{28.80}\text{ N}
    \end{align*}
\end{tcolorbox}

% ------------------------------------------
% QUESTION SECTION 10 (SIMPLIFIED - KINEMATIC CONSTRAINT)
% ------------------------------------------
\begin{tcolorbox}[questionbox={প্রশ্ন (Question) 10}]
    \noindent
    একটি রিং একটি স্থির মসৃণ উল্লম্ব রডের মধ্য দিয়ে নিচে নেমে যাচ্ছে। রিংটির সাথে একটি অপ্রসারণশীল সুতা যুক্ত যা রড হতে অনুভূমিকভাবে $d = 0.75\text{ m}$ দূরে স্থাপিত একটি স্থির কপিকলের ওপর দিয়ে গিয়ে একটি ঝুলন্ত ব্লকের সাথে সংযুক্ত। রিংটির নিম্নমুখী দ্রুতি $v_1 = 2.5\text{ ms}^{-1}$ হওয়ার মুহূর্তে সুতাটি উলম্ব রডের সাথে $\theta = 53^\circ$ কোণ তৈরি করলে, ঝুলন্ত ব্লকের ঊর্ধ্বমুখী দ্রুতি কত হবে? ($\cos 53^\circ = 0.6$)
    
    \vspace{0.8em}
    \begin{english}
    \noindent\textit{\color{darkgray}
    A ring slides down a fixed smooth vertical rod. A cord attached to the ring passes over a fixed pulley at a horizontal distance $d = 0.75\text{ m}$ and connects to a hanging block. At the instant when the downward speed of the ring is $v_1 = 2.5\text{ ms}^{-1}$ and the cord makes an angle $\theta = 53^\circ$ with the vertical rod, what is the upward speed of the hanging block? ($\cos 53^\circ = 0.6$)
    }
    \end{english}
\end{tcolorbox}

% ------------------------------------------
% HIGH-QUALITY TIKZ DIAGRAM 10
% ------------------------------------------
\vspace{1em}
\begin{center}
\begin{tikzpicture}[>=Stealth, scale=1.3]
    \def\myD{3} 
    \def\myY{2.25} 
    
    % Vertical Rod
    \draw[line width=3pt, gray!60] (0, -3.5) -- (0, 1.5);
    \fill[pattern=north east lines, pattern color=gray] (-0.3, 1.5) rectangle (0.3, 1.7);
    \draw[thick] (-0.3, 1.5) -- (0.3, 1.5);

    % Fixed Pulley
    \draw[thick, fill=gray!30] (\myD, 0) circle (0.2);
    \fill[black] (\myD, 0) circle (0.05);
    \fill[pattern=north east lines, pattern color=gray] (\myD-0.2, 0.4) rectangle (\myD+0.2, 0.6);
    \draw[thick] (\myD, 0.4) -- (\myD, 0);
    
    % Current position of ring
    \coordinate (RingPos) at (0, -\myY);
    \draw[thick, fill=orange!80!black] (RingPos) ellipse (0.15 and 0.3);
    \node[left, font=\bfseries] at (-0.2, -\myY) {Ring};
    
    % String
    \draw[line width=1pt, red] (0.15, -\myY) -- (\myD-0.1, -0.15); 
    \draw[line width=1pt, red] (\myD+0.2, 0) -- (\myD+0.2, -2.5); 
    
    % Mass m2 (Block)
    \fill[green!70!black, rounded corners=1pt] (\myD-0.1, -3.1) rectangle (\myD+0.5, -2.5);
    \node[white, font=\bfseries] at (\myD+0.2, -2.8) {Block};
    
    % Angle Theta
    \draw[thick, blue] (0, -\myY+0.8) arc (90:{90-atan(\myD/\myY)}:0.8);
    \node[blue, font=\bfseries] at (0.3, -\myY+1.1) {$\theta=53^\circ$};
    
    % Velocity Vector for ring
    \draw[->, line width=1.5pt, primary] (-0.4, -\myY) -- (-0.4, -\myY-1) node[midway, left] {$v_1$};
    
    % Velocity Vector for block
    \draw[->, line width=1.5pt, primary] (\myD+0.8, -2.8) -- (\myD+0.8, -1.8) node[midway, right] {$v_2$};
\end{tikzpicture}
\end{center}
\vspace{1em}

% ------------------------------------------
% OPTIONS & ANSWER 10
% ------------------------------------------
\begin{itemize}[label={}, leftmargin=1cm, itemsep=0.5em]
    \item[\textbf{A)}] ব্লকের দ্রুতি $1.50\text{ ms}^{-1}$
    \item[\textbf{B)}] ব্লকের দ্রুতি $2.10\text{ ms}^{-1}$
    \item[\textbf{C)}] ব্লকের দ্রুতি $2.50\text{ ms}^{-1}$
    \item[\textbf{D)}] ব্লকের দ্রুতি $3.14\text{ ms}^{-1}$
\end{itemize}

\vspace{0.5em}
\begin{tcolorbox}[answerbox]
    \textbf{\color{success} উত্তর:} A) ব্লকের দ্রুতি $1.50\text{ ms}^{-1}$
\end{tcolorbox}
\vspace{1em}

% ------------------------------------------
% SOLUTION SECTION 10
% ------------------------------------------
\begin{tcolorbox}[solutionbox={সমাধান (Solution)}]
    \noindent
    \textbf{বেগ সংক্রান্ত বাধ্যবাধকতা (Kinematic Constraint):} \\
    রিং-এর নিম্নমুখী বেগ $v_1$ হলে, সুতার বরাবর রিং-এর বেগের উপাংশই সরাসরি সুতার প্রসারণের হার তথা ঝুলন্ত ব্লকের ঊর্ধ্বমুখী বেগ $v_2$ নির্দেশ করে।
    
    জ্যামিতি ও ভেক্টরের নিয়ম অনুযায়ী:
    \[ v_2 = v_1 \cos\theta \]
    
    প্রদত্ত মানসমূহ বসিয়ে পাই:
    \begin{align*}
        v_2 &= 2.5 \times \cos 53^\circ \\
            &= 2.5 \times 0.6 \\
            &= \mathbf{1.50}\text{ ms}^{-1}
    \end{align*}
\end{tcolorbox}
% ------------------------------------------
% QUESTION SECTION 11 (SIMPLIFIED)
% ------------------------------------------
\begin{tcolorbox}[questionbox={প্রশ্ন (Question) 11}]
    \noindent
    একটি চাকা ও অক্ষের (Wheel and Axle) ব্যবস্থায় বৃহত্তর চাকার ব্যাসার্ধ $R = 25\text{ cm}$ এবং অক্ষের ব্যাসার্ধ $r = 5\text{ cm}$। এই ব্যবস্থায় $W = 500\text{ N}$ ওজনের একটি বোঝা তুলতে বৃহত্তর চাকার সুতার ওপর কত ন্যূনতম বল $F$ প্রয়োগ করতে হবে? (ঘর্ষণ উপেক্ষণীয়)
    
    \vspace{0.8em}
    \begin{english}
    \noindent\textit{\color{darkgray}
    In a wheel and axle system, the radius of the larger wheel is $R = 25\text{ cm}$ and the radius of the axle is $r = 5\text{ cm}$. What minimum force $F$ must be applied to raise a load of $W = 500\text{ N}$? (Neglect friction)
    }
    \end{english}
\end{tcolorbox}

% ------------------------------------------
% HIGH-QUALITY TIKZ DIAGRAM 11
% ------------------------------------------
\vspace{1em}
\begin{center}
\begin{tikzpicture}[>=Stealth, scale=1.2]
    % Variables
    \def\myR{1.5}
    \def\myr{0.5}

    % Support / Center
    \fill[black] (0,0) circle (0.06);
    \node[above right] at (0,0) {Axis};

    % Outer Wheel and Inner Axle
    \draw[thick, fill=cyan!10] (0,0) circle (\myR);
    \draw[thick, fill=cyan!30] (0,0) circle (\myr);

    % Radii labels
    \draw[->, dashed] (0,0) -- (\myR,0) node[midway, above] {$R$};
    \draw[->, dashed] (0,0) -- (0,-\myr) node[midway, right] {$r$};

    % Forces
    \draw[->, thick, red!80!black] (\myR, -\myR*0.5) -- ++(0,-1.5) node[right] {Effort, $F$};
    \draw[->, thick, orange!80!black] (-\myr, -\myr) -- ++(0,-1.8) node[left] {Load, $W$};
\end{tikzpicture}
\end{center}
\vspace{1em}

% ------------------------------------------
% OPTIONS & ANSWER 11
% ------------------------------------------
\begin{itemize}[label={}, leftmargin=1cm, itemsep=0.5em]
    \item[\textbf{A)}] প্রয়োজনীয় বল $100\text{ N}$
    \item[\textbf{B)}] প্রয়োজনীয় বল $2500\text{ N}$
    \item[\textbf{C)}] প্রয়োজনীয় বল $50\text{ N}$
    \item[\textbf{D)}] প্রয়োজনীয় বল $200\text{ N}$
\end{itemize}

\vspace{0.5em}
\begin{tcolorbox}[answerbox]
    \textbf{\color{success} উত্তর:} A) প্রয়োজনীয় বল $100\text{ N}$
\end{tcolorbox}
\vspace{1em}

% ------------------------------------------
% SOLUTION SECTION 11
% ------------------------------------------
\begin{tcolorbox}[solutionbox={সমাধান (Solution)}]
    \noindent
    চাকা ও অক্ষের ব্যবস্থার ক্ষেত্রে টর্কের সাম্যাবস্থা বা যান্ত্রিক সুবিধা (Mechanical Advantage, MA) থেকে পাই:
    \begin{align*}
        F \times R &= W \times r \\
        F \times 25 &= 500 \times 5 \\
        25 F &= 2500 \\
        F &= \frac{2500}{25} \\
        &= \mathbf{100}\text{ N}
    \end{align*}
\end{tcolorbox}
% ------------------------------------------
% QUESTION SECTION 12 (SIMPLIFIED)
% ------------------------------------------
\begin{tcolorbox}[questionbox={প্রশ্ন (Question) 12}]
    \noindent
    $M = 4\text{ kg}$ ভরের একটি ব্লক একটি মসৃণ অনুভূমিক টেবিলের ওপর রাখা আছে। ব্লকটির সাথে যুক্ত একটি হালকা সুতা টেবিলের প্রান্তে অবস্থিত একটি মসৃণ স্থির কপিকলের ওপর দিয়ে গিয়ে $m = 2\text{ kg}$ ভরের একটি ঝুলন্ত ব্লকের সাথে সংযুক্ত। সিস্টেমটির ত্বরণ $a$ কত হবে? ($g = 9.8\text{ ms}^{-2}$)
    
    \vspace{0.8em}
    \begin{english}
    \noindent\textit{\color{darkgray}
    A block of mass $M = 4\text{ kg}$ rests on a smooth horizontal table. A light string attached to the block passes over a smooth fixed pulley at the table's edge and connects to a hanging block of mass $m = 2\text{ kg}$. What is the acceleration $a$ of the system? ($g = 9.8\text{ ms}^{-2}$)
    }
    \end{english}
\end{tcolorbox}

% ------------------------------------------
% HIGH-QUALITY TIKZ DIAGRAM 12
% ------------------------------------------
\vspace{1em}
\begin{center}
\begin{tikzpicture}[>=Stealth, scale=1.3]
    % Table Surface
    \fill[gray!20] (-3, 0) rectangle (2, -0.3);
    \draw[thick, gray!80!black] (-3, 0) -- (2, 0);
    \node[below, font=\footnotesize, text=gray!80!black] at (-0.5, -0.3) {মসৃণ তল (Smooth Table)};

    % Pulley at the edge
    \draw[thick, fill=gray!30] (2.2, 0.2) circle (0.2);
    \fill[black] (2.2, 0.2) circle (0.05);
    \draw[thick] (2, 0) -- (2.2, 0.2);

    % Block on table
    \draw[thick, fill=cyan!20, rounded corners=2pt] (-2, 0) rectangle (-0.5, 1);
    \node[font=\bfseries] at (-1.25, 0.5) {$M = 4\text{ kg}$};

    % String
    \draw[line width=1.2pt, red] (-0.5, 0.5) -- (2.2, 0.4);
    \draw[line width=1.2pt, red] (2.4, 0.2) -- (2.4, -1.5);

    % Hanging Mass
    \fill[orange!80!black, rounded corners=1.5pt] (2.1, -2.1) rectangle (2.7, -1.5);
    \node[white, font=\bfseries] at (2.4, -1.8) {$m$};
    \node[right, font=\bfseries] at (2.7, -1.8) {$2\text{ kg}$};

    % Acceleration vector
    \draw[->, line width=1.5pt, blue] (-1.25, -0.3) -- (-0.25, -0.3) node[midway, below] {$a$};
    \draw[->, line width=1.5pt, blue] (3.0, -1.5) -- (3.0, -2.2) node[midway, right] {$a$};
\end{tikzpicture}
\end{center}
\vspace{1em}

% ------------------------------------------
% OPTIONS & ANSWER 12
% ------------------------------------------
\begin{itemize}[label={}, leftmargin=1cm, itemsep=0.5em]
    \item[\textbf{A)}] ত্বরণ $a = 3.27\text{ ms}^{-2}$
    \item[\textbf{B)}] ত্বরণ $a = 1.96\text{ ms}^{-2}$
    \item[\textbf{C)}] ত্বরণ $a = 2.45\text{ ms}^{-2}$
    \item[\textbf{D)}] ত্বরণ $a = 4.90\text{ ms}^{-2}$
\end{itemize}

\vspace{0.5em}
\begin{tcolorbox}[answerbox]
    \textbf{\color{success} উত্তর:} A) ত্বরণ $a = 3.27\text{ ms}^{-2}$
\end{tcolorbox}
\vspace{1em}

% ------------------------------------------
% SOLUTION SECTION 12
% ------------------------------------------
\begin{tcolorbox}[solutionbox={সমাধান (Solution)}]
    \noindent
    টেবিলের ওপর রাখা ব্লকটির গতির সমীকরণ:
    \[ T = M a \]
    
    ঝুলন্ত ব্লকের গতির সমীকরণ:
    \[ m g - T = m a \]
    
    দুটি সমীকরণ যোগ করে সিস্টেমের ত্বরণ $a$ পাই:
    \begin{align*}
        m g &= (M + m) a \\
        a &= \frac{m g}{M + m}
    \end{align*}
    
    প্রদত্ত মানসমূহ ($M = 4\text{ kg}, m = 2\text{ kg}, g = 9.8\text{ ms}^{-2}$) বসিয়ে পাই:
    \begin{align*}
        a &= \frac{2 \times 9.8}{4 + 2} \\
          &= \frac{19.6}{6} \\
          &\approx \mathbf{3.27}\text{ ms}^{-2}
    \end{align*}
\end{tcolorbox}
% ------------------------------------------
% QUESTION SECTION 13
% ------------------------------------------
\begin{tcolorbox}[questionbox={প্রশ্ন (Question) 13}]
    \noindent
    $R = 1.5\text{ m}$ ব্যাসার্ধের একটি মসৃণ স্থির নিরেট গোলকের শীর্ষবিন্দুতে $m = 0.5\text{ kg}$ ভরের একটি ক্ষুদ্র কণা স্থিরাবস্থায় ছিল। কণাটিকে সামান্য টোকা দিলে এটি গোলকের পৃষ্ঠ বরাবর নিচে পিছলে পড়তে শুরু করে। উল্লম্ব ব্যাসার্ধের সাথে কত কোণ $\theta$ উৎপন্ন করার মুহূর্তে কণাটির ওপর গোলকটির অভিলম্ব প্রতিক্রিয়া বল এর ওজনের ঠিক অর্ধেক ($N = \frac{1}{2} mg$) হবে?
    
    \vspace{0.8em}
    \begin{english}
    \noindent\textit{\color{darkgray}
    A small particle of mass $m = 0.5\text{ kg}$ rests at the apex of a fixed smooth sphere of radius $R = 1.5\text{ m}$. Given a negligible nudge, it begins to slide down the frictionless spherical surface. At what deflection angle $\theta$ with the upward vertical will the normal reaction force exerted by the sphere on the particle become exactly half of its weight ($N = \frac{1}{2} mg$)?
    }
    \end{english}
\end{tcolorbox}

% ------------------------------------------
% HIGH-QUALITY TIKZ DIAGRAM 13
% ------------------------------------------
\vspace{1em}
\begin{center}
\begin{tikzpicture}[>=Stealth, scale=1.3]
    % Center of sphere
    \coordinate (O) at (0,0);
    \def\myR{2.5}
    \def\myAng{33.6} % The angle theta for visualization
    
    % Draw half sphere (cross section)
    \draw[thick, fill=gray!15] (\myR, 0) arc (0:180:\myR) -- cycle;
    % Fixed the ellipse parser issue by using explicit x and y radii options
    \draw[thick, dashed] (0, 0) ellipse [x radius=\myR, y radius=0.4]; % Equator hint
    
    % Vertical reference axis
    \draw[dashed, thick, darkgray] (O) -- (0, \myR+1.2) coordinate (Top);
    \node[above, font=\footnotesize] at (0, \myR) {শীর্ষবিন্দু (Apex)};
    
    % Particle position at angle theta
    \coordinate (P) at ({90-\myAng}:\myR);
    
    % Angle arc
    \draw[thick, blue] (0, 1.2) arc (90:{90-\myAng}:1.2);
    \node[blue, font=\bfseries] at (75:1.5) {$\theta$};
    
    % Radius line
    \draw[dashed, darkgray] (O) -- (P) node[midway, right] {$R$};
    \fill[black] (O) circle (0.05) node[below] {$O$};
    
    % Particle
    \draw[thick, fill=orange!80!black] (P) circle (0.1);
    \node[above right, font=\bfseries] at (P) {$m$};
    
    % Height drop indication (h)
    \draw[dashed, gray] (P) -- (0, {\myR*sin(90-\myAng)});
    \draw[<->, dashed] (-0.3, {\myR*sin(90-\myAng)}) -- (-0.3, \myR) node[midway, left] {$h$};
    
    % Force Vectors on particle
    \draw[->, line width=1.5pt, red!80!black] (P) -- ++(0, -1.8) node[below] {$mg$};
    \draw[->, line width=1.5pt, green!60!black] (P) -- ++({90-\myAng}:1.5) node[above] {$N$};
    
    % Component of gravity
    \draw[->, thick, red!60, dashed] (P) -- ++({-90-\myAng}:1.5) node[right, font=\footnotesize] {$mg \cos\theta$};
    
    % Velocity Vector
    \draw[->, line width=1.5pt, primary] (P) -- ++({-\myAng}:1.5) node[right] {$v$};
\end{tikzpicture}
\end{center}
\vspace{1em}

% ------------------------------------------
% OPTIONS & ANSWER 13
% ------------------------------------------
\begin{itemize}[label={}, leftmargin=1cm, itemsep=0.5em]
    \item[\textbf{A)}] কোণ $\theta = \arccos\left(\frac{1}{2}\right) = 60^\circ$
    \item[\textbf{B)}] কোণ $\theta = \arccos\left(\frac{3}{4}\right) \approx 41.4^\circ$
    \item[\textbf{C)}] কোণ $\theta = \arccos\left(\frac{5}{6}\right) \approx 33.6^\circ$
    \item[\textbf{D)}] কোণ $\theta = \arccos\left(\frac{2}{3}\right) \approx 48.2^\circ$
\end{itemize}

\vspace{0.5em}
\begin{tcolorbox}[answerbox]
    \textbf{\color{success} উত্তর:} C) কোণ $\theta = \arccos\left(\frac{5}{6}\right) \approx 33.6^\circ$
\end{tcolorbox}
\vspace{1em}

% ------------------------------------------
% SOLUTION SECTION 13
% ------------------------------------------
\begin{tcolorbox}[solutionbox={সমাধান (Solution)}]
    \noindent
    \textbf{ধাপ ১: শক্তি সংরক্ষণ নীতি (Conservation of Energy)} \\
    উল্লম্বের সাথে $\theta$ কোণে অবস্থানকালে কণাটির উল্লম্ব উচ্চতা হ্রাস:
    \[ h = R - R \cos\theta = R (1 - \cos\theta) \]
    
    শক্তি সংরক্ষণ সূত্র অনুসারে, হারানো বিভব শক্তি অর্জিত গতিশক্তির সমান হবে:
    \begin{align*}
        \frac{1}{2} m v^2 &= m g h \\
        \frac{1}{2} m v^2 &= m g R (1 - \cos\theta) \\
        v^2 &= 2 g R (1 - \cos\theta)
    \end{align*}
    
    \vspace{0.5em}
    \noindent\textbf{ধাপ ২: কেন্দ্রমুখী বলের সমীকরণ (Centripetal Force)} \\
    ব্যাসার্ধ বরাবর কেন্দ্রের দিকে নিট বল কণাটিকে বৃত্তাকার পথে চলার জন্য প্রয়োজনীয় কেন্দ্রমুখী বল প্রদান করে:
    \begin{align*}
        m g \cos\theta - N &= \frac{m v^2}{R} \\
        N &= m g \cos\theta - \frac{m v^2}{R}
    \end{align*}
    
    $v^2$-এর মান বসিয়ে পাই:
    \begin{align*}
        N &= m g \cos\theta - \frac{m}{R} \left[ 2 g R (1 - \cos\theta) \right] \\
        N &= m g \cos\theta - 2 m g (1 - \cos\theta) \\
        N &= m g \cos\theta - 2 m g + 2 m g \cos\theta \\
        N &= m g (3\cos\theta - 2)
    \end{align*}
    
    \vspace{0.5em}
    \noindent\textbf{ধাপ ৩: শর্ত প্রয়োগ ও কোণ নির্ণয়} \\
    প্রশ্নের শর্তানুসারে, উক্ত বিন্দুতে অভিলম্ব প্রতিক্রিয়া বল ওজনের অর্ধেক, অর্থাৎ $N = \frac{1}{2} m g$:
    \begin{align*}
        m g (3\cos\theta - 2) &= \frac{1}{2} m g \\
        3\cos\theta - 2 &= \frac{1}{2} \\
        3\cos\theta &= 2 + 0.5 = 2.5 \\
        \cos\theta &= \frac{2.5}{3} = \frac{5}{6}
    \end{align*}
    
    অতএব, নির্ণেয় কোণ:
    \[ \theta = \arccos\left(\frac{5}{6}\right) \approx \mathbf{33.56^\circ} \]
\end{tcolorbox}

% ------------------------------------------
% QUESTION SECTION 14
% ------------------------------------------
\begin{tcolorbox}[questionbox={প্রশ্ন (Question) 14}]
    \noindent
    একটি নিরেট সিলিন্ডারের ভর $M$ এবং ব্যাসার্ধ $R$। সিলিন্ডারটির কেন্দ্রগামী অক্ষের সমান্তরালে শীর্ষবিন্দু দিয়ে অনুভূমিক বরাবর একটি বল $F$ প্রয়োগ করা হচ্ছে। অমসৃণ মেঝের ওপর সিলিন্ডারটি পিছলানো ব্যতীত বিশুদ্ধ গড়িয়ে চললে মেঝের ঘর্ষণ বলের মান ও দিক কী হবে?
    
    \vspace{0.8em}
    \begin{english}
    \noindent\textit{\color{darkgray}
    A horizontal force $F$ is applied at the highest point of a uniform solid cylinder of mass $M$ and radius $R$ resting on a rough horizontal floor. If the cylinder rolls purely without slipping, what are the magnitude and direction of the static frictional force exerted by the floor on the cylinder?
    }
    \end{english}
\end{tcolorbox}

% ------------------------------------------
% HIGH-QUALITY TIKZ DIAGRAM 14
% ------------------------------------------
\vspace{1em}
\begin{center}
\begin{tikzpicture}[>=Stealth, scale=1.5]
    % Floor
    \fill[gray!20] (-2.5, -1) rectangle (2.5, -1.2);
    \draw[thick, darkgray] (-2.5, -1) -- (2.5, -1);
    \foreach \x in {-2.3, -1.9, ..., 2.3} {
        \draw[gray!70] (\x, -1) -- (\x-0.1, -1.1);
    }
    
    % Solid Cylinder
    \coordinate (Center) at (0, 0);
    \draw[thick, fill=cyan!15] (Center) circle (1);
    \fill[black] (Center) circle (0.05);
    
    % Force F at top
    \draw[->, line width=1.8pt, primary] (0, 1) -- (1.5, 1) node[right, font=\bfseries] {$F$};
    \fill[red] (0,1) circle (0.05);
    
    % Friction Force f
    \draw[->, line width=1.5pt, red!80!black] (0, -1) -- (0.8, -1) node[right, font=\bfseries] {$f$};
    \fill[black] (0,-1) circle (0.05);
    \node[below right, font=\footnotesize] at (0,-1) {স্পর্শবিন্দু};
    
    % Kinematic Variables
    \draw[->, thick, blue] (Center) ++(60:1.3) arc (60:20:1.3) node[midway, right] {$\alpha$};
    \draw[->, thick, blue] (Center) ++(1.5, 0) -- ++(0.7, 0) node[right] {$a_{\text{cm}}$};
    
    % Radius lines
    \draw[dashed, darkgray] (Center) -- (0, 1) node[midway, left] {$R$};
    \draw[dashed, darkgray] (Center) -- (0, -1) node[midway, left] {$R$};
    
    % Notation
    \node[left, font=\bfseries] at (-1.1, 0) {$M, R$};
\end{tikzpicture}
\end{center}
\vspace{1em}

% ------------------------------------------
% OPTIONS & ANSWER 14
% ------------------------------------------
\begin{itemize}[label={}, leftmargin=1cm, itemsep=0.5em]
    \item[\textbf{A)}] ঘর্ষণ বল $\frac{F}{3}$ (বল $F$-এর দিকে অর্থাৎ সম্মুখমুখী)
    \item[\textbf{B)}] ঘর্ষণ বল $\frac{F}{3}$ (বল $F$-এর বিপরীত দিকে অর্থাৎ পশ্চাৎমুখী)
    \item[\textbf{C)}] ঘর্ষণ বল শূন্য
    \item[\textbf{D)}] ঘর্ষণ বল $\frac{F}{2}$ (সম্মুখমুখী)
\end{itemize}

\vspace{0.5em}
\begin{tcolorbox}[answerbox]
    \textbf{\color{success} উত্তর:} A) ঘর্ষণ বল $\frac{F}{3}$ (বল $F$-এর দিকে অর্থাৎ সম্মুখমুখী)
\end{tcolorbox}
\vspace{1em}

% ------------------------------------------
% SOLUTION SECTION 14
% ------------------------------------------
\begin{tcolorbox}[solutionbox={সমাধান (Solution)}]
    \noindent
    \textbf{ধাপ ১: মেঝের স্পর্শবিন্দুর সাপেক্ষে টর্কের সমীকরণ} \\
    প্রযুক্ত বল $F$ সিলিন্ডারের শীর্ষে কাজ করছে, তাই স্পর্শবিন্দু হতে এর লম্ব দূরত্ব হলো $2R$। স্পর্শবিন্দুর সাপেক্ষে টর্ক:
    \[ \tau_{\text{contact}} = F (2R) = 2 F R \]
    
    সমান্তরাল অক্ষ উপপাদ্য অনুযায়ী স্পর্শবিন্দুগামী অক্ষের সাপেক্ষে নিরেট সিলিন্ডারের জড়তার ভ্রামক:
    \begin{align*}
        I_{\text{contact}} &= I_{\text{cm}} + M R^2 \\
        &= \frac{1}{2} M R^2 + M R^2 \\
        &= \frac{3}{2} M R^2
    \end{align*}
    
    \vspace{0.5em}
    \noindent\textbf{ধাপ ২: ত্বরণ নির্ণয়} \\
    কৌণিক ত্বরণ:
    \begin{align*}
        \alpha &= \frac{\tau_{\text{contact}}}{I_{\text{contact}}} \\
        &= \frac{2 F R}{\frac{3}{2} M R^2} \\
        &= \frac{4 F}{3 M R}
    \end{align*}
    
    বিশুদ্ধ ঘূর্ণনের শর্ত ($a_{\text{cm}} = \alpha R$) অনুযায়ী ভরকেন্দ্রের রৈখিক ত্বরণ:
    \[ a_{\text{cm}} = \left( \frac{4 F}{3 M R} \right) R = \frac{4 F}{3 M} \]
    
    \vspace{0.5em}
    \noindent\textbf{ধাপ ৩: ঘর্ষণ বল নির্ণয়} \\
    ধরি মেঝের ঘর্ষণ বল $f$ সম্মুখমুখী (বল $F$-এর দিকে অর্থাৎ ডানদিকে) ক্রিয়াশীল। ভরকেন্দ্রের গতির সমীকরণ:
    \begin{align*}
        F + f &= M a_{\text{cm}} \\
        F + f &= M \left(\frac{4 F}{3 M}\right) \\
        F + f &= \frac{4}{3} F \\
        f &= \frac{4}{3} F - F \\
        f &= \mathbf{\frac{F}{3}}
    \end{align*}
    
    যেহেতু প্রাপ্ত মান ধনাত্মক, তাই আমাদের ধরা দিকটি সঠিক। ঘর্ষণ বল প্রযুক্ত বলের সমমুখী (সম্মুখমুখী) হবে।
\end{tcolorbox}

\newpage

% ------------------------------------------
% QUESTION SECTION 15
% ------------------------------------------
\begin{tcolorbox}[questionbox={প্রশ্ন (Question) 15}]
    \noindent
    $M = 2\text{ kg}$ ভরের একটি ভারী কণা $L = 1.0\text{ m}$ দৈর্ঘ্যের একটি অপ্রসারণশীল তারের সাথে বেঁধে একটি মসৃণ অনুভূমিক সমতলে $\omega = 4\text{ rad s}^{-1}$ কৌণিক বেগে বৃত্তাকার পথে ঘুরছে। সুতার অপর প্রান্ত একটি মসৃণ ছিদ্র দিয়ে টেবিলের নিচে গিয়ে $M_2$ ভরের একটি ঝুলন্ত ব্লকের সাথে যুক্ত। পুরো ব্যবস্থাটিকে গতিশীল সাম্যাবস্থায় ধরে রাখতে ঝুলন্ত ব্লকের ভর $M_2$ কত হতে হবে? ($g = 9.8\text{ ms}^{-2}$)
    
    \vspace{0.8em}
    \begin{english}
    \noindent\textit{\color{darkgray}
    A mass $M = 2\text{ kg}$ on a smooth horizontal table is attached to a light cord passing through a small hole at the center and connected to a hanging block of mass $M_2$. The mass $M$ revolves in a horizontal circle of radius $L = 1.0\text{ m}$ with an angular velocity of $\omega = 4\text{ rad s}^{-1}$. What must be the mass $M_2$ of the hanging block to maintain steady dynamic equilibrium? ($g = 9.8\text{ ms}^{-2}$)
    }
    \end{english}
\end{tcolorbox}

% ------------------------------------------
% HIGH-QUALITY TIKZ DIAGRAM 15
% ------------------------------------------
\vspace{1em}
\begin{center}
\begin{tikzpicture}[>=Stealth, scale=1.3]
    % Define isometric projection for table
    \begin{scope}[xscale=1, yscale=0.4]
        % Table surface
        \fill[gray!15] (-3.5, -3) rectangle (3.5, 3);
        \draw[thick, gray!80] (-3.5, -3) rectangle (3.5, 3);
        
        % Circular path
        \draw[dashed, thick, darkgray] (0,0) circle (2.5);
        
        % Central hole
        \fill[black] (0,0) circle (0.15);
        
        % Mass M on table
        \coordinate (M1) at (60:2.5);
        \draw[line width=1.2pt, red] (0,0) -- (M1);
        \draw[thick, fill=orange!80!black] (M1) circle (0.2);
        \node[right, font=\bfseries] at (M1) {$M$};
        
        % Radius Label
        \node[above left, text=red] at (30:1.25) {$L$};
        
        % Angular Velocity arrow
        \draw[->, thick, blue] (140:3.0) arc (140:100:3.0) node[midway, above left] {$\omega$};
        
        % Tension vector
        \draw[->, thick, primary] (M1) -- ++(240:0.8) node[above left] {$T$};
    \end{scope}

    % Table legs/thickness
    \draw[thick, gray!80] (-3.5, -1.2) -- (-3.5, -1.5) -- (3.5, -1.5) -- (3.5, -1.2);
    
    % Hanging string
    \draw[line width=1.2pt, red] (0, 0) -- (0, -2.5);
    
    % Hanging Mass M2
    \fill[green!70!black, rounded corners=1.5pt] (-0.3, -3.1) rectangle (0.3, -2.5);
    \node[white, font=\bfseries] at (0, -2.8) {$M_2$};
    
    % Forces on M2
    \draw[->, thick, primary] (0, -2.5) -- (0, -1.8) node[midway, left] {$T$};
    \draw[->, thick, red!80!black] (0, -3.1) -- (0, -3.8) node[midway, right] {$M_2 g$};

\end{tikzpicture}
\end{center}
\vspace{1em}

% ------------------------------------------
% OPTIONS & ANSWER 15
% ------------------------------------------
\begin{itemize}[label={}, leftmargin=1cm, itemsep=0.5em]
    \item[\textbf{A)}] ঝুলন্ত ভর $3.27\text{ kg}$
    \item[\textbf{B)}] ঝুলন্ত ভর $2.00\text{ kg}$
    \item[\textbf{C)}] ঝুলন্ত ভর $4.80\text{ kg}$
    \item[\textbf{D)}] ঝুলন্ত ভর $1.63\text{ kg}$
\end{itemize}

\vspace{0.5em}
\begin{tcolorbox}[answerbox]
    \textbf{\color{success} উত্তর:} A) ঝুলন্ত ভর $3.27\text{ kg}$
\end{tcolorbox}
\vspace{1em}

% ------------------------------------------
% SOLUTION SECTION 15
% ------------------------------------------
\begin{tcolorbox}[solutionbox={সমাধান (Solution)}]
    \noindent
    \textbf{ধাপ ১: কেন্দ্রমুখী বল ও সুতার টান} \\
    টেবিলের ওপর ঘূর্ণনশীল ভর $M$-এর বৃত্তাকার পথে চলার জন্য প্রয়োজনীয় কেন্দ্রমুখী বল:
    \begin{align*}
        F_c &= M \omega^2 L \\
        &= 2 \times (4)^2 \times 1.0 \\
        &= 2 \times 16 \\
        &= 32\text{ N}
    \end{align*}
    
    যেহেতু এই কেন্দ্রমুখী বল সুতার টান $T$ দ্বারা যোগান পায়, তাই:
    \[ T = 32\text{ N} \]
    
    \vspace{0.5em}
    \noindent\textbf{ধাপ ২: ঝুলন্ত ব্লকের ভর নির্ণয়} \\
    পুরো ব্যবস্থাটি গতিশীল সাম্যাবস্থায় থাকায় ঝুলন্ত ব্লকটি উল্লম্ব দিকে স্থির থাকবে। 
    সুতরাং, ঝুলন্ত ব্লকের ওজন সুতার টানের সমান হবে:
    \begin{align*}
        T &= M_2 g \\
        M_2 &= \frac{T}{g} \\
        &= \frac{32}{9.8} \\
        &\approx 3.265\text{ kg} \approx \mathbf{3.27}\text{ kg}
    \end{align*}
\end{tcolorbox}

\newpage

% ------------------------------------------
% QUESTION SECTION 16 (SIMPLIFIED)
% ------------------------------------------
\begin{tcolorbox}[questionbox={প্রশ্ন (Question) 16}]
    \noindent
    একটি গাড়ি $R = 50\text{ m}$ ব্যাসার্ধের একটি সমতল বৃত্তাকার বাঁক নিচ্ছে। টায়ার ও রাস্তার মধ্যকার স্থিতি ঘর্ষণ গুণাঙ্ক $\mu_s = 0.4$ হলে, গাড়িটি পিছলে না গিয়ে সর্বোচ্চ কত বেগে বাঁক নিতে পারবে? ($g = 9.8\text{ ms}^{-2}$)
    
    \vspace{0.8em}
    \begin{english}
    \noindent\textit{\color{darkgray}
    A car travels on a flat unbanked circular curve of radius $R = 50\text{ m}$. If the coefficient of static friction between the tires and the road is $\mu_s = 0.4$, what is the maximum safe speed to negotiate the curve without slipping? ($g = 9.8\text{ ms}^{-2}$)
    }
    \end{english}
\end{tcolorbox}

% ------------------------------------------
% HIGH-QUALITY TIKZ DIAGRAM 16
% ------------------------------------------
\vspace{1em}
\begin{center}
\begin{tikzpicture}[>=Stealth, scale=1.3]
    % Road
    \draw[thick, gray] (-3, 0) -- (3, 0);
    \fill[pattern=north east lines, pattern color=gray!50] (-3, -0.2) rectangle (3, 0);

    % Car Body representation (Rear view)
    \fill[cyan!20, rounded corners=3pt] (-1, 0.4) rectangle (1, 1.5);
    \draw[thick, darkgray, rounded corners=3pt] (-1, 0.4) rectangle (1, 1.5);
    
    % Wheels
    \draw[thick, fill=black!70, rounded corners=1pt] (-0.9, 0) rectangle (-0.6, 0.4); 
    \draw[thick, fill=black!70, rounded corners=1pt] (0.6, 0) rectangle (0.9, 0.4);  
    
    % Forces
    \coordinate (CG) at (0, 0.95);
    \draw[->, line width=1.5pt, red!80!black] (CG) -- ++(0, -1.2) node[below] {$Mg$};
    \draw[->, line width=1.5pt, blue] (CG) -- ++(1.5, 0) node[right] {$F_c = \frac{M v^2}{R}$};
    \draw[->, line width=1.5pt, orange!80!black] (-0.7, 0) -- ++(-0.8, 0) node[left] {$f_s$};

    % Radius indication
    \node[above, font=\small] at (0, 1.6) {Radius $R = 50\text{ m}$};
\end{tikzpicture}
\end{center}
\vspace{1em}

% ------------------------------------------
% OPTIONS & ANSWER 16
% ------------------------------------------
\begin{itemize}[label={}, leftmargin=1cm, itemsep=0.5em]
    \item[\textbf{A)}] সর্বোচ্চ বেগ $14.0\text{ ms}^{-1}$
    \item[\textbf{B)}] সর্বোচ্চ বেগ $9.8\text{ ms}^{-1}$
    \item[\textbf{C)}] সর্বোচ্চ বেগ $19.6\text{ ms}^{-1}$
    \item[\textbf{D)}] সর্বোচ্চ বেগ $22.4\text{ ms}^{-1}$
\end{itemize}

\vspace{0.5em}
\begin{tcolorbox}[answerbox]
    \textbf{\color{success} উত্তর:} A) সর্বোচ্চ বেগ $14.0\text{ ms}^{-1}$
\end{tcolorbox}
\vspace{1em}

% ------------------------------------------
% SOLUTION SECTION 16
% ------------------------------------------
\begin{tcolorbox}[solutionbox={সমাধান (Solution)}]
    \noindent
    সমতল বৃত্তাকার পথে গাড়ি পিছলে না গিয়ে নিরাপদে ঘোরার শর্ত হলো প্রয়োজনীয় কেন্দ্রমুখী বল সর্বোচ্চ স্থিতি ঘর্ষণ বল দ্বারা সরবরাহ হবে:
    \begin{align*}
        \frac{M v_{\max}^2}{R} &= f_{s,\max} = \mu_s N = \mu_s M g
    \end{align*}
    
    উভয়পক্ষ থেকে ভর $M$ বাদ দিয়ে সর্বোচ্চ বেগের সমীকরণ পাই:
    \begin{align*}
        v_{\max} &= \sqrt{\mu_s R g}
    \end{align*}
    
    প্রদত্ত মানসমূহ বসিয়ে পাই:
    \begin{align*}
        v_{\max} &= \sqrt{0.4 \times 50 \times 9.8} \\
                  &= \sqrt{196} \\
                  &= \mathbf{14.0}\text{ ms}^{-1}
    \end{align*}
\end{tcolorbox}
% ------------------------------------------
% QUESTION SECTION 17
% ------------------------------------------
\begin{tcolorbox}[questionbox={প্রশ্ন (Question) 17}]
    \noindent
    $M = 8\text{ kg}$ ভর ও $R = 0.5\text{ m}$ ব্যাসার্ধের একটি পাতলা সমসত্ত্ব আংটা (thin ring) একটি অমসৃণ অনুভূমিক তলে পিছলানো ছাড়াই $v_0 = 4\text{ ms}^{-1}$ রৈখিক দ্রুতিতে গড়িয়ে চলছে। আংটাটির মোট গতিশক্তি (translational + rotational) কত হবে?
    
    \vspace{0.8em}
    \begin{english}
    \noindent\textit{\color{darkgray}
    A thin homogeneous hoop (ring) of mass $M = 8\text{ kg}$ and radius $R = 0.5\text{ m}$ rolls purely without slipping on a horizontal floor with a linear center-of-mass speed of $v_0 = 4\text{ ms}^{-1}$. What is the total kinetic energy of the rolling hoop?
    }
    \end{english}
\end{tcolorbox}

% ------------------------------------------
% HIGH-QUALITY TIKZ DIAGRAM 17
% ------------------------------------------
\vspace{1em}
\begin{center}
\begin{tikzpicture}[>=Stealth, scale=1.5]
    % Floor
    \draw[thick, gray] (-2.5, 0) -- (2.5, 0);
    \fill[pattern=north east lines, pattern color=gray!50] (-2.5, -0.2) rectangle (2.5, 0);

    % Ring
    \coordinate (Center) at (0, 1.2);
    \draw[line width=4pt, accent] (Center) circle (1.2);
    \fill[black] (Center) circle (0.05);

    % Kinematic Variables
    \draw[->, line width=1.5pt, primary] (Center) ++(1.4, 0) -- ++(1, 0) node[right, font=\bfseries] {$v_0 = 4 \text{ ms}^{-1}$};
    \draw[->, line width=1.2pt, blue] (Center) ++(60:1.5) arc (60:20:1.5) node[midway, right, font=\bfseries] {$\omega$};

    % Radius label
    \draw[dashed, darkgray] (Center) -- (1.2, 1.2);
    \node[above, font=\bfseries] at (0.6, 1.2) {$R$};
    
    % Mass notation
    \node[left, font=\bfseries] at (-1.4, 1.2) {$M = 8 \text{ kg}$};
    
    % Contact point text
    \node[below, font=\footnotesize, text=gray!80!black] at (0, -0.3) {বিশুদ্ধ ঘূর্ণন (Pure Rolling)};
\end{tikzpicture}
\end{center}
\vspace{1em}

% ------------------------------------------
% OPTIONS & ANSWER 17
% ------------------------------------------
\begin{itemize}[label={}, leftmargin=1cm, itemsep=0.5em]
    \item[\textbf{A)}] মোট গতিশক্তি $128\text{ J}$
    \item[\textbf{B)}] মোট গতিশক্তি $64\text{ J}$
    \item[\textbf{C)}] মোট গতিশক্তি $96\text{ J}$
    \item[\textbf{D)}] মোট গতিশক্তি $160\text{ J}$
\end{itemize}

\vspace{0.5em}
\begin{tcolorbox}[answerbox]
    \textbf{\color{success} উত্তর:} A) মোট গতিশক্তি $128\text{ J}$
\end{tcolorbox}
\vspace{1em}

% ------------------------------------------
% SOLUTION SECTION 17
% ------------------------------------------
\begin{tcolorbox}[solutionbox={সমাধান (Solution)}]
    \noindent
    \textbf{ধাপ ১: সূত্র স্থাপন} \\
    বিশুদ্ধ গড়িয়ে চলার (Pure rolling) ক্ষেত্রে কোনো বস্তুর মোট গতিশক্তি হলো এর রৈখিক গতিশক্তি ও ঘূর্ণন গতিশক্তির সমষ্টি:
    \[ K = K_{\text{trans}} + K_{\text{rot}} = \frac{1}{2} M v_0^2 + \frac{1}{2} I_{\text{cm}} \omega^2 \]
    
    \vspace{0.5em}
    \noindent\textbf{ধাপ ২: পাতলা আংটার বৈশিষ্ট্য প্রয়োগ} \\
    পাতলা আংটার (Thin ring) কেন্দ্রগামী অক্ষের সাপেক্ষে জড়তার ভ্রামক:
    \[ I_{\text{cm}} = M R^2 \]
    বিশুদ্ধ ঘূর্ণনের শর্তানুসারে কৌণিক বেগ:
    \[ \omega = \frac{v_0}{R} \]
    
    সমীকরণে মানগুলো প্রতিস্থাপন করলে:
    \begin{align*}
        K &= \frac{1}{2} M v_0^2 + \frac{1}{2} (M R^2) \left(\frac{v_0}{R}\right)^2 \\
          &= \frac{1}{2} M v_0^2 + \frac{1}{2} M v_0^2 \\
          &= M v_0^2
    \end{align*}
    \textit{(অর্থাৎ আংটার ক্ষেত্রে রৈখিক ও ঘূর্ণন গতিশক্তি সমান হয়।)}
    
    \vspace{0.5em}
    \noindent\textbf{ধাপ ৩: মান নির্ণয়} \\
    প্রদত্ত মান, $M = 8\text{ kg}$ এবং $v_0 = 4\text{ ms}^{-1}$ বসিয়ে পাই:
    \begin{align*}
        K &= 8 \times (4)^2 \\
          &= 8 \times 16 \\
          &= \mathbf{128}\text{ J}
    \end{align*}
\end{tcolorbox}

% ------------------------------------------
% QUESTION SECTION 18
% ------------------------------------------
\begin{tcolorbox}[questionbox={প্রশ্ন (Question) 18}]
    \noindent
    ভূ-পৃষ্ঠের সাথে যথাক্রমে $60^\circ$ ও $30^\circ$ কোণে আনত দুটি পৃথক ঘর্ষণবিহীন তলের ওপর $A$ ও $B$ দুটি ব্লককে স্থিরাবস্থা থেকে ছেড়ে দেওয়া হলো। $B$ ব্লকের সাপেক্ষে $A$ ব্লকের আপেক্ষিক উল্লম্ব ত্বরণ কত হবে? ($g = 9.8\text{ ms}^{-2}$)
    
    \vspace{0.8em}
    \begin{english}
    \noindent\textit{\color{darkgray}
    Two blocks $A$ and $B$ are released from rest on two separate frictionless inclined planes making angles of $60^\circ$ and $30^\circ$ with the horizontal, respectively. What is the magnitude of the relative vertical acceleration of block $A$ with respect to block $B$? ($g = 9.8\text{ ms}^{-2}$)
    }
    \end{english}
\end{tcolorbox}

% ------------------------------------------
% HIGH-QUALITY TIKZ DIAGRAM 18
% ------------------------------------------
\vspace{1em}
\begin{center}
\begin{tikzpicture}[>=Stealth, scale=1.1]
    % Plane A (60 degrees)
    \begin{scope}[shift={(0,0)}]
        \draw[thick, fill=gray!20] (0, 3.464) -- (0,0) -- (2,0) -- cycle;
        \draw (1.5, 0) arc (180:120:0.5);
        \node at (1.2, 0.35) {$60^\circ$};
        
        \coordinate (PosA) at (1.0, 1.732);
        \begin{scope}[shift={(PosA)}, rotate=-60]
            \draw[thick, fill=orange!80!black] (-0.4, 0) rectangle (0.4, 0.5);
            \node[white, font=\bfseries] at (0, 0.25) {$A$};
        \end{scope}
        
        % Acceleration vectors
        \draw[->, line width=1.5pt, primary] (PosA) ++(0.3, 0.2) -- ++( {1.2*cos(-60)}, {1.2*sin(-60)} ) node[right] {$a$};
        \draw[->, line width=1.5pt, red] (PosA) ++(0.3, 0.2) -- ++( 0, {1.2*sin(-60)} ) node[left] {$a_{A,y}$};
    \end{scope}
    
    % Plane B (30 degrees)
    \begin{scope}[shift={(3.5, 0)}]
        \draw[thick, fill=gray!20] (0, 2) -- (0,0) -- (3.464,0) -- cycle;
        \draw (2.7, 0) arc (180:150:0.76);
        \node at (2.2, 0.3) {$30^\circ$};
        
        \coordinate (PosB) at (1.732, 1.0);
        \begin{scope}[shift={(PosB)}, rotate=-30]
            \draw[thick, fill=cyan!70!black] (-0.4, 0) rectangle (0.4, 0.5);
            \node[white, font=\bfseries] at (0, 0.25) {$B$};
        \end{scope}
        
        % Acceleration vectors
        \draw[->, line width=1.5pt, primary] (PosB) ++(0.2, 0.4) -- ++( {1.2*cos(-30)}, {1.2*sin(-30)} ) node[right] {$a$};
        \draw[->, line width=1.5pt, red] (PosB) ++(0.2, 0.4) -- ++( 0, {1.2*sin(-30)} ) node[left] {$a_{B,y}$};
    \end{scope}
\end{tikzpicture}
\end{center}
\vspace{1em}

% ------------------------------------------
% OPTIONS & ANSWER 18
% ------------------------------------------
\begin{itemize}[label={}, leftmargin=1cm, itemsep=0.5em]
    \item[\textbf{A)}] আপেক্ষিক উল্লম্ব ত্বরণ $4.9\text{ ms}^{-2}$
    \item[\textbf{B)}] আপেক্ষিক উল্লম্ব ত্বরণ $9.8\text{ ms}^{-2}$
    \item[\textbf{C)}] আপেক্ষিক উল্লম্ব ত্বরণ $0\text{ ms}^{-2}$
    \item[\textbf{D)}] আপেক্ষিক উল্লম্ব ত্বরণ $2.45\text{ ms}^{-2}$
\end{itemize}

\vspace{0.5em}
\begin{tcolorbox}[answerbox]
    \textbf{\color{success} উত্তর:} A) আপেক্ষিক উল্লম্ব ত্বরণ $4.9\text{ ms}^{-2}$
\end{tcolorbox}
\vspace{1em}

% ------------------------------------------
% SOLUTION SECTION 18
% ------------------------------------------
\begin{tcolorbox}[solutionbox={সমাধান (Solution)}]
    \noindent
    $\theta$ নতিকোণের আনত তল বেয়ে পতনশীল ব্লকের তলের সমান্তরাল ত্বরণ $a = g \sin\theta$।
    আনত তলটি ভূমির সাথে $\theta$ কোণে থাকায় উল্লম্ব দিকের সাথে এর কোণ $(90^\circ - \theta)$।
    
    ত্বরণের উল্লম্ব উপাংশ:
    \[ a_y = a \cos(90^\circ - \theta) = (g \sin\theta) \sin\theta = g \sin^2\theta \]
    
    $A$ ব্লকের উল্লম্ব ত্বরণ:
    \begin{align*}
        a_{A,y} &= g \sin^2 60^\circ \\
                &= g \left(\frac{\sqrt{3}}{2}\right)^2 = \frac{3}{4}g
    \end{align*}
    
    $B$ ব্লকের উল্লম্ব ত্বরণ:
    \begin{align*}
        a_{B,y} &= g \sin^2 30^\circ \\
                &= g \left(\frac{1}{2}\right)^2 = \frac{1}{4}g
    \end{align*}
    
    অতএব, $B$-এর সাপেক্ষে $A$-এর আপেক্ষিক উল্লম্ব ত্বরণ:
    \begin{align*}
        a_{\text{rel}, y} &= a_{A,y} - a_{B,y} \\
                          &= \frac{3}{4}g - \frac{1}{4}g = \frac{1}{2}g \\
                          &= \frac{9.8}{2} = \mathbf{4.9}\text{ ms}^{-2}
    \end{align*}
\end{tcolorbox}

\newpage

% ------------------------------------------
% QUESTION SECTION 19
% ------------------------------------------
\begin{tcolorbox}[questionbox={প্রশ্ন (Question) 19}]
    \noindent
    মাটিতে পোঁতা একটি পেরেক উল্লম্ব বরাবর সর্বোচ্চ $40\text{ N}$ ঊর্ধ্বমুখী বল সহ্য করতে পারে। একটি হালকা অপ্রসারণশীল দড়ি ঘর্ষণহীন কপিকলের ওপর দিয়ে গেছে; দড়ির এক প্রান্ত অনুভূমিকের সাথে $60^\circ$ কোণে ঐ পেরেকের সাথে বাঁধা। $2\text{ kg}$ ভরের একটি বানর দড়ির অপর উল্লম্ব অংশ বেয়ে সর্বোচ্চ কত ঊর্ধ্বমুখী ত্বরণে উপরে উঠতে পারবে যাতে পেরেকটি মাটি থেকে উপড়ে না যায়? ($g = 9.8\text{ ms}^{-2}$)
    
    \vspace{0.8em}
    \begin{english}
    \noindent\textit{\color{darkgray}
    A nail driven into the ground can withstand a maximum vertical pulling force of $40\text{ N}$ before being pulled out. A light rope passes over a frictionless pulley; one end is fixed to the nail making an angle of $60^\circ$ with the horizontal, while a monkey of mass $2\text{ kg}$ climbs up the other vertical section. What is the maximum acceleration with which the monkey can climb up without pulling the nail out? ($g = 9.8\text{ ms}^{-2}$)
    }
    \end{english}
\end{tcolorbox}

% ------------------------------------------
% HIGH-QUALITY TIKZ DIAGRAM 19
% ------------------------------------------
\vspace{1em}
\begin{center}
\begin{tikzpicture}[>=Stealth, scale=1.2]
    % Ground
    \fill[gray!20] (-1, -0.3) rectangle (4, 0);
    \draw[thick, darkgray] (-1, 0) -- (4, 0);
    
    % Nail
    \draw[thick, fill=gray] (0, 0.2) -- (-0.1, 0.2) -- (-0.1, -0.3) -- (0, -0.5) -- (0.1, -0.3) -- (0.1, 0.2) -- cycle;
    \fill[black] (0, 0.1) circle (0.05);
    
    % Pulley
    \coordinate (Pulley) at (1.5, 2.6); % 1.5 * tan(60) = 2.598
    \draw[thick, fill=gray!30] (Pulley) circle (0.2);
    \fill[black] (Pulley) circle (0.05);
    
    % Support for pulley
    \draw[thick] (Pulley) -- (1.5, 3.2);
    \fill[pattern=north east lines, pattern color=gray] (1.2, 3.2) rectangle (1.8, 3.4);
    
    % Rope
    \draw[line width=1.2pt, red] (0, 0.1) -- (1.4, 2.45);
    \draw[line width=1.2pt, red] (1.7, 2.6) -- (1.7, 0.5);
    
    % Monkey (Block representation)
    \fill[brown!80!black, rounded corners=3pt] (1.4, 0.5) rectangle (2.0, 1.3);
    \node[white, font=\bfseries] at (1.7, 0.9) {$2\text{kg}$};
    \node[below, font=\bfseries] at (1.7, 0.4) {Monkey};
    
    % Angles and Forces
    \draw[dashed] (0, 0.1) -- (1.5, 0.1);
    \draw (0.6, 0.1) arc (0:60:0.6);
    \node at (0.9, 0.4) {$60^\circ$};
    
    % Acceleration arrow
    \draw[->, line width=1.5pt, primary] (2.2, 0.5) -- (2.2, 1.5) node[midway, right] {$a$};
    
    % Force on nail
    \draw[->, line width=1.5pt, blue] (0, 0.1) -- (0, 1.2) node[left] {$T_y \le 40\text{ N}$};
\end{tikzpicture}
\end{center}
\vspace{1em}

% ------------------------------------------
% OPTIONS & ANSWER 19
% ------------------------------------------
\begin{itemize}[label={}, leftmargin=1cm, itemsep=0.5em]
    \item[\textbf{A)}] সর্বোচ্চ ত্বরণ $13.29\text{ ms}^{-2}$
    \item[\textbf{B)}] সর্বোচ্চ ত্বরণ $10.20\text{ ms}^{-2}$
    \item[\textbf{C)}] সর্বোচ্চ ত্বরণ $15.45\text{ ms}^{-2}$
    \item[\textbf{D)}] সর্বোচ্চ ত্বরণ $8.50\text{ ms}^{-2}$
\end{itemize}

\vspace{0.5em}
\begin{tcolorbox}[answerbox]
    \textbf{\color{success} উত্তর:} A) সর্বোচ্চ ত্বরণ $13.29\text{ ms}^{-2}$
\end{tcolorbox}
\vspace{1em}

% ------------------------------------------
% SOLUTION SECTION 19
% ------------------------------------------
\begin{tcolorbox}[solutionbox={সমাধান (Solution)}]
    \noindent
    পেরেকের ওপর প্রযুক্ত বলের উল্লম্ব উপাংশ সর্বোচ্চ $40\text{ N}$ হতে পারে:
    \begin{align*}
        T \sin 60^\circ &\le 40\text{ N} \\
        T \times \frac{\sqrt{3}}{2} &\le 40 \\
        T &\le \frac{80}{\sqrt{3}} \approx 46.188\text{ N}
    \end{align*}
    
    বানরটি $a$ ত্বরণে উপরে ওঠার ক্ষেত্রে দড়ির টান ($T$):
    \begin{align*}
        T - mg &= ma \\
        ma &= T - mg
    \end{align*}
    
    সর্বোচ্চ ত্বরণ:
    \begin{align*}
        a &= \frac{T}{m} - g \\
          &= \frac{46.188}{2} - 9.8 \\
          &= 23.094 - 9.8 \\
          &= \mathbf{13.29}\text{ ms}^{-2}
    \end{align*}
\end{tcolorbox}

\newpage

% ------------------------------------------
% QUESTION SECTION 20
% ------------------------------------------
\begin{tcolorbox}[questionbox={প্রশ্ন (Question) 20}]
    \noindent
    $2\text{ kg}$ ভরের একটি বস্তু অনুভূমিক মসৃণ মেঝেতে $4\text{ ms}^{-1}$ বেগে গতিশীল অবস্থায় $x = 0.5\text{ m}$ হতে $x = 1.5\text{ m}$ সীমার একটি অমসৃণ অঞ্চলে প্রবেশ করল। এই অঞ্চলে বস্তুটির ওপর অবস্থান-নির্ভর বাধাদানকারী বল $F = -kx$ ক্রিয়া করে, যেখানে $k = 12\text{ Nm}^{-1}$। অমসৃণ অঞ্চল থেকে বের হওয়ার মুহূর্তে ($x = 1.5\text{ m}$) বস্তুর বেগ কত হবে?
    
    \vspace{0.8em}
    \begin{english}
    \noindent\textit{\color{darkgray}
    A body of mass $2\text{ kg}$ moving horizontally at $4\text{ ms}^{-1}$ on a smooth floor enters a rough zone extending from $x = 0.5\text{ m}$ to $x = 1.5\text{ m}$. Inside this zone, it experiences a position-dependent resistive force $F = -kx$, where $k = 12\text{ Nm}^{-1}$. What is the velocity of the body as it exits the rough zone at $x = 1.5\text{ m}$?
    }
    \end{english}
\end{tcolorbox}

% ------------------------------------------
% HIGH-QUALITY TIKZ DIAGRAM 20
% ------------------------------------------
\vspace{1em}
\begin{center}
\begin{tikzpicture}[>=Stealth, scale=1.5]
    % Floor
    \draw[thick] (-1, 0) -- (3.5, 0);
    
    % Smooth Region 1
    \fill[cyan!10] (-1, 0) rectangle (0.5, -0.2);
    \node[below, font=\footnotesize] at (-0.25, -0.2) {মসৃণ (Smooth)};
    
    % Rough Region
    \fill[pattern=north west lines, pattern color=red!50] (0.5, 0) rectangle (1.5, -0.2);
    \draw[thick, red] (0.5, 0) -- (1.5, 0);
    \node[below, text=red!80!black, font=\footnotesize] at (1.0, -0.2) {অমসৃণ ($F = -kx$)};
    
    % Smooth Region 2
    \fill[cyan!10] (1.5, 0) rectangle (3.5, -0.2);
    \node[below, font=\footnotesize] at (2.5, -0.2) {মসৃণ (Smooth)};
    
    % X markers
    \draw[dashed] (0.5, 0.5) -- (0.5, -0.8) node[below] {$x = 0.5$};
    \draw[dashed] (1.5, 0.5) -- (1.5, -0.8) node[below] {$x = 1.5$};
    
    % Block at entry
    \draw[thick, fill=orange!80!black] (-0.8, 0) rectangle (-0.2, 0.4);
    \node[white, font=\bfseries] at (-0.5, 0.2) {$2\text{kg}$};
    \draw[->, line width=1.5pt, primary] (-0.2, 0.2) -- (0.3, 0.2) node[above] {$u = 4 \text{ ms}^{-1}$};
    
    % Block at exit
    \draw[thick, fill=orange!80!black, opacity=0.5] (1.5, 0) rectangle (2.1, 0.4);
    \draw[->, line width=1.5pt, primary] (2.1, 0.2) -- (2.6, 0.2) node[above] {$v = ?$};
\end{tikzpicture}
\end{center}
\vspace{1em}

% ------------------------------------------
% OPTIONS & ANSWER 20
% ------------------------------------------
\begin{itemize}[label={}, leftmargin=1cm, itemsep=0.5em]
    \item[\textbf{A)}] শেষ বেগ $2.0\text{ ms}^{-1}$
    \item[\textbf{B)}] শেষ বেগ $1.5\text{ ms}^{-1}$
    \item[\textbf{C)}] শেষ বেগ $2.5\text{ ms}^{-1}$
    \item[\textbf{D)}] শেষ বেগ $0\text{ ms}^{-1}$
\end{itemize}

\vspace{0.5em}
\begin{tcolorbox}[answerbox]
    \textbf{\color{success} উত্তর:} A) শেষ বেগ $2.0\text{ ms}^{-1}$
\end{tcolorbox}
\vspace{1em}

% ------------------------------------------
% SOLUTION SECTION 20
% ------------------------------------------
\begin{tcolorbox}[solutionbox={সমাধান (Solution)}]
    \noindent
    \textbf{ধাপ ১: পরিবর্তনশীল বল দ্বারা কৃতকাজ নির্ণয়} \\
    বাধাদানকারী বল দ্বারা কৃতকাজ:
    \begin{align*}
        W &= \int_{x_i}^{x_f} F \, dx \\
          &= \int_{0.5}^{1.5} (-12x) \, dx \\
          &= -12 \left[ \frac{x^2}{2} \right]_{0.5}^{1.5} \\
          &= -6 \left[ (1.5)^2 - (0.5)^2 \right] \\
          &= -6 [2.25 - 0.25] \\
          &= -6 \times 2 = -12\text{ J}
    \end{align*}
    
    \vspace{0.5em}
    \noindent\textbf{ধাপ ২: কাজ-শক্তি উপপাদ্য প্রয়োগ} \\
    কাজ-শক্তি উপপাদ্য অনুসারে, কৃতকাজ = গতিশক্তির পরিবর্তন:
    \begin{align*}
        W &= \Delta K = \frac{1}{2} m v^2 - \frac{1}{2} m u^2 \\
        -12 &= \frac{1}{2} (2) v^2 - \frac{1}{2} (2) (4)^2 \\
        -12 &= v^2 - 16 \\
        v^2 &= 16 - 12 = 4 \\
        v &= \sqrt{4} = \mathbf{2.0}\text{ ms}^{-1}
    \end{align*}
\end{tcolorbox}

\newpage

% ------------------------------------------
% QUESTION SECTION 21
% ------------------------------------------
\begin{tcolorbox}[questionbox={প্রশ্ন (Question) 21}]
    \noindent
    $L$ দৈর্ঘ্য এবং $M$ মোট ভরের একটি অসম সরু রডের রৈখিক ভর ঘনত্ব প্রান্ত $A$ হতে $B$-এর দিকে $\lambda(x) = \lambda_0 \left(1 + \frac{x}{L}\right)$ নিয়ম অনুযায়ী বৃদ্ধি পায়, যেখানে $x$ হলো $A$ প্রান্ত হতে দূরত্ব। রডটির $A$ প্রান্তগামী এবং দৈর্ঘ্যের সাথে লম্ব অক্ষের সাপেক্ষে জড়তার ভ্রামক কত হবে?
    
    \vspace{0.8em}
    \begin{english}
    \noindent\textit{\color{darkgray}
    A thin non-uniform rod of length $L$ and total mass $M$ has a linear mass density that varies from end $A$ to end $B$ as $\lambda(x) = \lambda_0 \left(1 + \frac{x}{L}\right)$, where $x$ is the distance from end $A$. What is the moment of inertia of the rod about an axis passing through end $A$ and perpendicular to its length?
    }
    \end{english}
\end{tcolorbox}

% ------------------------------------------
% HIGH-QUALITY TIKZ DIAGRAM 21
% ------------------------------------------
\vspace{1em}
\begin{center}
\begin{tikzpicture}[>=Stealth, scale=1.5]
    % Axis of rotation
    \draw[dash dot, line width=1.5pt, primary] (0, -1) -- (0, 1) node[above] {Axis};
    
    % Rod with gradient
    \shade[left color=gray!30, right color=gray!80] (0, -0.15) rectangle (4, 0.15);
    \draw[thick] (0, -0.15) rectangle (4, 0.15);
    
    % Labels A and B
    \node[above, font=\bfseries] at (0, 0.2) {$A$};
    \node[above, font=\bfseries] at (4, 0.2) {$B$};
    
    % Distance x
    \draw[->, thick] (0, -0.4) -- (1.5, -0.4) node[midway, below] {$x$};
    
    % Small element dm
    \fill[red!50] (1.5, -0.15) rectangle (1.7, 0.15);
    \node[below, text=red!80!black, font=\footnotesize] at (1.6, -0.4) {$dx, dm$};
    
    % Length L
    \draw[<->, dashed] (0, 0.6) -- (4, 0.6) node[midway, fill=white] {$L$};
    
    % Density function
    \node[above] at (2, 0.8) {$\lambda(x) = \lambda_0 \left(1 + \frac{x}{L}\right)$};
\end{tikzpicture}
\end{center}
\vspace{1em}

% ------------------------------------------
% OPTIONS & ANSWER 21
% ------------------------------------------
\begin{itemize}[label={}, leftmargin=1cm, itemsep=0.5em]
    \item[\textbf{A)}] জড়তার ভ্রামক $\frac{7}{18} M L^2$
    \item[\textbf{B)}] জড়তার ভ্রামক $\frac{5}{12} M L^2$
    \item[\textbf{C)}] জড়তার ভ্রামক $\frac{2}{5} M L^2$
    \item[\textbf{D)}] জড়তার ভ্রামক $\frac{3}{7} M L^2$
\end{itemize}

\vspace{0.5em}
\begin{tcolorbox}[answerbox]
    \textbf{\color{success} উত্তর:} A) জড়তার ভ্রামক $\frac{7}{18} M L^2$
\end{tcolorbox}
\vspace{1em}

% ------------------------------------------
% SOLUTION SECTION 21
% ------------------------------------------
\begin{tcolorbox}[solutionbox={সমাধান (Solution)}]
    \noindent
    \textbf{ধাপ ১: মোট ভরের সাথে $\lambda_0$ এর সম্পর্ক} \\
    রডের মোট ভর $M$ হলো ক্ষুদ্র ভরগুলোর ($dm = \lambda(x) dx$) যোগফল:
    \begin{align*}
        M &= \int_0^L \lambda(x) \, dx \\
          &= \lambda_0 \int_0^L \left(1 + \frac{x}{L}\right) \, dx \\
          &= \lambda_0 \left[ x + \frac{x^2}{2L} \right]_0^L \\
          &= \lambda_0 \left( L + \frac{L^2}{2L} \right) \\
          &= \lambda_0 \left( L + \frac{L}{2} \right) = \frac{3}{2} \lambda_0 L
    \end{align*}
    এখান থেকে $\lambda_0$-এর মান:
    \[ \lambda_0 = \frac{2M}{3L} \]
    
    \vspace{0.5em}
    \noindent\textbf{ধাপ ২: জড়তার ভ্রামক নির্ণয়} \\
    $A$ প্রান্তের সাপেক্ষে জড়তার ভ্রামক:
    \begin{align*}
        I_A &= \int_0^L x^2 \, dm \\
            &= \int_0^L x^2 \lambda(x) \, dx \\
            &= \lambda_0 \int_0^L \left( x^2 + \frac{x^3}{L} \right) \, dx \\
            &= \lambda_0 \left[ \frac{x^3}{3} + \frac{x^4}{4L} \right]_0^L \\
            &= \lambda_0 \left( \frac{L^3}{3} + \frac{L^4}{4L} \right) \\
            &= \lambda_0 L^3 \left( \frac{1}{3} + \frac{1}{4} \right) \\
            &= \frac{7}{12} \lambda_0 L^3
    \end{align*}
    
    $\lambda_0$-এর মান প্রতিস্থাপন করে পাই:
    \begin{align*}
        I_A &= \frac{7}{12} \left(\frac{2M}{3L}\right) L^3 \\
            &= \frac{14}{36} M L^2 \\
            &= \mathbf{\frac{7}{18} M L^2}
    \end{align*}
\end{tcolorbox}

\newpage

% ------------------------------------------
% QUESTION SECTION 22
% ------------------------------------------
\begin{tcolorbox}[questionbox={প্রশ্ন (Question) 22}]
    \noindent
    প্রতিটি $M$ ভর এবং $R$ ব্যাসার্ধের $7$টি অভিন্ন সুষম বৃত্তাকার চাকতি একই সমতলে এমনভাবে জোড়া লাগানো হয়েছে যাতে একটি কেন্দ্রীয় চাকতিকে ঘিরে $6$টি চাকতি পরস্পর স্পর্শ করে অবস্থান করে। কেন্দ্রীয় চাকতির কেন্দ্র থেকে $3R$ দূরত্বে অবস্থিত যেকোনো একটি প্রান্তীয় চাকতির পরিধিস্থ বিন্দু $P$ দিয়ে গমনকারী এবং চাকতিগুলোর তলের ওপর লম্ব অক্ষের সাপেক্ষে সম্পূর্ণ সংস্থাটির জড়তার ভ্রামক কত হবে?
    
    \vspace{0.8em}
    \begin{english}
    \noindent\textit{\color{darkgray}
    Seven identical uniform circular discs, each of mass $M$ and radius $R$, are welded together symmetrically in a single plane such that a central disc is surrounded by six touching outer discs. What is the moment of inertia of the entire system about an axis passing through point $P$ on the outer rim of an exterior disc (at a distance $3R$ from the central disc's center) and perpendicular to the plane of the discs?
    }
    \end{english}
\end{tcolorbox}

% ------------------------------------------
% HIGH-QUALITY TIKZ DIAGRAM 22
% ------------------------------------------
\vspace{1em}
\begin{center}
\begin{tikzpicture}[>=Stealth, scale=1.1]
    \def\myR{0.8}
    
    % Central disc
    \draw[thick, fill=cyan!15] (0,0) circle (\myR);
    \fill[black] (0,0) circle (0.05) node[below] {$O$};
    
    % 6 surrounding discs
    \foreach \angle in {60, 120, 180, 240, 300} {
        \draw[thick, fill=cyan!5] ({\angle}:2*\myR) circle (\myR);
        \fill[black] ({\angle}:2*\myR) circle (0.04);
    }
    
    % Highlight one outer disc and point P
    \draw[thick, fill=orange!20] (0:2*\myR) circle (\myR);
    \fill[black] (0:2*\myR) circle (0.05) node[above left] {$O_{\text{out}}$};
    
    % Point P (on the outer rim)
    \fill[red] (0:3*\myR) circle (0.06) node[right, font=\bfseries] {$P$};
    
    % Distance annotations
    \draw[dashed, thick] (0,0) -- (0:3*\myR);
    \draw[<->] (0, -1.2*\myR) -- (2*\myR, -1.2*\myR) node[midway, below] {$2R$};
    \draw[<->] (0, -1.8*\myR) -- (3*\myR, -1.8*\myR) node[midway, below] {$3R$};
    
    % Axis at P
    \draw[line width=1.5pt, red] (0:3*\myR) circle (0.15);
    \node[red, above, align=center, font=\footnotesize] at (0:3*\myR) {Axis \\ $\perp$ plane};
\end{tikzpicture}
\end{center}
\vspace{1em}

% ------------------------------------------
% OPTIONS & ANSWER 22
% ------------------------------------------
\begin{itemize}[label={}, leftmargin=1cm, itemsep=0.5em]
    \item[\textbf{A)}] জড়তার ভ্রামক $\frac{181}{2} M R^2$
    \item[\textbf{B)}] জড়তার ভ্রামক $\frac{55}{2} M R^2$
    \item[\textbf{C)}] জড়তার ভ্রামক $\frac{73}{2} M R^2$
    \item[\textbf{D)}] জড়তার ভ্রামক $90 M R^2$
\end{itemize}

\vspace{0.5em}
\begin{tcolorbox}[answerbox]
    \textbf{\color{success} উত্তর:} A) জড়তার ভ্রামক $\frac{181}{2} M R^2$
\end{tcolorbox}
\vspace{1em}

% ------------------------------------------
% SOLUTION SECTION 22
% ------------------------------------------
\begin{tcolorbox}[solutionbox={সমাধান (Solution)}]
    \noindent
    \textbf{ধাপ ১: কেন্দ্র $O$-এর সাপেক্ষে জড়তার ভ্রামক} \\
    কেন্দ্রীয় চাকতির নিজস্ব জড়তার ভ্রামক:
    \[ I_{\text{center}} = \frac{1}{2} M R^2 \]
    
    বাইরের প্রতিটি চাকতির কেন্দ্র $O$ থেকে $d = 2R$ দূরত্বে অবস্থিত। সমান্তরাল অক্ষ উপপাদ্য মতে প্রতিটি বাইরের চাকতির ভ্রামক:
    \begin{align*}
        I_{\text{outer}, O} &= I_{\text{cm}} + M d^2 \\
        &= \frac{1}{2} M R^2 + M (2R)^2 \\
        &= \frac{1}{2} M R^2 + 4 M R^2 = \frac{9}{2} M R^2
    \end{align*}
    
    $O$ বিন্দুর সাপেক্ষে সম্পূর্ণ সংস্থার (৭টি চাকতি) জড়তার ভ্রামক:
    \begin{align*}
        I_O &= I_{\text{center}} + 6 \times I_{\text{outer}, O} \\
        &= \frac{1}{2} M R^2 + 6 \left(\frac{9}{2} M R^2\right) \\
        &= \frac{1}{2} M R^2 + 27 M R^2 = \frac{55}{2} M R^2
    \end{align*}
    
    \vspace{0.5em}
    \noindent\textbf{ধাপ ২: বিন্দু $P$-এর সাপেক্ষে জড়তার ভ্রামক} \\
    সংস্থাটির মোট ভর $M_{\text{total}} = 7M$। 
    বিন্দু $P$ কেন্দ্র $O$ থেকে $3R$ দূরত্বে অবস্থিত। সমান্তরাল অক্ষ উপপাদ্য প্রয়োগ করে বিন্দু $P$-গামী অক্ষের সাপেক্ষে জড়তার ভ্রামক:
    \begin{align*}
        I_P &= I_O + M_{\text{total}} (3R)^2 \\
        &= \frac{55}{2} M R^2 + (7M) (9 R^2) \\
        &= \frac{55}{2} M R^2 + 63 M R^2 \\
        &= \left(\frac{55 + 126}{2}\right) M R^2 \\
        &= \mathbf{\frac{181}{2} M R^2}
    \end{align*}
\end{tcolorbox}

\newpage

% ------------------------------------------
% QUESTION SECTION 23
% ------------------------------------------
\begin{tcolorbox}[questionbox={প্রশ্ন (Question) 23}]
    \noindent
    একটি অনুভূমিক বৃত্তাকার চাকতি $D_1$-এর ভর $M$ এবং ব্যাসার্ধ $R$। এর ব্যাসের দুই বিপরীত প্রান্তে অপর দুটি খাড়া উল্লম্ব বৃত্তাকার চাকতি $D_2$ ও $D_3$ (উভয়েরই ভর $M$ ও ব্যাসার্ধ $R$) লম্বভাবে সংযুক্ত আছে। $D_1$-এর কেন্দ্রগামী এবং এর তলের ওপর লম্ব উল্লম্ব অক্ষ $OO'$-এর সাপেক্ষে তিন চাকতির এই ব্যবস্থার মোট জড়তার ভ্রামক কত?
    
    \vspace{0.8em}
    \begin{english}
    \noindent\textit{\color{darkgray}
    A horizontal circular disc $D_1$ of mass $M$ and radius $R$ has two identical vertical circular discs $D_2$ and $D_3$ (each of mass $M$ and radius $R$) attached perpendicularly at the opposite ends of its diameter. What is the total moment of inertia of this composite three-disc structure about a vertical axis $OO'$ passing through the center of $D_1$ and perpendicular to its face?
    }
    \end{english}
\end{tcolorbox}

% ------------------------------------------
% HIGH-QUALITY TIKZ DIAGRAM 23
% ------------------------------------------
\vspace{1em}
\begin{center}
\begin{tikzpicture}[>=Stealth, scale=1.5]
    % Vertical Axis OO'
    \draw[dash dot, line width=1.5pt, primary] (0, -2) -- (0, 2) node[above] {$O'$};
    \node[below, text=primary, font=\bfseries] at (0, -2) {$O$};
    
    % Horizontal Disc D1 (Ellipse for 3D perspective)
    \draw[thick, fill=cyan!20, fill opacity=0.7] (0,0) ellipse [x radius=2, y radius=0.6];
    \node[above, font=\bfseries] at (0, 0.2) {$D_1$};
    
    % Vertical Disc D2 at (-2,0)
    \draw[thick, fill=orange!30, fill opacity=0.8] (-2,0) ellipse [x radius=0.3, y radius=1];
    \node[left, font=\bfseries] at (-2.3, 0) {$D_2$};
    \fill[black] (-2,0) circle (0.04);
    
    % Vertical Disc D3 at (2,0)
    \draw[thick, fill=orange!30, fill opacity=0.8] (2,0) ellipse [x radius=0.3, y radius=1];
    \node[right, font=\bfseries] at (2.3, 0) {$D_3$};
    \fill[black] (2,0) circle (0.04);
    
    % Center O
    \fill[black] (0,0) circle (0.05);
    
    % Radii annotations
    \draw[dashed, thick] (0,0) -- (2,0) node[midway, below] {$R$};
    \draw[dashed, thick] (0,0) -- (-2,0) node[midway, below] {$R$};
    \draw[dashed, thick] (2,0) -- (2, 1) node[midway, right] {$R$};
\end{tikzpicture}
\end{center}
\vspace{1em}

% ------------------------------------------
% OPTIONS & ANSWER 23
% ------------------------------------------
\begin{itemize}[label={}, leftmargin=1cm, itemsep=0.5em]
    \item[\textbf{A)}] জড়তার ভ্রামক $3 M R^2$
    \item[\textbf{B)}] জড়তার ভ্রামক $\frac{5}{2} M R^2$
    \item[\textbf{C)}] জড়তার ভ্রামক $2 M R^2$
    \item[\textbf{D)}] জড়তার ভ্রামক $\frac{7}{2} M R^2$
\end{itemize}

\vspace{0.5em}
\begin{tcolorbox}[answerbox]
    \textbf{\color{success} উত্তর:} A) জড়তার ভ্রামক $3 M R^2$
\end{tcolorbox}
\vspace{1em}

% ------------------------------------------
% SOLUTION SECTION 23
% ------------------------------------------
\begin{tcolorbox}[solutionbox={সমাধান (Solution)}]
    \noindent
    অক্ষ $OO'$ চাকতি $D_1$-এর তলের ওপর লম্ব এবং এর কেন্দ্রগামী। 
    অতএব $D_1$-এর জড়তার ভ্রামক:
    \[ I_{D_1} = \frac{1}{2} M R^2 \]
    
    উল্লম্ব চাকতি $D_2$ ও $D_3$-এর ক্ষেত্রে অক্ষ $OO'$ তাদের নিজ নিজ তলের সমান্তরাল এবং তাদের কেন্দ্র থেকে $R$ দূরত্বে অবস্থিত।
    আমরা জানি, একটি চাকতির তার ব্যাসের সাপেক্ষে জড়তার ভ্রামক:
    \[ I_{\text{dia}} = \frac{1}{4} M R^2 \]
    
    সমান্তরাল অক্ষ উপপাদ্য প্রয়োগ করে $OO'$ অক্ষের সাপেক্ষে $D_2$ (ও অনুরূপভাবে $D_3$)-এর জড়তার ভ্রামক:
    \begin{align*}
        I_{D_2} &= I_{\text{dia}} + M R^2 \\
        &= \frac{1}{4} M R^2 + M R^2 \\
        &= \frac{5}{4} M R^2
    \end{align*}
    
    সম্পূর্ণ ব্যবস্থার মোট জড়তার ভ্রামক:
    \begin{align*}
        I_{\text{total}} &= I_{D_1} + 2 \times I_{D_2} \\
        &= \frac{1}{2} M R^2 + 2 \left(\frac{5}{4} M R^2\right) \\
        &= \frac{1}{2} M R^2 + \frac{5}{2} M R^2 \\
        &= \frac{6}{2} M R^2 \\
        &= \mathbf{3 M R^2}
    \end{align*}
\end{tcolorbox}

\newpage

% ------------------------------------------
% QUESTION SECTION 24
% ------------------------------------------
\begin{tcolorbox}[questionbox={প্রশ্ন (Question) 24}]
    \noindent
    $50\text{ m}$ বাঁকের ব্যাসার্ধ বিশিষ্ট একটি পাহাড়ি সড়কের ব্যাংকিং কোণ $30^\circ$। টায়ার ও রাস্তার মধ্যকার স্থিতি ঘর্ষণ গুণাঙ্ক $\mu_s = 0.2$ হলে, রাস্তা থেকে নিচের দিকে পিছলে না পড়ার জন্য একটি গাড়ির সর্বনিম্ন সম্ভাব্য নিরাপদ দ্রুতি কত হওয়া প্রয়োজন? ($g = 9.8\text{ ms}^{-2}$)
    
    \vspace{0.8em}
    \begin{english}
    \noindent\textit{\color{darkgray}
    A curved mountain road with a radius of curvature of $50\text{ m}$ is banked at an angle of $30^\circ$. If the coefficient of static friction between the tires and the road surface is $\mu_s = 0.2$, what is the minimum safe speed required for a car so that it does not slide downwards along the inclined surface? ($g = 9.8\text{ ms}^{-2}$)
    }
    \end{english}
\end{tcolorbox}

% ------------------------------------------
% HIGH-QUALITY TIKZ DIAGRAM 24
% ------------------------------------------
\vspace{1em}
\begin{center}
\begin{tikzpicture}[>=Stealth, scale=1.5]
    % Banked road
    \draw[thick, fill=gray!20] (0,0) -- (4,0) -- (4, 2.309) coordinate (Top) -- cycle;
    \draw (1,0) arc (180:150:1);
    \node at (0.7, 0.25) {$30^\circ$};
    
    % Car Block
    \coordinate (Car) at (2, 1.154);
    \begin{scope}[shift={(Car)}, rotate=30]
        \draw[thick, fill=cyan!70!black] (-0.4, 0) rectangle (0.4, 0.5);
        \fill[black] (-0.25, -0.1) circle (0.1); % Wheel
        \fill[black] (0.25, -0.1) circle (0.1);  % Wheel
    \end{scope}
    
    % Forces
    \draw[->, line width=1.5pt, red] (Car) ++(0, 0.25) -- ++(0, -1.2) node[below] {$mg$};
    \draw[->, line width=1.5pt, green!60!black] (Car) ++(0, 0.25) -- ++({0.8*cos(120)}, {0.8*sin(120)}) node[left] {$N$};
    
    % Friction UP the incline (to prevent sliding down)
    \draw[->, line width=1.5pt, orange!80!black] (Car) ++(0, 0.1) -- ++({1.0*cos(30)}, {1.0*sin(30)}) node[midway, above left] {$f_s$};
    
    % Center of curve direction
    \draw[->, thick, dashed, primary] (Car) ++(0, 0.25) -- ++(-1.5, 0) node[left] {বাঁকের কেন্দ্র ($R$)};
\end{tikzpicture}
\end{center}
\vspace{1em}

% ------------------------------------------
% OPTIONS & ANSWER 24
% ------------------------------------------
\begin{itemize}[label={}, leftmargin=1cm, itemsep=0.5em]
    \item[\textbf{A)}] সর্বনিম্ন দ্রুতি $12.87\text{ ms}^{-1}$
    \item[\textbf{B)}] সর্বনিম্ন দ্রুতি $16.82\text{ ms}^{-1}$
    \item[\textbf{C)}] সর্বনিম্ন দ্রুতি $9.50\text{ ms}^{-1}$
    \item[\textbf{D)}] সর্বনিম্ন দ্রুতি $14.25\text{ ms}^{-1}$
\end{itemize}

\vspace{0.5em}
\begin{tcolorbox}[answerbox]
    \textbf{\color{success} উত্তর:} A) সর্বনিম্ন দ্রুতি $12.87\text{ ms}^{-1}$
\end{tcolorbox}
\vspace{1em}

% ------------------------------------------
% SOLUTION SECTION 24
% ------------------------------------------
\begin{tcolorbox}[solutionbox={সমাধান (Solution)}]
    \noindent
    ব্যাংকিং যুক্ত অমসৃণ তলে না পিছলিয়ে চলার (নিচের দিকে) সর্বনিম্ন বেগের তাত্ত্বিক সমীকরণ:
    \[ v_{\min} = \sqrt{r g \left(\frac{\tan\theta - \mu_s}{1 + \mu_s \tan\theta}\right)} \]
    
    প্রদত্ত মানসমূহ: $r = 50\text{ m}$, $\theta = 30^\circ$, $\mu_s = 0.2$, $g = 9.8\text{ ms}^{-2}$।
    
    এখানে, $\tan 30^\circ = \frac{1}{\sqrt{3}} \approx 0.57735$
    \begin{align*}
        \tan 30^\circ - \mu_s &= 0.57735 - 0.2 = 0.37735 \\
        1 + \mu_s \tan 30^\circ &= 1 + (0.2 \times 0.57735) \\
                                &= 1 + 0.11547 = 1.11547
    \end{align*}
    
    অনুপাতের মান:
    \[ \frac{0.37735}{1.11547} \approx 0.33829 \]
    
    এখন বেগের সমীকরণে মানগুলো বসালে:
    \begin{align*}
        r g &= 50 \times 9.8 = 490\text{ m}^2\text{s}^{-2} \\
        v_{\min} &= \sqrt{490 \times 0.33829} \\
                 &= \sqrt{165.76} \\
                 &\approx \mathbf{12.87}\text{ ms}^{-1}
    \end{align*}
\end{tcolorbox}

% ------------------------------------------
% QUESTION SECTION 25
% ------------------------------------------
\begin{tcolorbox}[questionbox={প্রশ্ন (Question) 25}]
    \noindent
    $5\text{ kg}$ ভরের একটি পাথরকে ভূমি হতে খাড়া উপরের দিকে নিক্ষেপ করা হলো। পাথরটির সম্পূর্ণ গতিপথ জুড়ে বাতাসের ঘর্ষণের কারণে $10\text{ N}$ ধ্রুবক বাধাদানকারী বল ক্রিয়াশীল থাকে। সর্বোচ্চ উচ্চতায় ওঠার সময়কাল $t_u$ এবং সর্বোচ্চ উচ্চতা থেকে পতনের সময়কাল $t_d$ হলে, $\frac{t_u}{t_d}$ অনুপাতটি কত হবে? ($g = 10\text{ ms}^{-2}$)
    
    \vspace{0.8em}
    \begin{english}
    \noindent\textit{\color{darkgray}
    A stone of mass $5\text{ kg}$ is thrown vertically upwards from the ground. Throughout its upward and downward motion, a constant resistive force of $10\text{ N}$ acts on it due to air friction. If $t_u$ is the time of ascent to reach the maximum height and $t_d$ is the time of descent back to the initial point, what is the ratio $\frac{t_u}{t_d}$? ($g = 10\text{ ms}^{-2}$)
    }
    \end{english}
\end{tcolorbox}

% ------------------------------------------
% HIGH-QUALITY TIKZ DIAGRAM 25
% ------------------------------------------
\vspace{1em}
\begin{center}
\begin{tikzpicture}[>=Stealth, scale=1.2]
    % Ground
    \draw[thick, gray] (-3.5, 0) -- (3.5, 0);
    \fill[pattern=north east lines, pattern color=gray!50] (-3.5, -0.2) rectangle (3.5, 0);
    
    % Ascent (Left Side)
    \node[font=\bfseries, primary] at (-1.8, 3.8) {উপরে ওঠা (Ascent)};
    \draw[dashed, gray] (-1.8, 0) -- (-1.8, 3.5);
    \draw[thick, fill=orange!80!black] (-1.8, 1.5) circle (0.25);
    \node[right, font=\bfseries] at (-1.5, 1.5) {$5\text{ kg}$};
    
    % Ascent Vectors
    \draw[->, line width=1.5pt, blue] (-1.8, 1.8) -- (-1.8, 2.8) node[above] {$v$ (বেগ)};
    \draw[->, line width=1.5pt, red] (-1.95, 1.25) -- (-1.95, 0.2) node[left] {$mg$};
    \draw[->, line width=1.5pt, red!50!black] (-1.65, 1.25) -- (-1.65, 0.5) node[right] {$F_{\text{air}}$};
    
    % Descent (Right Side)
    \node[font=\bfseries, primary] at (1.8, 3.8) {নিচে নামা (Descent)};
    \draw[dashed, gray] (1.8, 0) -- (1.8, 3.5);
    \draw[thick, fill=orange!80!black] (1.8, 2.5) circle (0.25);
    \node[right, font=\bfseries] at (2.1, 2.5) {$5\text{ kg}$};
    
    % Descent Vectors
    \draw[->, line width=1.5pt, blue] (1.8, 2.2) -- (1.8, 1.2) node[below] {$v$ (বেগ)};
    \draw[->, line width=1.5pt, red] (1.65, 2.2) -- (1.65, 1.15) node[left] {$mg$};
    \draw[->, line width=1.5pt, red!50!black] (1.95, 2.75) -- (1.95, 3.5) node[right] {$F_{\text{air}}$};

\end{tikzpicture}
\end{center}
\vspace{1em}

% ------------------------------------------
% OPTIONS & ANSWER 25
% ------------------------------------------
\begin{itemize}[label={}, leftmargin=1cm, itemsep=0.5em]
    \item[\textbf{A)}] সময়ের অনুপাত $\sqrt{2} : \sqrt{3}$
    \item[\textbf{B)}] সময়ের অনুপাত $\sqrt{3} : \sqrt{2}$
    \item[\textbf{C)}] সময়ের অনুপাত $2 : 3$
    \item[\textbf{D)}] সময়ের অনুপাত $1 : 1$
\end{itemize}

\vspace{0.5em}
\begin{tcolorbox}[answerbox]
    \textbf{\color{success} উত্তর:} A) সময়ের অনুপাত $\sqrt{2} : \sqrt{3}$
\end{tcolorbox}
\vspace{1em}

% ------------------------------------------
% SOLUTION SECTION 25
% ------------------------------------------
\begin{tcolorbox}[solutionbox={সমাধান (Solution)}]
    \noindent
    \textbf{ধাপ ১: উপরে ওঠার ক্ষেত্রে ত্বরণ (মন্দন) নির্ণয়} \\
    উপরে ওঠার সময় অভিকর্ষ বল এবং বাতাসের বাধা উভয়ই নিচের দিকে ক্রিয়া করে। ফলে মোট নিম্নমুখী বল:
    \begin{align*}
        F_{\text{net}, u} &= mg + F_{\text{air}} \\
        &= (5 \times 10) + 10 \\
        &= 50 + 10 = 60\text{ N}
    \end{align*}
    উপরে ওঠার মন্দন ($a_u$):
    \[ a_u = \frac{F_{\text{net}, u}}{m} = \frac{60}{5} = 12\text{ ms}^{-2} \]
    
    \vspace{0.5em}
    \noindent\textbf{ধাপ ২: নিচে নামার ক্ষেত্রে ত্বরণ নির্ণয়} \\
    নিচে নামার সময় অভিকর্ষ নিচের দিকে কিন্তু বাতাসের বাধা গতির বিপরীতে অর্থাৎ উপরের দিকে কাজ করে। ফলে মোট নিম্নমুখী বল:
    \begin{align*}
        F_{\text{net}, d} &= mg - F_{\text{air}} \\
        &= (5 \times 10) - 10 \\
        &= 50 - 10 = 40\text{ N}
    \end{align*}
    নিচে নামার ত্বরণ ($a_d$):
    \[ a_d = \frac{F_{\text{net}, d}}{m} = \frac{40}{5} = 8\text{ ms}^{-2} \]
    
    \vspace{0.5em}
    \noindent\textbf{ধাপ ৩: সময়ের অনুপাত নির্ণয়} \\
    উভয় ক্ষেত্রে অতিক্রান্ত উচ্চতা $H$ সমান। স্থির অবস্থা থেকে শুরু ও শেষের ক্ষেত্রে সমীকরণ:
    \[ H = \frac{1}{2} a_u t_u^2 = \frac{1}{2} a_d t_d^2 \]
    এখান থেকে সময়ের বর্গের অনুপাত:
    \begin{align*}
        \frac{t_u^2}{t_d^2} &= \frac{a_d}{a_u} \\
        &= \frac{8}{12} = \frac{2}{3}
    \end{align*}
    অতএব, সময়ের অনুপাত:
    \[ \frac{t_u}{t_d} = \sqrt{\frac{2}{3}} = \mathbf{\frac{\sqrt{2}}{\sqrt{3}}} \]
\end{tcolorbox}

\newpage

% ------------------------------------------
% QUESTION SECTION 26
% ------------------------------------------
\begin{tcolorbox}[questionbox={প্রশ্ন (Question) 26}]
    \noindent
    একটি অনুভূমিক মসৃণ মেঝের ওপর $M$ ভরের একটি কিলক (wedge) রাখা আছে যার আনত তলের নতিকোণ $\theta$। কিলকটির ওপর $m$ ভরের একটি ছোট ব্লককে রাখা হলো। কিলকটির ওপর অনুভূমিক বরাবর কত মানের বাহ্যিক বল $F$ প্রয়োগ করলে ব্লকটি আনত তলের সাপেক্ষে স্থির থাকবে (নিচে বা উপরে পিছলে যাবে না)? (কিলক ও ব্লকের মধ্যবর্তী তল সম্পূর্ণ ঘর্ষণহীন)
    
    \vspace{0.8em}
    \begin{english}
    \noindent\textit{\color{darkgray}
    A wedge of mass $M$ with an inclination angle $\theta$ is placed on a smooth horizontal floor. A small block of mass $m$ is placed on the inclined surface of the wedge. What horizontal external force $F$ must be applied to the wedge so that the block remains stationary relative to the inclined plane? (The interface between the wedge and the block is frictionless)
    }
    \end{english}
\end{tcolorbox}

% ------------------------------------------
% HIGH-QUALITY TIKZ DIAGRAM 26
% ------------------------------------------
\vspace{1em}
\begin{center}
\begin{tikzpicture}[>=Stealth, scale=1.3]
    \def\mytheta{35} % Angle for drawing
    
    % Ground
    \draw[thick, gray] (-2.5, 0) -- (4, 0);
    \fill[pattern=north east lines, pattern color=gray!50] (-2.5, -0.2) rectangle (4, 0);
    
    % Wedge
    \draw[thick, fill=cyan!15] (0, 0) -- (3, 0) -- (0, {3*tan(\mytheta)}) -- cycle;
    \node[font=\bfseries] at (1, 0.5) {$M$};
    
    % Angle Theta
    \draw (2.5, 0) arc (180:{180-\mytheta}:0.5);
    \node at (2.1, 0.2) {$\theta$};
    
    % Force F
    \draw[->, line width=1.8pt, primary] (-1.5, 1) -- (0, 1) node[midway, above] {$F$};
    \draw[->, thick, primary] (3.2, 0.5) -- (4, 0.5) node[midway, above] {$a$};
    
    % Block m
    \coordinate (Block) at (1.5, {1.5*tan(\mytheta)});
    \begin{scope}[shift={(Block)}, rotate=-\mytheta]
        \draw[thick, fill=orange!80!black] (-0.4, 0) rectangle (0.4, 0.5);
        \node[white, font=\bfseries] at (0, 0.25) {$m$};
    \end{scope}
    
    % Forces on block m (FBD)
    \draw[->, line width=1.5pt, red] (Block) ++(0, 0.25) -- ++(0, -1.5) node[below] {$mg$};
    \draw[->, line width=1.5pt, purple] (Block) ++(0, 0.25) -- ++(-1.5, 0) node[above left] {$ma$ (Pseudo)};
    \draw[->, line width=1.5pt, green!60!black] (Block) ++(0, 0.25) -- ++({1.5*sin(\mytheta)}, {1.5*cos(\mytheta)}) node[right] {$N$};
\end{tikzpicture}
\end{center}
\vspace{1em}

% ------------------------------------------
% OPTIONS & ANSWER 26
% ------------------------------------------
\begin{itemize}[label={}, leftmargin=1cm, itemsep=0.5em]
    \item[\textbf{A)}] প্রযুক্ত বল $F = (M + m)g \tan\theta$
    \item[\textbf{B)}] প্রযুক্ত বল $F = (M + m)g \cot\theta$
    \item[\textbf{C)}] প্রযুক্ত বল $F = mg \sin\theta$
    \item[\textbf{D)}] প্রযুক্ত বল $F = M g \tan\theta$
\end{itemize}

\vspace{0.5em}
\begin{tcolorbox}[answerbox]
    \textbf{\color{success} উত্তর:} A) প্রযুক্ত বল $F = (M + m)g \tan\theta$
\end{tcolorbox}
\vspace{1em}

% ------------------------------------------
% SOLUTION SECTION 26
% ------------------------------------------
\begin{tcolorbox}[solutionbox={সমাধান (Solution)}]
    \noindent
    \textbf{ধাপ ১: সিস্টেমের ত্বরণ নির্ণয়} \\
    যেহেতু ব্লকটি কিলকের সাপেক্ষে স্থির, তাই কিলক এবং ব্লক একত্রে একই অনুভূমিক ত্বরণ $a$ নিয়ে গতিশীল হবে। পুরো সিস্টেমের (ভর $= M + m$) ওপর নিউটনের দ্বিতীয় সূত্র প্রয়োগ করে পাই:
    \[ F = (M + m)a \]
    
    \vspace{0.5em}
    \noindent\textbf{ধাপ ২: ব্লকের মুক্ত বস্তু চিত্র (FBD) বিশ্লেষণ} \\
    ত্বরণশীল কিলকের সাপেক্ষে (অজড়ীয় প্রসঙ্গ কাঠামো) ব্লকের ওপর একটি অনুভূমিক বিপরীতমুখী অবাস্তব বল (Pseudo force) $ma$ ক্রিয়া করবে। 
    
    ব্লকের ওপর ক্রিয়াশীল বলসমূহ হলো:
    \begin{enumerate}[topsep=0pt, itemsep=0pt]
        \item ওজন $mg$ (খাড়া নিচের দিকে)
        \item অবাস্তব বল $ma$ (অনুভূমিক বামদিকে)
        \item অভিলম্ব প্রতিক্রিয়া $N$ (তলের ওপর লম্বভাবে)
    \end{enumerate}
    
    \vspace{0.5em}
    \noindent\textbf{ধাপ ৩: তলের সমান্তরালে বলের সাম্যাবস্থা} \\
    ব্লকটি আনত তলের সমান্তরালে স্থির থাকতে হলে, তলের সমান্তরাল বরাবর বলের উপাংশগুলোর লব্ধি শূন্য হতে হবে। 
    $mg$-এর তলের সমান্তরাল উপাংশ (নিচের দিকে) $= mg \sin\theta$ 
    $ma$-এর তলের সমান্তরাল উপাংশ (উপরের দিকে) $= ma \cos\theta$
    
    সাম্যাবস্থার শর্তানুসারে:
    \begin{align*}
        mg \sin\theta &= ma \cos\theta \\
        a &= \frac{g \sin\theta}{\cos\theta} = g \tan\theta
    \end{align*}
    
    \vspace{0.5em}
    \noindent\textbf{ধাপ ৪: প্রয়োজনীয় বাহ্যিক বল $F$} \\
    ত্বরণের মান $F$-এর সমীকরণে বসিয়ে পাই:
    \begin{align*}
        F &= (M + m) a \\
          &= \mathbf{(M + m) g \tan\theta}
    \end{align*}
\end{tcolorbox}

\newpage

% ------------------------------------------
% QUESTION SECTION 27
% ------------------------------------------
\begin{tcolorbox}[questionbox={প্রশ্ন (Question) 27}]
    \noindent
    একটি ঘর্ষণহীন অনুভূমিক সমতলে $m_1 = 4\text{ kg}$ ও $m_2 = 2\text{ kg}$ ভরের দুটি ব্লক একটি ভরহীন স্প্রিং দ্বারা যুক্ত। ব্লক দুটিকে বিপরীত দিকে টেনে ছেড়ে দেওয়া হলো। কোনো এক মুহূর্তে যদি $4\text{ kg}$ ভরের ব্লকের ত্বরণ ডানদিকে $3\text{ ms}^{-2}$ হয়, তবে ঐ একই মুহূর্তে $2\text{ kg}$ ভরের ব্লকের ত্বরণ কত হবে?
    
    \vspace{0.8em}
    \begin{english}
    \noindent\textit{\color{darkgray}
    Two blocks of masses $m_1 = 4\text{ kg}$ and $m_2 = 2\text{ kg}$ on a frictionless horizontal plane are connected by a massless spring. The blocks are pulled apart and released. If at a certain instant the acceleration of the $4\text{ kg}$ block is $3\text{ ms}^{-2}$ towards the right, what is the acceleration of the $2\text{ kg}$ block at that same instant?
    }
    \end{english}
\end{tcolorbox}

% ------------------------------------------
% HIGH-QUALITY TIKZ DIAGRAM 27
% ------------------------------------------
\vspace{1em}
\begin{center}
\begin{tikzpicture}[>=Stealth, scale=1.3]
    % Floor
    \draw[thick, gray] (-4, 0) -- (4, 0);
    \fill[pattern=north east lines, pattern color=gray!50] (-4, -0.2) rectangle (4, 0);
    \node[below, font=\footnotesize, text=gray!80!black] at (0, -0.3) {ঘর্ষণহীন তল (Frictionless)};
    
    % Block m2 (2 kg) on the left
    \draw[thick, fill=cyan!70!black, rounded corners=2pt] (-2.5, 0) rectangle (-1.2, 0.8);
    \node[white, font=\bfseries] at (-1.85, 0.4) {$2\text{ kg}$};
    \node[above, font=\bfseries] at (-1.85, 0.8) {$m_2$};
    
    % Block m1 (4 kg) on the right
    \draw[thick, fill=orange!80!black, rounded corners=2pt] (1.2, 0) rectangle (2.8, 1.0);
    \node[white, font=\bfseries] at (2.0, 0.5) {$4\text{ kg}$};
    \node[above, font=\bfseries] at (2.0, 1.0) {$m_1$};
    
    % Spring (Stretched)
    \draw[thick, decoration={coil, aspect=0.4, segment length=6pt, amplitude=5pt}, decorate] (-1.2, 0.4) -- (1.2, 0.4);
    
    % Spring Force
    \draw[->, thick, red] (-1.2, 0.6) -- (-0.2, 0.6) node[above] {$F_s$};
    \draw[->, thick, red] (1.2, 0.6) -- (0.2, 0.6) node[above] {$F_s$};
    
    % Acceleration Vectors
    \draw[->, line width=1.5pt, primary] (2.9, 0.5) -- (4.2, 0.5) node[above] {$a_1 = 3\text{ ms}^{-2}$};
    \draw[->, line width=1.5pt, primary] (-2.6, 0.4) -- (-3.8, 0.4) node[above] {$a_2 = ?$};
\end{tikzpicture}
\end{center}
\vspace{1em}

% ------------------------------------------
% OPTIONS & ANSWER 27
% ------------------------------------------
\begin{itemize}[label={}, leftmargin=1cm, itemsep=0.5em]
    \item[\textbf{A)}] ত্বরণ বামদিকে $6\text{ ms}^{-2}$
    \item[\textbf{B)}] ত্বরণ ডানদিকে $6\text{ ms}^{-2}$
    \item[\textbf{C)}] ত্বরণ বামদিকে $3\text{ ms}^{-2}$
    \item[\textbf{D)}] ত্বরণ বামদিকে $1.5\text{ ms}^{-2}$
\end{itemize}

\vspace{0.5em}
\begin{tcolorbox}[answerbox]
    \textbf{\color{success} উত্তর:} A) ত্বরণ বামদিকে $6\text{ ms}^{-2}$
\end{tcolorbox}
\vspace{1em}

% ------------------------------------------
% SOLUTION SECTION 27
% ------------------------------------------
\begin{tcolorbox}[solutionbox={সমাধান (Solution)}]
    \noindent
    \textbf{ধাপ ১: ভরকেন্দ্রের ত্বরণ নীতি} \\
    স্প্রিং-ব্লক ব্যবস্থায় অনুভূমিক বরাবর কোনো বাহ্যিক বল ক্রিয়াশীল নেই (অর্থাৎ $F_{\text{ext}} = 0$)। স্প্রিং-এর টান বলটি সম্পূর্ণ অভ্যন্তরীণ। ফলে ব্যবস্থার ভরকেন্দ্রের ত্বরণ সর্বদা শূন্য হবে:
    \[ a_{\text{cm}} = \frac{m_1 \vec{a}_1 + m_2 \vec{a}_2}{m_1 + m_2} = 0 \]
    
    \vspace{0.5em}
    \noindent\textbf{ধাপ ২: ত্বরণের সম্পর্ক ও মান নির্ণয়} \\
    উপরের সমীকরণ থেকে পাই:
    \begin{align*}
        m_1 \vec{a}_1 + m_2 \vec{a}_2 &= 0 \\
        m_2 \vec{a}_2 &= -m_1 \vec{a}_1 \\
        \vec{a}_2 &= -\frac{m_1}{m_2} \vec{a}_1
    \end{align*}
    
    প্রশ্নমতে, $4\text{ kg}$ ভরের ব্লকের ত্বরণ ডানদিকে, অর্থাৎ $\vec{a}_1 = +3\hat{i}\text{ ms}^{-2}$। 
    মানগুলো বসালে:
    \begin{align*}
        \vec{a}_2 &= -\frac{4}{2} \times (+3\hat{i}) \\
        &= -2 \times (+3\hat{i}) \\
        &= -6\hat{i}\text{ ms}^{-2}
    \end{align*}
    
    ঋণাত্মক চিহ্ন নির্দেশ করে যে $2\text{ kg}$ ভরের ব্লকের ত্বরণ বিপরীত দিকে অর্থাৎ \textbf{বামদিকে $6\text{ ms}^{-2}$} হবে। 
    \textit{(বিকল্পভাবে, স্প্রিং-এর পুনরুদ্ধারকারী বল $F_s$ উভয়ের ওপর সমান কিন্তু বিপরীতমুখী। $m_1$ এর ওপর $F_s = 4 \times 3 = 12\text{ N}$ বামদিকে, ফলে $m_2$ এর ওপর $12\text{ N}$ বামদিকে। $a_2 = 12 / 2 = 6\text{ ms}^{-2}$ বামদিকে।)}
\end{tcolorbox}

% ------------------------------------------
% QUESTION SECTION 28
% ------------------------------------------
\begin{tcolorbox}[questionbox={প্রশ্ন (Question) 28}]
    \noindent
    $3\text{ kg}$ ভরের একটি কণার অবস্থান ভেক্টর সময়ের অপেক্ষক হিসেবে $\vec{r}(t) = (2t^2)\hat{i} + (3t)\hat{j} - (5t^3)\hat{k}\text{ m}$ দ্বারা নির্দেশিত হয়। $t = 2\text{ s}$ সময়ে মূলবিন্দুর সাপেক্ষে কণাটির ওপর ক্রিয়াশীল টর্কের মান কত হবে?
    
    \vspace{0.8em}
    \begin{english}
    \noindent\textit{\color{darkgray}
    The position vector of a $3\text{ kg}$ particle as a function of time is given by $\vec{r}(t) = (2t^2)\hat{i} + (3t)\hat{j} - (5t^3)\hat{k}\text{ m}$. What is the magnitude of the torque acting on the particle about the origin at $t = 2\text{ s}$?
    }
    \end{english}
\end{tcolorbox}

% ------------------------------------------
% HIGH-QUALITY TIKZ DIAGRAM 28
% ------------------------------------------
\vspace{1em}
\begin{center}
\begin{tikzpicture}[>=Stealth, scale=1.3]
    % Coordinate axes (3D effect)
    \draw[->, thick, gray] (0,0,0) -- (3,0,0) node[right] {$y$};
    \draw[->, thick, gray] (0,0,0) -- (0,3,0) node[above] {$z$};
    \draw[->, thick, gray] (0,0,0) -- (0,0,4) node[below left] {$x$};
    
    % Origin
    \fill[black] (0,0,0) circle (0.06) node[below right] {$O (0,0,0)$};
    
    % Particle Path (curved line)
    \draw[dashed, darkgray] (0,0,2) .. controls (1,1,1) and (2,1,-1) .. (2.5, 2, -2) node[right] {গতিপথ};
    
    % Position at t=2
    \coordinate (P) at (2, 1.5, -1);
    \fill[orange!80!black] (P) circle (0.08) node[above right] {$m = 3\text{ kg}$};
    
    % Position vector r
    \draw[->, line width=1.2pt, blue] (0,0,0) -- (P) node[midway, above left] {$\vec{r}(2)$};
    
    % Force vector
    \draw[->, line width=1.5pt, red] (P) -- ++(1, -1.5, -0.5) node[right] {$\vec{F}(2)$};
    
    % Torque vector (at origin, roughly perpendicular to r and F)
    \draw[->, line width=1.8pt, green!60!black] (0,0,0) -- (-1.5, 2, 1) node[above left] {$\vec{\tau} = \vec{r} \times \vec{F}$};
\end{tikzpicture}
\end{center}
\vspace{1em}

% ------------------------------------------
% OPTIONS & ANSWER 28
% ------------------------------------------
\begin{itemize}[label={}, leftmargin=1cm, itemsep=0.5em]
    \item[\textbf{A)}] টর্কের মান $36\sqrt{409}\text{ Nm}$
    \item[\textbf{B)}] টর্কের মান $540\text{ Nm}$
    \item[\textbf{C)}] টর্কের মান $18\sqrt{305}\text{ Nm}$
    \item[\textbf{D)}] টর্কের মান $720\text{ Nm}$
\end{itemize}

\vspace{0.5em}
\begin{tcolorbox}[answerbox]
    \textbf{\color{success} উত্তর:} A) টর্কের মান $36\sqrt{409}\text{ Nm}$
\end{tcolorbox}
\vspace{1em}

% ------------------------------------------
% SOLUTION SECTION 28
% ------------------------------------------
\begin{tcolorbox}[solutionbox={সমাধান (Solution)}]
    \noindent
    \textbf{ধাপ ১: বেগ, ত্বরণ এবং বল নির্ণয়} \\
    সময়ের সাপেক্ষে কণার অবস্থান ভেক্টর:
    \[ \vec{r}(t) = (2t^2)\hat{i} + (3t)\hat{j} - (5t^3)\hat{k} \]
    
    কণার বেগ:
    \[ \vec{v}(t) = \frac{d\vec{r}}{dt} = (4t)\hat{i} + 3\hat{j} - (15t^2)\hat{k} \]
    
    ত্বরণ:
    \[ \vec{a}(t) = \frac{d\vec{v}}{dt} = 4\hat{i} + 0\hat{j} - (30t)\hat{k} \]
    
    নিউটনের দ্বিতীয় সূত্র অনুসারে কণাটির ওপর বল ($m = 3\text{ kg}$):
    \begin{align*}
        \vec{F}(t) &= m\vec{a}(t) \\
        &= 3 \left[ 4\hat{i} - (30t)\hat{k} \right] \\
        &= 12\hat{i} - (90t)\hat{k}
    \end{align*}
    
    \vspace{0.5em}
    \noindent\textbf{ধাপ ২: $t = 2\text{ s}$ সময়ে ভেক্টরসমূহের মান} \\
    $t = 2\text{ s}$ সময়ে অবস্থান ভেক্টর:
    \begin{align*}
        \vec{r}(2) &= 2(2)^2\hat{i} + 3(2)\hat{j} - 5(2)^3\hat{k} \\
        &= 8\hat{i} + 6\hat{j} - 40\hat{k}
    \end{align*}
    
    $t = 2\text{ s}$ সময়ে বল ভেক্টর:
    \begin{align*}
        \vec{F}(2) &= 12\hat{i} + 0\hat{j} - 90(2)\hat{k} \\
        &= 12\hat{i} - 180\hat{k}
    \end{align*}
    
    \vspace{0.5em}
    \noindent\textbf{ধাপ ৩: টর্ক নির্ণয়} \\
    মূলবিন্দুর সাপেক্ষে টর্ক, $\vec{\tau} = \vec{r} \times \vec{F}$:
    \begin{align*}
        \vec{\tau} &= \begin{vmatrix} 
        \hat{i} & \hat{j} & \hat{k} \\ 
        8 & 6 & -40 \\ 
        12 & 0 & -180 
        \end{vmatrix} \\
        &= \hat{i}\big[6(-180) - (-40)(0)\big] - \hat{j}\big[8(-180) - (-40)(12)\big] + \hat{k}\big[8(0) - 6(12)\big] \\
        &= \hat{i}(-1080 - 0) - \hat{j}(-1440 + 480) + \hat{k}(0 - 72) \\
        &= -1080\hat{i} - (-960)\hat{j} - 72\hat{k} \\
        &= -1080\hat{i} + 960\hat{j} - 72\hat{k}
    \end{align*}
    
    টর্কের মান:
    \begin{align*}
        |\vec{\tau}| &= \sqrt{(-1080)^2 + (960)^2 + (-72)^2} \\
        &= \sqrt{1166400 + 921600 + 5184} \\
        &= \sqrt{2093184} \\
        &= \mathbf{36\sqrt{409}}\text{ Nm} \approx 728\text{ Nm}
    \end{align*}
\end{tcolorbox}

\newpage

% ------------------------------------------
% QUESTION SECTION 29
% ------------------------------------------
\begin{tcolorbox}[questionbox={প্রশ্ন (Question) 29}]
    \noindent
    $m$ ভরের একটি গতিশীল বস্তু $2m$ ভরের একটি স্থির বস্তুর সাথে পূর্ণ অস্থিতিস্থাপক (completely inelastic) মুখোমুখি সংঘর্ষে লিপ্ত হলো। এই সংঘর্ষের ফলে আদি গতিশক্তির শতকরা কত অংশ তাপ ও শব্দ শক্তিতে রূপান্তরিত বা অপচয় হবে?
    
    \vspace{0.8em}
    \begin{english}
    \noindent\textit{\color{darkgray}
    A moving body of mass $m$ undergoes a completely inelastic one-dimensional head-on collision with a stationary body of mass $2m$. What percentage of the initial kinetic energy is dissipated into heat and sound during this collision?
    }
    \end{english}
\end{tcolorbox}

% ------------------------------------------
% HIGH-QUALITY TIKZ DIAGRAM 29
% ------------------------------------------
\vspace{1em}
\begin{center}
\begin{tikzpicture}[>=Stealth, scale=1.2]
    % Before collision
    \node[font=\bfseries, primary] at (0, 1.5) {সংঘর্ষের পূর্বে (Before Collision)};
    \draw[thick, gray] (-3, 0) -- (3, 0);
    \draw[thick, fill=orange!80!black] (-2, 0) rectangle (-1, 0.8);
    \node[white, font=\bfseries] at (-1.5, 0.4) {$m$};
    \draw[->, line width=1.5pt, primary] (-1, 0.4) -- (0, 0.4) node[above] {$u$};

    \draw[thick, fill=cyan!70!black] (1, 0) rectangle (2.5, 1.2);
    \node[white, font=\bfseries] at (1.75, 0.6) {$2m$};
    \node[above] at (1.75, 1.2) {$v=0$};

    % After collision
    \begin{scope}[shift={(0, -3)}]
        \node[font=\bfseries, primary] at (0, 1.5) {সংঘর্ষের পরে (After Collision)};
        \draw[thick, gray] (-3, 0) -- (3, 0);
        \draw[thick, fill=purple!70!black] (-0.75, 0) rectangle (1.25, 1.2);
        \node[white, font=\bfseries] at (0.25, 0.6) {$3m$};
        \draw[->, line width=1.5pt, primary] (1.25, 0.6) -- (2.5, 0.6) node[above] {$V = u/3$};

        % Energy loss representation
        \draw[->, decorate, decoration={snake, amplitude=2pt, segment length=5pt}, red] (0.25, 1.2) -- (0.25, 2.2) node[above, font=\bfseries, text=red!80!black] {তাপ ও শব্দ (Energy Loss)};
    \end{scope}
\end{tikzpicture}
\end{center}
\vspace{1em}

% ------------------------------------------
% OPTIONS & ANSWER 29
% ------------------------------------------
\begin{itemize}[label={}, leftmargin=1cm, itemsep=0.5em]
    \item[\textbf{A)}] গতিশক্তির অপচয় $66.67\%$
    \item[\textbf{B)}] গতিশক্তির অপচয় $33.33\%$
    \item[\textbf{C)}] গতিশক্তির অপচয় $50.00\%$
    \item[\textbf{D)}] গতিশক্তির অপচয় $75.00\%$
\end{itemize}

\vspace{0.5em}
\begin{tcolorbox}[answerbox]
    \textbf{\color{success} উত্তর:} A) গতিশক্তির অপচয় $66.67\%$
\end{tcolorbox}
\vspace{1em}

% ------------------------------------------
% SOLUTION SECTION 29
% ------------------------------------------
\begin{tcolorbox}[solutionbox={সমাধান (Solution)}]
    \noindent
    \textbf{ধাপ ১: আদি গতিশক্তি ও ভরবেগ সংরক্ষণ} \\
    সংঘর্ষের পূর্বে মোট আদি গতিশক্তি:
    \[ K_i = \frac{1}{2} m u^2 \]
    
    যেহেতু সংঘর্ষটি পূর্ণ অস্থিতিস্থাপক (completely inelastic), তাই বস্তুদুটি সংঘর্ষের পর একত্রে একই সাধারণ বেগ $V$-তে চলতে থাকবে। রৈখিক ভরবেগের সংরক্ষণ সূত্র অনুসারে:
    \begin{align*}
        m u + 2m (0) &= (m + 2m) V \\
        m u &= 3m V \\
        V &= \frac{u}{3}
    \end{align*}
    
    \vspace{0.5em}
    \noindent\textbf{ধাপ ২: শেষ গতিশক্তি ও শক্তির অপচয়} \\
    সংঘর্ষের পর শেষ গতিশক্তি:
    \begin{align*}
        K_f &= \frac{1}{2} (3m) V^2 \\
            &= \frac{1}{2} (3m) \left(\frac{u}{3}\right)^2 \\
            &= \frac{1}{2} (3m) \frac{u^2}{9} \\
            &= \frac{1}{6} m u^2 = \frac{1}{3} \left(\frac{1}{2} m u^2\right) \\
            &= \frac{1}{3} K_i
    \end{align*}
    
    গতিশক্তির অপচয়ের পরিমাণ (যা তাপ ও শব্দে রূপান্তরিত হয়):
    \begin{align*}
        \Delta K &= K_i - K_f \\
                 &= K_i - \frac{1}{3} K_i = \frac{2}{3} K_i
    \end{align*}
    
    \vspace{0.5em}
    \noindent\textbf{ধাপ ৩: শতকরা হার} \\
    শতকরা অপচয়ের পরিমাণ:
    \begin{align*}
        \frac{\Delta K}{K_i} \times 100\% &= \frac{2}{3} \times 100\% \\
        &= 66.666...\% \\
        &\approx \mathbf{66.67\%}
    \end{align*}
\end{tcolorbox}

\newpage

% ------------------------------------------
% QUESTION SECTION 30
% ------------------------------------------
\begin{tcolorbox}[questionbox={প্রশ্ন (Question) 30}]
    \noindent
    একটি সরল দোলকের ববের ভর $m$ এবং সুতার দৈর্ঘ্য $L$। ববটিকে উলম্ব অবস্থান থেকে $60^\circ$ কোণে টেনে ছেড়ে দেওয়া হলো। দোলনের সর্বনিম্ন বিন্দু অতিক্রম করার মুহূর্তে সুতায় টান বল কত হবে? ($g$ হলো অভিকর্ষজ ত্বরণ)
    
    \vspace{0.8em}
    \begin{english}
    \noindent\textit{\color{darkgray}
    A simple pendulum consists of a bob of mass $m$ and a string of length $L$. The bob is pulled aside to make an angle of $60^\circ$ with the vertical and released from rest. What is the tension in the string as the bob passes through its lowest equilibrium position? ($g$ is the acceleration due to gravity)
    }
    \end{english}
\end{tcolorbox}

% ------------------------------------------
% HIGH-QUALITY TIKZ DIAGRAM 30
% ------------------------------------------
\vspace{1em}
\begin{center}
\begin{tikzpicture}[>=Stealth, scale=1.3]
    % Ceiling
    \draw[thick] (-1.5, 0) -- (1.5, 0);
    \fill[pattern=north east lines, pattern color=gray!50] (-1.5, 0) rectangle (1.5, 0.2);
    \coordinate (O) at (0,0);
    \fill[black] (O) circle (0.05);

    \def\myL{3.5}
    \def\myAng{60}

    % Lowest point
    \coordinate (P0) at (0, -\myL);
    % Initial point
    \coordinate (P1) at ({\myL*sin(\myAng)}, {-\myL*cos(\myAng)});

    % Arc Path
    \draw[dashed, gray!80, thick] (P1) arc ({-90+\myAng}:-90:\myL);

    % Strings
    \draw[line width=1pt, darkgray] (O) -- (P1) node[midway, right=2pt] {$L$};
    \draw[line width=1pt, darkgray] (O) -- (P0);

    % Angle
    \draw[dashed, darkgray] (O) -- (0, -2);
    \draw[thick, blue] (0, -1) arc (-90:{-90+\myAng}:1);
    \node[blue, font=\bfseries] at (15:1.2) [rotate=-75] {$60^\circ$};

    % Bobs
    \draw[thick, fill=orange!80!black, opacity=0.4] (P1) circle (0.15) node[above right, opacity=1, text=black] {স্থিরাবস্থা ($v=0$)};
    \draw[thick, fill=orange!80!black] (P0) circle (0.15);

    % Forces at lowest point P0
    \draw[->, line width=1.5pt, red] (P0) -- ++(0, -1.2) node[below] {$mg$};
    \draw[->, line width=1.5pt, primary] (P0) -- ++(0, 1.8) node[right] {$T = ?$};
    \draw[->, line width=1.5pt, green!60!black] (P0) ++(0.3, 0) -- ++(1, 0) node[right] {$v$};

    % Height difference h
    \draw[dashed, gray] (P1) -- (0, {-\myL*cos(\myAng)});
    \draw[<->, dashed, thick] (-0.4, {-\myL*cos(\myAng)}) -- (-0.4, -\myL) node[midway, left] {$h = \frac{L}{2}$};
\end{tikzpicture}
\end{center}
\vspace{1em}

% ------------------------------------------
% OPTIONS & ANSWER 30
% ------------------------------------------
\begin{itemize}[label={}, leftmargin=1cm, itemsep=0.5em]
    \item[\textbf{A)}] সুতার টান $2mg$
    \item[\textbf{B)}] সুতার টান $3mg$
    \item[\textbf{C)}] সুতার টান $mg$
    \item[\textbf{D)}] সুতার টান $1.5mg$
\end{itemize}

\vspace{0.5em}
\begin{tcolorbox}[answerbox]
    \textbf{\color{success} উত্তর:} A) সুতার টান $2mg$
\end{tcolorbox}
\vspace{1em}

% ------------------------------------------
% SOLUTION SECTION 30
% ------------------------------------------
\begin{tcolorbox}[solutionbox={সমাধান (Solution)}]
    \noindent
    \textbf{ধাপ ১: পতনশীল উচ্চতা নির্ণয়} \\
    উল্লম্বের সাথে $\theta = 60^\circ$ কোণে অবস্থানকালে ববটির উলম্ব উচ্চতা (সর্বনিম্ন বিন্দু সাপেক্ষে):
    \begin{align*}
        h &= L - L \cos 60^\circ \\
          &= L(1 - \cos 60^\circ) \\
          &= L\left(1 - \frac{1}{2}\right) = \frac{L}{2}
    \end{align*}
    
    \vspace{0.5em}
    \noindent\textbf{ধাপ ২: সর্বনিম্ন বিন্দুতে বেগ নির্ণয়} \\
    শক্তির সংরক্ষণ সূত্রানুসারে, সর্বোচ্চ অবস্থানে সঞ্চিত বিভব শক্তি সর্বনিম্ন বিন্দুতে গতিশক্তিতে রূপান্তরিত হবে:
    \begin{align*}
        \frac{1}{2} m v^2 &= mgh \\
        \frac{1}{2} m v^2 &= mg\left(\frac{L}{2}\right) \\
        \frac{m v^2}{L} &= mg
    \end{align*}
    
    \vspace{0.5em}
    \noindent\textbf{ধাপ ৩: কেন্দ্রমুখী বল ও সুতার টান} \\
    সর্বনিম্ন বিন্দুতে বৃত্তাকার গতির জন্য প্রয়োজনীয় কেন্দ্রমুখী বল প্রদান করে সুতার টান ($T$) এবং ববের ওজন ($mg$)-এর লব্ধি:
    \begin{align*}
        T - mg &= \frac{m v^2}{L}
    \end{align*}
    
    এখানে $\frac{m v^2}{L} = mg$ মান বসিয়ে পাই:
    \begin{align*}
        T &= mg + \frac{m v^2}{L} \\
          &= mg + mg \\
          &= \mathbf{2mg}
    \end{align*}
\end{tcolorbox}

\newpage

% ------------------------------------------
% QUESTION SECTION 31
% ------------------------------------------
\begin{tcolorbox}[questionbox={প্রশ্ন (Question) 31}]
    \noindent
    $60\text{ kg}$ ভরের এক ব্যক্তি $940\text{ kg}$ ভরের একটি লিফটের মধ্যে দাঁড়িয়ে আছেন। লিফটের তারটি ছিঁড়ে গেলে পুরো লিফটটি অভিকর্ষের অধীনে মুক্তভাবে নিচে পড়তে শুরু করল। পতনকালে মেঝে কর্তৃক ব্যক্তির ওপর প্রযুক্ত অভিলম্ব প্রতিক্রিয়া বল কত হবে? ($g = 9.8\text{ ms}^{-2}$)
    
    \vspace{0.8em}
    \begin{english}
    \noindent\textit{\color{darkgray}
    A man of mass $60\text{ kg}$ is standing inside an elevator of mass $940\text{ kg}$. If the supporting cable snaps and the elevator falls freely under gravity, what is the normal reaction force exerted by the elevator floor on the man during the fall? ($g = 9.8\text{ ms}^{-2}$)
    }
    \end{english}
\end{tcolorbox}

% ------------------------------------------
% HIGH-QUALITY TIKZ DIAGRAM 31
% ------------------------------------------
\vspace{1em}
\begin{center}
\begin{tikzpicture}[>=Stealth, scale=1.3]
    % Elevator Box
    \fill[cyan!10] (-1.5, 0) rectangle (1.5, 3.5);
    \draw[line width=2pt, darkgray] (-1.5, 3.5) -- (-1.5, 0) -- (1.5, 0) -- (1.5, 3.5); % Open top for drawing
    
    % Snapped Cable
    \draw[line width=1.5pt, decorate, decoration={zigzag, segment length=4pt, amplitude=1pt}, red] (0, 3.5) -- (0, 4.2) node[above, font=\bfseries] {তার ছিঁড়ে গেছে (Snapped!)};
    
    % Person (Stick figure)
    \draw[thick] (0, 1.8) circle (0.25); % Head
    \draw[thick] (0, 1.55) -- (0, 0.6); % Body
    \draw[thick] (0, 1.2) -- (-0.4, 0.8); % Arm L
    \draw[thick] (0, 1.2) -- (0.4, 0.8); % Arm R
    \draw[thick] (0, 0.6) -- (-0.3, 0); % Leg L
    \draw[thick] (0, 0.6) -- (0.3, 0); % Leg R
    \node[right, font=\bfseries] at (0.3, 1.8) {$m = 60\text{ kg}$};
    
    % Forces on Person
    \draw[->, line width=1.5pt, primary] (-0.5, 0) -- (-0.5, 1) node[left] {$R = ?$ (প্রতিক্রিয়া)};
    \draw[->, line width=1.5pt, red] (0.5, 0.8) -- (0.5, -0.2) node[right] {$mg$};
    
    % Global Elevator Acceleration
    \draw[->, line width=2pt, blue] (-2.2, 2.5) -- (-2.2, 0.5) node[midway, left, align=center, font=\bfseries] {মুক্ত পতন\\$a = g$};
\end{tikzpicture}
\end{center}
\vspace{1em}

% ------------------------------------------
% OPTIONS & ANSWER 31
% ------------------------------------------
\begin{itemize}[label={}, leftmargin=1cm, itemsep=0.5em]
    \item[\textbf{A)}] প্রতিক্রিয়া বল $0\text{ N}$
    \item[\textbf{B)}] প্রতিক্রিয়া বল $588\text{ N}$
    \item[\textbf{C)}] প্রতিক্রিয়া বল $9800\text{ N}$
    \item[\textbf{D)}] প্রতিক্রিয়া বল $1176\text{ N}$
\end{itemize}

\vspace{0.5em}
\begin{tcolorbox}[answerbox]
    \textbf{\color{success} উত্তর:} A) প্রতিক্রিয়া বল $0\text{ N}$
\end{tcolorbox}
\vspace{1em}

% ------------------------------------------
% SOLUTION SECTION 31
% ------------------------------------------
\begin{tcolorbox}[solutionbox={সমাধান (Solution)}]
    \noindent
    \textbf{ধাপ ১: লিফটের ত্বরণ} \\
    যেহেতু লিফটের তারটি ছিঁড়ে গেছে, পুরো লিফটটি অভিকর্ষের অধীনে মুক্তভাবে নিচে পড়ছে। সুতরাং লিফটের নিম্নমুখী ত্বরণ:
    \[ a = g \]
    
    \vspace{0.5em}
    \noindent\textbf{ধাপ ২: ব্যক্তির গতির সমীকরণ (FBD)} \\
    লিফটের মেঝের ওপর দাঁড়িয়ে থাকা ব্যক্তির জন্য নিউটনের দ্বিতীয় সূত্র প্রয়োগ করি। ব্যক্তির ওপর দুটি বল ক্রিয়াশীল:
    \begin{enumerate}[topsep=0pt, itemsep=0pt]
        \item ব্যক্তির ওজন, $mg$ (খাড়া নিচের দিকে)
        \item লিফটের মেঝে কর্তৃক অভিলম্ব প্রতিক্রিয়া বল, $R$ (খাড়া উপরের দিকে)
    \end{enumerate}
    যেহেতু ব্যক্তিটিও লিফটের সাথে $a$ ত্বরণে নিচে নামছেন, তাই উলম্ব গতির সমীকরণ হবে:
    \[ mg - R = ma \]
    
    \vspace{0.5em}
    \noindent\textbf{ধাপ ৩: প্রতিক্রিয়া বল $R$ নির্ণয়} \\
    মুক্ত পতনের কারণে $a = g$ বসিয়ে পাই:
    \begin{align*}
        mg - R &= mg \\
        R &= mg - mg \\
        R &= \mathbf{0}\text{ N}
    \end{align*}
    
    \textit{(অর্থাৎ, ব্যক্তিটি সম্পূর্ণ \textbf{আপাত ওজনহীনতা} (Weightlessness) অনুভব করবেন এবং মেঝে দ্বারা তার ওপর কোনো প্রতিক্রিয়া বল প্রযুক্ত হবে না। লিফটের ভরের ($940\text{ kg}$) এক্ষেত্রে কোনো প্রভাব নেই।)}
\end{tcolorbox}

\newpage

% ------------------------------------------
% QUESTION SECTION 32 (FIXED)
% ------------------------------------------
\begin{tcolorbox}[questionbox={প্রশ্ন (Question) 32}]
    \noindent
    একটি সুষম নিরেট শঙ্কুর (cone) উচ্চতা $H$ এবং ভূমির ব্যাসার্ধ $R$। এর ভূমির সমতল পৃষ্ঠের কেন্দ্রবিন্দু হতে শঙ্কুর ভরকেন্দ্রের উচ্চতা কত হবে?
    
    \vspace{0.8em}
    \begin{english}
    \noindent\textit{\color{darkgray}
    A uniform solid right circular cone has a height $H$ and base radius $R$. What is the distance of the center of mass of the cone from the center of its flat circular base?
    }
    \end{english}
\end{tcolorbox}

% ------------------------------------------
% HIGH-QUALITY TIKZ DIAGRAM 32 (FIXED)
% ------------------------------------------
\vspace{1em}
\begin{center}
\begin{tikzpicture}[>=Stealth, scale=1.3]
    \def\myH{4}
    \def\myR{2}
    \coordinate (Apex) at (0, \myH);
    \coordinate (BaseLeft) at (-\myR, 0);
    \coordinate (BaseRight) at (\myR, 0);
    
    \fill[cyan!15] (0,0) ellipse [x radius=\myR, y radius=0.5];
    \draw[thick, dashed] (\myR, 0) arc (0:180:{\myR} and 0.5);
    \draw[thick] (-\myR, 0) arc (180:360:{\myR} and 0.5);
    
    \fill[cyan!10, opacity=0.8] (BaseLeft) -- (Apex) -- (BaseRight) arc (0:-180:{\myR} and 0.5) -- cycle;
    \draw[thick] (BaseLeft) -- (Apex) -- (BaseRight);
    
    \draw[dash dot, thick, primary] (0, -0.8) -- (0, \myH+0.5) node[above] {$y$-axis};
    \fill[black] (0,0) circle (0.05) node[below right] {Base Center};
    
    % FIXED MATH PARSING HERE
    \def\yDepth{2.5}
    \pgfmathsetmacro{\rEle}{\myR*(\yDepth/\myH)}
    \pgfmathsetmacro{\yEle}{\rEle*0.25}
    
    \draw[fill=orange!40, opacity=0.8] (0, {\myH-\yDepth}) ellipse [x radius=\rEle, y radius=\yEle];
    \draw[<->] (1.2, \myH) -- (1.2, {\myH-\yDepth}) node[midway, right] {$y$};
    \node[left, font=\footnotesize] at (-\rEle, {\myH-\yDepth}) {$dy$};
    \draw[dashed] (0, {\myH-\yDepth}) -- (\rEle, {\myH-\yDepth}) node[midway, below] {$r$};
    
    \coordinate (CM) at (0, {\myH/4});
    \fill[red] (CM) circle (0.06);
    \node[right, text=red, font=\bfseries] at (0.2, {\myH/4}) {CM};
    
    \draw[<->, dashed] (-2.5, 0) -- (-2.5, \myH) node[midway, left, font=\bfseries] {$H$};
    \draw[dashed] (0, \myH) -- (-2.7, \myH);
    \draw[dashed] (0, 0) -- (-2.7, 0);
    \draw[<->, dashed] (0, -0.6) -- (\myR, -0.6) node[midway, below, font=\bfseries] {$R$};
    \draw[dashed] (\myR, 0) -- (\myR, -0.8);
    \draw[<->, line width=1.2pt, red] (-0.3, 0) -- (-0.3, {\myH/4}) node[midway, left] {$h=?$};
\end{tikzpicture}
\end{center}
\vspace{1em}

\begin{itemize}[label={}, leftmargin=1cm, itemsep=0.5em]
    \item[\textbf{A)}] ভরকেন্দ্রের উচ্চতা $\frac{H}{4}$
    \item[\textbf{B)}] ভরকেন্দ্রের উচ্চতা $\frac{H}{3}$
    \item[\textbf{C)}] ভরকেন্দ্রের উচ্চতা $\frac{3H}{8}$
    \item[\textbf{D)}] ভরকেন্দ্রের উচ্চতা $\frac{H}{2}$
\end{itemize}

\vspace{0.5em}
\begin{tcolorbox}[answerbox]
    \textbf{\color{success} উত্তর:} A) ভরকেন্দ্রের উচ্চতা $\frac{H}{4}$
\end{tcolorbox}
\vspace{1em}

\begin{tcolorbox}[solutionbox={সমাধান (Solution)}]
    \noindent
    \textbf{ধাপ ১: ভরকেন্দ্রের সমীকরণ গঠন} \\
    শঙ্কুর শীর্ষবিন্দুকে মূলবিন্দু ধরে শীর্ষ থেকে $y$ গভীরতায় $dy$ পুরুত্বের একটি বৃত্তাকার পাতলা চাকতি বিবেচনা করি। 
    সমরূপ ত্রিভুজের ধারণা থেকে ঐ চাকতির ব্যাসার্ধ $r$ হবে:
    \[ \frac{r}{y} = \frac{R}{H} \implies r = \frac{R}{H} y \]
    
    চাকতিটির আয়তন (যা ভরের সমানুপাতিক):
    \[ dV = \pi r^2 dy = \pi \left(\frac{R^2}{H^2}\right) y^2 \, dy \]
    
    \vspace{0.5em}
    \noindent\textbf{ধাপ ২: সমাকলনের মাধ্যমে ভরকেন্দ্রের অবস্থান} \\
    শীর্ষবিন্দু থেকে ভরকেন্দ্রের দূরত্ব ($y_{\text{cm}}$):
    \begin{align*}
        y_{\text{cm}} &= \frac{\int_0^H y^3 \, dy}{\int_0^H y^2 \, dy} \\
                      &= \frac{\frac{H^4}{4}}{\frac{H^3}{3}} = \frac{3}{4}H
    \end{align*}
    
    \vspace{0.5em}
    \noindent\textbf{ধাপ ৩: ভূমি থেকে উচ্চতা নির্ণয়} \\
    সমতল ভূমি থেকে ভরকেন্দ্রের উল্লম্ব দূরত্ব ($h$):
    \begin{align*}
        h &= H - y_{\text{cm}} \\
          &= H - \frac{3}{4}H = \mathbf{\frac{H}{4}}
    \end{align*}
\end{tcolorbox}

\newpage

% ------------------------------------------
% QUESTION SECTION 33
% ------------------------------------------
\begin{tcolorbox}[questionbox={প্রশ্ন (Question) 33}]
    \noindent
    একটি বন্দুক থেকে প্রতি মিনিটে $60$টি গুলি ছোড়া হচ্ছে। প্রতিটি গুলির ভর $20\text{ g}$ এবং এদের নির্গমন বেগ $400\text{ ms}^{-1}$। বন্দুকটিকে যথাস্থানে স্থির ধরে রাখতে ব্যবহারকারীকে গড়ে কত বল প্রয়োগ করতে হবে?
    
    \vspace{0.8em}
    \begin{english}
    \noindent\textit{\color{darkgray}
    A machine gun fires $60$ bullets per minute. Each bullet has a mass of $20\text{ g}$ and an exit muzzle velocity of $400\text{ ms}^{-1}$. What average force must the operator exert on the gun to hold it steadily in position?
    }
    \end{english}
\end{tcolorbox}

% ------------------------------------------
% HIGH-QUALITY TIKZ DIAGRAM 33
% ------------------------------------------
\vspace{1em}
\begin{center}
\begin{tikzpicture}[>=Stealth, scale=1.3]
    % Gun representation
    \fill[gray!60] (-1, -0.2) rectangle (1.5, 0.2); % Barrel
    \fill[gray!80!black] (-1.5, -0.5) rectangle (-0.8, 0.3); % Stock
    \fill[gray!90!black] (-0.3, -0.5) rectangle (0, -0.2); % Grip
    \draw[thick] (-1.5, -0.5) rectangle (-0.8, 0.3);
    \draw[thick] (-1, -0.2) rectangle (1.5, 0.2);
    
    % Operator Force (Holding it steady)
    \draw[->, line width=1.8pt, red] (-0.5, 0) -- (-2.5, 0) node[midway, above] {$F_{\text{avg}}$};
    
    % Bullets
    \foreach \x/\op in {2.0/1.0, 3.0/0.7, 4.0/0.4} {
        \fill[orange!80!black, opacity=\op] (\x, 0) ellipse [x radius=0.15, y radius=0.05];
        \draw[->, thick, primary, opacity=\op] (\x+0.2, 0) -- (\x+0.8, 0);
    }
    
    % Info Labels
    \node[above, font=\bfseries, primary] at (3.0, 0.3) {$v = 400 \text{ ms}^{-1}$};
    \node[below, font=\bfseries] at (3.0, -0.2) {$m = 20 \text{ g}$};
    \node[above, font=\bfseries, darkgray] at (0, 0.5) {হার: 60 bullets/min};
\end{tikzpicture}
\end{center}
\vspace{1em}

% ------------------------------------------
% OPTIONS & ANSWER 33
% ------------------------------------------
\begin{itemize}[label={}, leftmargin=1cm, itemsep=0.5em]
    \item[\textbf{A)}] গড় বল $8.0\text{ N}$
    \item[\textbf{B)}] গড় বল $480\text{ N}$
    \item[\textbf{C)}] গড় বল $80.0\text{ N}$
    \item[\textbf{D)}] গড় বল $4.0\text{ N}$
\end{itemize}

\vspace{0.5em}
\begin{tcolorbox}[answerbox]
    \textbf{\color{success} উত্তর:} A) গড় বল $8.0\text{ N}$
\end{tcolorbox}
\vspace{1em}

% ------------------------------------------
% SOLUTION SECTION 33
% ------------------------------------------
\begin{tcolorbox}[solutionbox={সমাধান (Solution)}]
    \noindent
    \textbf{ধাপ ১: প্রতি সেকেন্ডে গুলির সংখ্যা ও ভর নির্ণয়} \\
    প্রতি মিনিটে ৬০টি গুলি ছোড়া হয়, তাই প্রতি সেকেন্ডে নির্গত গুলির সংখ্যা (হার):
    \[ n = \frac{60}{60\text{ s}} = 1\text{ s}^{-1} \]
    
    প্রতিটি গুলির ভর $m = 20\text{ g} = 0.02\text{ kg}$।
    
    \vspace{0.5em}
    \noindent\textbf{ধাপ ২: ভরবেগের পরিবর্তন} \\
    প্রতিটি গুলির ভরবেগ:
    \[ p = m v = 0.02 \times 400 = 8\text{ kg ms}^{-1} \]
    
    \vspace{0.5em}
    \noindent\textbf{ধাপ ৩: প্রযুক্ত বল নির্ণয়} \\
    নিউটনের দ্বিতীয় সূত্র অনুসারে, ভরবেগের পরিবর্তনের হারই হলো বন্দুকের ওপর প্রযুক্ত পশ্চাৎমুখী গড় ঘাত বল (Recoil force):
    \begin{align*}
        F_{\text{avg}} &= \frac{dp_{\text{total}}}{dt} = n \times (mv) \\
                       &= 1 \times 8 \\
                       &= \mathbf{8.0}\text{ N}
    \end{align*}
    
    বন্দুকটিকে যথাস্থানে স্থির রাখতে ব্যবহারকারীকে ঠিক সমমানের $8.0\text{ N}$ গড় বল সামনের দিকে প্রয়োগ করতে হবে।
\end{tcolorbox}

\newpage

% ------------------------------------------
% QUESTION SECTION 34
% ------------------------------------------
\begin{tcolorbox}[questionbox={প্রশ্ন (Question) 34}]
    \noindent
    অনুভূমিক তলের ওপর অবস্থিত $M$ ভরের একটি ভারী কাঠের ব্লকে $m$ ভরের একটি বুলেট অনুভূমিকভাবে $u$ বেগে আঘাত করে ব্লকের ভেতরে আটকে গেল। কাঠের ব্লকের সাথে মেঝের চল ঘর্ষণ গুণাঙ্ক $\mu$ হলে, বুলেটসহ ব্লকটি থেমে যাওয়ার পূর্বে মেঝের ওপর কত দূরত্ব অতিক্রম করবে?
    
    \vspace{0.8em}
    \begin{english}
    \noindent\textit{\color{darkgray}
    A bullet of mass $m$ moving horizontally with speed $u$ strikes and gets embedded in a wooden block of mass $M$ resting on a horizontal surface. If the coefficient of kinetic friction between the block and the surface is $\mu$, what distance will the combined block-bullet system slide before coming to rest?
    }
    \end{english}
\end{tcolorbox}

% ------------------------------------------
% HIGH-QUALITY TIKZ DIAGRAM 34
% ------------------------------------------
\vspace{1em}
\begin{center}
\begin{tikzpicture}[>=Stealth, scale=1.2]
    % Surface
    \fill[pattern=north east lines, pattern color=gray!50] (-3, -0.2) rectangle (5, 0);
    \draw[thick, darkgray] (-3, 0) -- (5, 0);
    \node[below, font=\footnotesize, text=gray!80!black] at (1, -0.2) {ঘর্ষণযুক্ত তল ($\mu$)};

    % Initial Block
    \draw[thick, fill=brown!80!black] (-0.5, 0) rectangle (0.5, 1);
    \node[white, font=\bfseries] at (0, 0.5) {$M$};
    
    % Bullet incoming
    \draw[thick, fill=black] (-2.2, 0.4) rectangle (-1.7, 0.6);
    \node[above, font=\bfseries] at (-1.95, 0.6) {$m$};
    \draw[->, line width=1.5pt, primary] (-1.7, 0.5) -- (-0.8, 0.5) node[midway, above] {$u$};

    % Final sliding block (ghosted)
    \draw[thick, fill=brown!80!black, opacity=0.4] (3.5, 0) rectangle (4.5, 1);
    \draw[thick, fill=black, opacity=0.4] (3.8, 0.4) rectangle (4.3, 0.6); % embedded bullet
    \node[above, font=\bfseries, text=black!70] at (4, 1.0) {$V_f=0$ (স্থির)};
    
    % Velocity vector right after collision
    \draw[->, line width=1.5pt, primary] (0.6, 0.8) -- (1.8, 0.8) node[midway, above] {$V$ (যৌথ বেগ)};

    % Distance s
    \draw[<->, dashed, thick] (0, 1.5) -- (4, 1.5) node[midway, above, font=\bfseries] {$s$};
    \draw[dashed] (0, 1) -- (0, 1.7);
    \draw[dashed] (4, 1) -- (4, 1.7);

    % Kinetic Friction
    \draw[->, line width=1.5pt, red] (2.5, 0) -- (1.5, 0) node[midway, below] {$f_k = \mu (M+m)g$};
\end{tikzpicture}
\end{center}
\vspace{1em}

% ------------------------------------------
% OPTIONS & ANSWER 34
% ------------------------------------------
\begin{itemize}[label={}, leftmargin=1cm, itemsep=0.5em]
    \item[\textbf{A)}] অতিক্রান্ত দূরত্ব $\frac{m^2 u^2}{2\mu g (M + m)^2}$
    \item[\textbf{B)}] অতিক্রান্ত দূরত্ব $\frac{m u^2}{2\mu g (M + m)}$
    \item[\textbf{C)}] অতিক্রান্ত দূরত্ব $\frac{m^2 u^2}{2\mu g M^2}$
    \item[\textbf{D)}] অতিক্রান্ত দূরত্ব $\frac{(M + m) u^2}{2\mu g m}$
\end{itemize}

\vspace{0.5em}
\begin{tcolorbox}[answerbox]
    \textbf{\color{success} উত্তর:} A) অতিক্রান্ত দূরত্ব $\frac{m^2 u^2}{2\mu g (M + m)^2}$
\end{tcolorbox}
\vspace{1em}

% ------------------------------------------
% SOLUTION SECTION 34
% ------------------------------------------
\begin{tcolorbox}[solutionbox={সমাধান (Solution)}]
    \noindent
    \textbf{ধাপ ১: ভরবেগের সংরক্ষণ (Collision phase)} \\
    বুলেটটি ব্লকে আটকে যাওয়ার প্রক্রিয়াটি অত্যন্ত দ্রুত সম্পন্ন হয়, তাই বাহ্যিক ঘর্ষণ বলের ঘাত নগণ্য থাকে এবং রৈখিক ভরবেগ সংরক্ষিত থাকে। ধরি সংঘর্ষের ঠিক পর যৌথ ব্যবস্থার সাধারণ বেগ $V$।
    \begin{align*}
        m u + M (0) &= (M + m) V \\
        V &= \frac{m u}{M + m}
    \end{align*}
    
    \vspace{0.5em}
    \noindent\textbf{ধাপ ২: গতিশক্তি নির্ণয়} \\
    সংঘর্ষের ঠিক পর ব্লক ও বুলেটের মিলিত গতিশক্তি:
    \begin{align*}
        K &= \frac{1}{2} (M + m) V^2 \\
          &= \frac{1}{2} (M + m) \left(\frac{m u}{M + m}\right)^2 \\
          &= \frac{m^2 u^2}{2 (M + m)}
    \end{align*}
    
    \vspace{0.5em}
    \noindent\textbf{ধাপ ৩: কাজ-শক্তি উপপাদ্য (Sliding phase)} \\
    মেঝের চল ঘর্ষণ বল, $f_k = \mu N = \mu (M + m) g$। 
    কাজ-শক্তি উপপাদ্য অনুসারে, থেমে যাওয়ার পূর্বে ঘর্ষণ বলের বিরুদ্ধে কৃতকাজ বস্তুর গতিশক্তির সমান হবে:
    \begin{align*}
        W_{\text{friction}} &= \Delta K \\
        f_k \times s &= K \\
        \mu (M + m) g \times s &= \frac{m^2 u^2}{2 (M + m)}
    \end{align*}
    
    অতএব, অতিক্রান্ত দূরত্ব $s$:
    \begin{align*}
        s &= \mathbf{\frac{m^2 u^2}{2\mu g (M + m)^2}}
    \end{align*}
\end{tcolorbox}

% ------------------------------------------
% QUESTION SECTION 35
% ------------------------------------------
\begin{tcolorbox}[questionbox={প্রশ্ন (Question) 35}]
    \noindent
    $M = 3\text{ kg}$ ভরের একটি ব্লক একটি মসৃণ অনুভূমিক মেঝের ওপর রাখা আছে। একটি ভরহীন অপ্রসারণশীল সুতার এক প্রান্ত দৃঢ় দেওয়ালের সাথে বাঁধা এবং সুতাটি ব্লকের সাথে সংযুক্ত একটি ঘর্ষণহীন হালকা চলনক্ষম কপিকলের (movable pulley) ওপর দিয়ে গিয়ে অপর প্রান্তে $m = 2\text{ kg}$ ভরের একটি ঝুলন্ত ব্লকের সাথে যুক্ত। পুরো সংস্থাকে স্থির অবস্থা থেকে মুক্ত করে দিলে $M$ ভরের অনুভূমিক ব্লকটির ত্বরণ কত হবে? ($g = 9.8\text{ ms}^{-2}$)
    
    \vspace{0.8em}
    \begin{english}
    \noindent\textit{\color{darkgray}
    A block of mass $M = 3\text{ kg}$ rests on a smooth horizontal floor. A light inextensible cord has one end fixed to a rigid wall, passes around a frictionless massless movable pulley connected to the block $M$, and then hangs vertically supporting a suspended block of mass $m = 2\text{ kg}$. If the system is released from rest, what is the acceleration of the horizontal block $M$? ($g = 9.8\text{ ms}^{-2}$)
    }
    \end{english}
\end{tcolorbox}

% ------------------------------------------
% HIGH-QUALITY TIKZ DIAGRAM 35
% ------------------------------------------
\vspace{1em}
\begin{center}
\begin{tikzpicture}[>=Stealth, scale=1.3]
    % Floor
    \draw[thick, gray] (-4, 0) -- (2.5, 0);
    \fill[pattern=north east lines, pattern color=gray!50] (-4, -0.2) rectangle (2.5, 0);
    
    % Post (Rigid Wall) on the table
    \fill[gray!60] (1.5, 0) rectangle (2.0, 1.2);
    \draw[thick, darkgray] (1.5, 0) rectangle (2.0, 1.2);
    \fill[pattern=north east lines, pattern color=gray!80] (1.5, 0) rectangle (2.0, 1.2);
    
    % Fixed Pulley at the edge
    \draw[thick, fill=gray!30] (2.5, 0.2) circle (0.2);
    \fill[black] (2.5, 0.2) circle (0.05);
    \draw[thick] (2.2, 0) -- (2.5, 0.2); % Mount
    
    % Block M
    \draw[thick, fill=orange!80!black, rounded corners=2pt] (-2.2, 0) rectangle (-0.8, 0.8);
    \node[white, font=\bfseries] at (-1.5, 0.4) {$M$};
    
    % Movable Pulley
    \draw[thick, fill=gray!30] (-0.3, 0.4) circle (0.2);
    \fill[black] (-0.3, 0.4) circle (0.05);
    \draw[line width=2pt, darkgray] (-0.8, 0.4) -- (-0.3, 0.4); % Connect to block
    \node[above, font=\footnotesize] at (-0.3, 0.6) {Movable};
    
    % Cord
    \draw[line width=1.2pt, red] (1.5, 0.6) -- (-0.3, 0.6); % Post to movable top
    \draw[line width=1.2pt, red] (-0.3, 0.2) -- (2.5, 0.2); % Movable bot to fixed top (approx)
    \draw[line width=1.2pt, red] (2.7, 0.2) -- (2.7, -1.5); % Fixed down to mass m
    
    \fill[black] (1.5, 0.6) circle (0.04); % Anchor point
    
    % Hanging Mass m
    \draw[thick, fill=cyan!70!black, rounded corners=1pt] (2.4, -2.1) rectangle (3.0, -1.5);
    \node[white, font=\bfseries] at (2.7, -1.8) {$m$};
    
    % Acceleration Vectors
    \draw[->, line width=1.5pt, primary] (-1.5, 1.0) -- (-0.5, 1.0) node[midway, above] {$a_M = a$};
    \draw[->, line width=1.5pt, primary] (3.2, -1.2) -- (3.2, -2.2) node[midway, right] {$a_m = 2a$};
    
    % Tensions
    \node[above, text=red, font=\footnotesize] at (0.6, 0.6) {$T$};
    \node[below, text=red, font=\footnotesize] at (0.6, 0.2) {$T$};
    \draw[->, thick, red!50!black] (-0.3, 0.6) -- (0.2, 0.6);
    \draw[->, thick, red!50!black] (-0.3, 0.2) -- (0.2, 0.2);
\end{tikzpicture}
\end{center}
\vspace{1em}

% ------------------------------------------
% OPTIONS & ANSWER 35
% ------------------------------------------
\begin{itemize}[label={}, leftmargin=1cm, itemsep=0.5em]
    \item[\textbf{A)}] ব্লকের ত্বরণ $3.56\text{ ms}^{-2}$
    \item[\textbf{B)}] ব্লকের ত্বরণ $1.78\text{ ms}^{-2}$
    \item[\textbf{C)}] ব্লকের ত্বরণ $4.90\text{ ms}^{-2}$
    \item[\textbf{D)}] ব্লকের ত্বরণ $2.45\text{ ms}^{-2}$
\end{itemize}

\vspace{0.5em}
\begin{tcolorbox}[answerbox]
    \textbf{\color{success} উত্তর:} A) ব্লকের ত্বরণ $3.56\text{ ms}^{-2}$
\end{tcolorbox}
\vspace{1em}

% ------------------------------------------
% SOLUTION SECTION 35
% ------------------------------------------
\begin{tcolorbox}[solutionbox={সমাধান (Solution)}]
    \noindent
    \textbf{ধাপ ১: বাধ্যবাধকতা সম্পর্ক (Constraint Relation)} \\
    ধরি, চলনক্ষম কপিকলসহ অনুভূমিক ব্লক $M$-এর ডানদিকে সরণ $x$। সুতাটি চলনক্ষম কপিকলের ওপর দিয়ে ঘুরে আসায়, ব্লকটি $x$ পরিমাণ সরলে সুতার উপরের ও নিচের উভয় অংশ থেকে $x$ করে মোট $2x$ দৈর্ঘ্য মুক্ত হবে। 
    ফলে ঝুলন্ত ভর $m$-এর উল্লম্ব নিম্নমুখী সরণ হবে $y = 2x$। ত্বরণের ক্ষেত্রেও একই সম্পর্ক প্রযোজ্য:
    \[ a_m = 2 a_M \]
    ধরি, $a_M = a$, তাহলে $a_m = 2a$।
    
    \vspace{0.5em}
    \noindent\textbf{ধাপ ২: গতির সমীকরণ (Equations of Motion)} \\
    সুতার টান $T$ হলে, চলনক্ষম কপিকলের মাধ্যমে ব্লক $M$-এর ওপর অনুভূমিক ডানদিকে মোট টান বল হবে $2T$ (কারণ সুতার দুটি অংশই কপিকলকে ডানদিকে টানছে)। 
    ব্লক $M$-এর অনুভূমিক গতির সমীকরণ:
    \begin{align*}
        2T &= M a \\
        T &= \frac{M a}{2}
    \end{align*}
    
    ঝুলন্ত ভর $m$-এর উল্লম্ব গতির সমীকরণ:
    \[ mg - T = m (2a) \]
    
    \vspace{0.5em}
    \noindent\textbf{ধাপ ৩: ত্বরণ নির্ণয়} \\
    দ্বিতীয় সমীকরণে $T$-এর মান বসিয়ে পাই:
    \begin{align*}
        mg - \frac{M a}{2} &= 2ma \\
        mg &= \left(2m + \frac{M}{2}\right) a \\
        mg &= \left(\frac{4m + M}{2}\right) a \\
        a &= \frac{2mg}{4m + M}
    \end{align*}
    
    প্রদত্ত মানসমূহ ($M = 3\text{ kg}, m = 2\text{ kg}, g = 9.8\text{ ms}^{-2}$) বসিয়ে:
    \begin{align*}
        a &= \frac{2 \times 2 \times 9.8}{(4 \times 2) + 3} \\
          &= \frac{39.2}{8 + 3} \\
          &= \frac{39.2}{11} \approx 3.564\text{ ms}^{-2} \approx \mathbf{3.56}\text{ ms}^{-2}
    \end{align*}
\end{tcolorbox}

\newpage

% ------------------------------------------
% QUESTION SECTION 36
% ------------------------------------------
\begin{tcolorbox}[questionbox={প্রশ্ন (Question) 36}]
    \noindent
    উল্লম্ব সমতলে স্থাপিত $R = 2.5\text{ m}$ ব্যাসার্ধের একটি মসৃণ বৃত্তাকার লুপের ভেতরে তলদেশ থেকে একটি ছোট ব্লক কত ন্যূনতম বেগে প্রবেশ করলে এটি লুপের সম্পূর্ণ উলম্ব পথ অতিক্রম করতে পারবে এবং লুপের অনুভূমিক ব্যাসের প্রান্তবিন্দুতে (যেখানে স্পর্শক খাড়া উলম্ব) ব্লকটির ওপর ট্র্যাক কর্তৃক প্রযুক্ত অভিলম্ব প্রতিক্রিয়া বলের মান কত হবে? (ব্লকের ভর $m = 1\text{ kg}, g = 9.8\text{ ms}^{-2}$)
    
    \vspace{0.8em}
    \begin{english}
    \noindent\textit{\color{darkgray}
    A small block of mass $m = 1\text{ kg}$ enters the bottom of a smooth vertical circular loop of radius $R = 2.5\text{ m}$. What minimum speed must it possess at the lowest point to complete the full vertical loop, and what will be the normal reaction force exerted by the track on the block when it reaches the horizontal diameter level (where the tangent is vertical)? ($g = 9.8\text{ ms}^{-2}$)
    }
    \end{english}
\end{tcolorbox}

% ------------------------------------------
% HIGH-QUALITY TIKZ DIAGRAM 36
% ------------------------------------------
\vspace{1em}
\begin{center}
\begin{tikzpicture}[>=Stealth, scale=1.1]
    \def\myR{2.5}
    
    % Loop
    \draw[line width=2pt, cyan!30] (0, \myR) circle (\myR);
    \draw[dashed, gray] (0, \myR) circle (\myR-0.2);
    
    % Axes
    \draw[dashed, darkgray] (-\myR-0.5, \myR) -- (\myR+0.5, \myR) node[right] {অনুভূমিক ব্যাস};
    \draw[dashed, darkgray] (0, -0.5) -- (0, 2*\myR+0.5);
    \fill[black] (0, \myR) circle (0.05) node[below left] {$O$};
    
    % Block at bottom
    \coordinate (B_bot) at (0, 0.15);
    \draw[thick, fill=orange!80!black] (B_bot) circle (0.15);
    \draw[->, line width=1.5pt, primary] (B_bot) ++(0.2, 0) -- ++(1.5, 0) node[right] {$u_{\min}$};
    
    % Block at horizontal diameter (right)
    \coordinate (B_mid) at (\myR-0.15, \myR);
    \draw[thick, fill=orange!80!black] (B_mid) circle (0.15);
    \draw[->, line width=1.5pt, primary] (B_mid) ++(0, 0.2) -- ++(0, 1.5) node[above] {$v$};
    
    % Forces at horizontal diameter
    \draw[->, line width=1.5pt, green!60!black] (B_mid) ++(-0.2, 0) -- ++(-1.5, 0) node[above] {$N$};
    \draw[->, line width=1.5pt, red] (B_mid) ++(0, -0.2) -- ++(0, -1.2) node[right] {$mg$};
    
    % Radius indicator
    \draw[->, dashed] (0, \myR) -- ({-\myR*cos(45)}, {\myR+\myR*sin(45)}) node[midway, below left] {$R$};
\end{tikzpicture}
\end{center}
\vspace{1em}

% ------------------------------------------
% OPTIONS & ANSWER 36
% ------------------------------------------
\begin{itemize}[label={}, leftmargin=1cm, itemsep=0.5em]
    \item[\textbf{A)}] সর্বনিম্ন বেগ $11.07\text{ ms}^{-1}$ এবং প্রতিক্রিয়া বল $29.4\text{ N}$
    \item[\textbf{B)}] সর্বনিম্ন বেগ $11.07\text{ ms}^{-1}$ এবং প্রতিক্রিয়া বল $49.0\text{ N}$
    \item[\textbf{C)}] সর্বনিম্ন বেগ $7.00\text{ ms}^{-1}$ এবং প্রতিক্রিয়া বল $19.6\text{ N}$
    \item[\textbf{D)}] সর্বনিম্ন বেগ $15.65\text{ ms}^{-1}$ এবং প্রতিক্রিয়া বল $39.2\text{ N}$
\end{itemize}

\vspace{0.5em}
\begin{tcolorbox}[answerbox]
    \textbf{\color{success} উত্তর:} A) সর্বনিম্ন বেগ $11.07\text{ ms}^{-1}$ এবং প্রতিক্রিয়া বল $29.4\text{ N}$
\end{tcolorbox}
\vspace{1em}

% ------------------------------------------
% SOLUTION SECTION 36
% ------------------------------------------
\begin{tcolorbox}[solutionbox={সমাধান (Solution)}]
    \noindent
    \textbf{ধাপ ১: ক্রান্তি বেগ নির্ণয়} \\
    উল্লম্ব বৃত্তাকার লুপ সম্পূর্ণ করার জন্য সর্বনিম্ন বিন্দুতে ন্যূনতম ক্রান্তি বেগের তাত্ত্বিক সমীকরণ:
    \begin{align*}
        u_{\min} &= \sqrt{5gR} \\
        &= \sqrt{5 \times 9.8 \times 2.5} \\
        &= \sqrt{122.5} \\
        &\approx 11.068\text{ ms}^{-1} \approx \mathbf{11.07}\text{ ms}^{-1}
    \end{align*}
    
    \vspace{0.5em}
    \noindent\textbf{ধাপ ২: অনুভূমিক ব্যাসের স্তরে বেগ নির্ণয়} \\
    অনুভূমিক ব্যাসের স্তরে ব্লকটি সর্বনিম্ন বিন্দু হতে $h = R$ উচ্চতায় ওঠে। শক্তির সংরক্ষণ সূত্র প্রয়োগ করে ঐ বিন্দুতে বেগ $v$ নির্ণয় করা যায়:
    \begin{align*}
        \frac{1}{2} m u^2 &= \frac{1}{2} m v^2 + m g h \\
        v^2 &= u^2 - 2gR
    \end{align*}
    যেহেতু $u^2 = 5gR$:
    \begin{align*}
        v^2 &= 5gR - 2gR = 3gR
    \end{align*}
    
    \vspace{0.5em}
    \noindent\textbf{ধাপ ৩: অভিলম্ব প্রতিক্রিয়া বল নির্ণয়} \\
    অনুভূমিক ব্যাসের স্তরে, অভিলম্ব প্রতিক্রিয়া বল $N$ সম্পূর্ণরূপে অনুভূমিক বরাবর কাজ করে এবং বৃত্তাকার গতির জন্য প্রয়োজনীয় কেন্দ্রমুখী বল প্রদান করে (ওজন $mg$ এক্ষেত্রে খাড়া নিম্নমুখী হওয়ায় এর কোনো কেন্দ্রমুখী উপাংশ থাকে না)।
    \begin{align*}
        N &= \frac{m v^2}{R} \\
          &= \frac{m (3gR)}{R} \\
          &= 3mg
    \end{align*}
    
    মান বসিয়ে পাই:
    \begin{align*}
        N &= 3 \times 1 \times 9.8 \\
          &= \mathbf{29.4}\text{ N}
    \end{align*}
\end{tcolorbox}

\newpage

% ------------------------------------------
% QUESTION SECTION 37 (MORE COMPLEX)
% ------------------------------------------
\begin{tcolorbox}[questionbox={প্রশ্ন (Question) 37}]
\noindent
$M_0 = 400\text{ kg}$ ভরের একটি খোলা ট্রলি অনুভূমিক অমসৃণ ট্র্যাকে $v_0 = 20\text{ ms}^{-1}$ আদি বেগে চলছে। ট্রলি ও ট্র্যাকের মধ্যকার গতীয় ঘর্ষণ গুণাঙ্ক $\mu_k = 0.05$। হঠাৎ উল্লম্বভাবে খাড়া নিচের দিকে প্রতি সেকেন্ডে ভূমি সাপেক্ষে $r = 2\text{ kgs}^{-1}$ হারে বৃষ্টি পড়তে শুরু করল এবং ট্রলিতে পানি জমা হতে লাগল। একই সাথে ট্রলির গতির অভিমুখে $F = 250\text{ N}$ মানের একটি ধ্রুব অনুভূমিক বাহ্যিক বল প্রয়োগ করা হলো। বৃষ্টি শুরুর $t = 50\text{ s}$ পর ট্রলির বেগ কত হবে? ($g = 9.8\text{ ms}^{-2}$)
\vspace{0.8em}
\begin{english}
\noindent\textit{\color{darkgray}
An open trolley of initial mass $M_0 = 400\text{ kg}$ is traveling horizontally on a rough track with a kinetic friction coefficient $\mu_k = 0.05$ at an initial velocity $v_0 = 20\text{ ms}^{-1}$. Suddenly, rain starts falling vertically downwards into the trolley at a constant rate of $r = 2\text{ kgs}^{-1}$ and accumulates. Simultaneously, a constant horizontal external pulling force $F = 250\text{ N}$ is applied in the direction of motion. What will be the velocity of the trolley at $t = 50\text{ s}$ after the rain began? ($g = 9.8\text{ ms}^{-2}$)
}
\end{english}
\end{tcolorbox}

% ------------------------------------------
% HIGH-QUALITY TIKZ DIAGRAM 37
% ------------------------------------------
\vspace{1em}
\begin{center}
\begin{tikzpicture}[>=Stealth, scale=1.2]
% Track with friction indication
\draw[thick, gray] (-3, 0) -- (4, 0);
\foreach \x in {-2.8, -2.4, ..., 3.8} {
    \draw[gray] (\x, 0) -- (\x-0.1, -0.15);
}
\node[below, font=\small, text=gray!80!black] at (0.5, -0.3) {অমসৃণ অনুভূমিক তল ($\mu_k = 0.05$)};

% Trolley
\fill[cyan!15] (-1.5, 0.4) rectangle (1.5, 1.5);
\draw[line width=2pt, darkgray] (-1.5, 1.5) -- (-1.5, 0.4) -- (1.5, 0.4) -- (1.5, 1.5);

% Wheels
\draw[thick, fill=gray!40] (-1, 0.2) circle (0.2);
\draw[thick, fill=gray!40] (1, 0.2) circle (0.2);
\fill[black] (-1, 0.2) circle (0.05);
\fill[black] (1, 0.2) circle (0.05);

% Rain accumulation
\fill[blue!40, opacity=0.7] (-1.5, 0.4) rectangle (1.5, 0.7);
\draw[blue] (-1.5, 0.7) .. controls (-0.5, 0.75) and (0.5, 0.65) .. (1.5, 0.7);

% Rain falling
\foreach \x in {-1.2, -0.6, 0, 0.6, 1.2} {
    \draw[->, dashed, thick, blue] (\x, 2.5) -- (\x, 1.6);
}
\node[above, text=blue, font=\bfseries] at (0, 2.5) {বৃষ্টি $r = 2 \text{ kgs}^{-1}$};

% External Force
\draw[->, line width=1.5pt, red] (1.5, 0.95) -- (3.0, 0.95) node[right, font=\bfseries] {$F = 250\text{ N}$};

% Friction Force
\draw[->, line width=1.5pt, orange!80!black] (-1.5, 0.2) -- (-2.5, 0.2) node[left, font=\bfseries] {$f_k = \mu_k N$};

% Velocity
\draw[->, line width=1.5pt, primary] (0, 1.8) -- (1.5, 1.8) node[right, font=\bfseries] {$v(t)$};

% Mass Label
\node[font=\bfseries] at (0, 1.1) {$M(t) = M_0 + rt$};
\end{tikzpicture}
\end{center}
\vspace{1em}

% ------------------------------------------
% OPTIONS & ANSWER 37
% ------------------------------------------
\begin{itemize}[label={}, leftmargin=1cm, itemsep=0.5em]
\item[\textbf{A)}] ট্রলির বেগ $18.95\text{ ms}^{-1}$
\item[\textbf{B)}] ট্রলির বেগ $15.40\text{ ms}^{-1}$
\item[\textbf{C)}] ট্রলির বেগ $21.20\text{ ms}^{-1}$
\item[\textbf{D)}] ট্রলির বেগ $12.50\text{ ms}^{-1}$
\end{itemize}
\vspace{0.5em}
\begin{tcolorbox}[answerbox]
\textbf{\color{success} উত্তর:} A) ট্রলির বেগ $18.95\text{ ms}^{-1}$
\end{tcolorbox}
\vspace{1em}

% ------------------------------------------
% SOLUTION SECTION 37
% ------------------------------------------
\begin{tcolorbox}[solutionbox={সমাধান (Solution)}]
\noindent
\textbf{ধাপ ১: গতির অন্তরক সমীকরণ গঠন} \\
বৃষ্টির ফোঁটার অনুভূমিক বেগ শূন্য হওয়ায় এবং ট্র্যাকে ঘর্ষণ বিদ্যমান থাকায় গতিশীল ব্যবস্থার ওপর প্রযুক্ত নিট অনুভূমিক বল হবে বাহ্যিক বল ও ঘর্ষণ বলের লব্ধি। পরিবর্তনশীল ভর $M(t)$ এবং গতীয় ঘর্ষণ বল $f_k = \mu_k N = \mu_k M(t)g$-কে বিবেচনা করে নিউটনের দ্বিতীয় সূত্রানুসারে পাই:
\begin{align*}
\frac{d(Mv)}{dt} &= F - f_k \\
\frac{d(Mv)}{dt} &= F - \mu_k M(t) g
\end{align*}
এখানে যেকোনো সময়ে ট্রলির মোট ভর $M(t) = M_0 + rt$, ফলে $\frac{dM}{dt} = r$ হওয়ায় সময়ের সাপেক্ষে অন্তরক রূপটি ভরের সাপেক্ষে রূপান্তর করা যায় ($dt = \frac{dM}{r}$):
\begin{align*}
    r \frac{d(Mv)}{dM} &= F - \mu_k M g \\
    d(Mv) &= \left(\frac{F}{r} - \frac{\mu_k g}{r} M\right) dM
\end{align*}

\vspace{0.5em}
\noindent\textbf{ধাপ ২: সমাকলন (Integration)} \\
প্রাথমিক ভর $M_0$-তে বেগ $v_0$ এবং $t$ সময়ে ভর $M$-এ বেগ স্পষ্ট করে সমাকলন করে পাই:
\begin{align*}
    \int_{M_0 v_0}^{M v} d(Mv) &= \int_{M_0}^{M} \left(\frac{F}{r} - \frac{\mu_k g}{r} M\right) dM \\
    M v - M_0 v_0 &= \frac{F}{r}(M - M_0) - \frac{\mu_k g}{2r} (M^2 - M_0^2)
\end{align*}

\vspace{0.5em}
\noindent\textbf{ধাপ ৩: মান বসিয়ে বেগ নির্ণয়} \\
প্রদত্ত মানসমূহ: $M_0 = 400\text{ kg}$, $v_0 = 20\text{ ms}^{-1}$, $r = 2\text{ kgs}^{-1}$, $F = 250\text{ N}$, $\mu_k = 0.05$, $g = 9.8\text{ ms}^{-2}$ এবং $t = 50\text{ s}$।

$t = 50\text{ s}$ পর ট্রলির মোট ভর:
\[M = M_0 + rt = 400 + (2 \times 50) = 500\text{ kg}\]

সমীকরণে মানগুলো বসিয়ে পাই:
\begin{align*}
    (500 \times v) - (400 \times 20) &= \frac{250}{2}(500 - 400) - \frac{0.05 \times 9.8}{2 \times 2}(500^2 - 400^2) \\
    500 v - 8000 &= 125 \times 100 - \frac{0.49}{4}(250000 - 160000) \\
    500 v - 8000 &= 12500 - 0.1225 \times 90000 \\
    500 v - 8000 &= 12500 - 11025 \\
    500 v - 8000 &= 1475 \\
    500 v &= 9475 \\
    v &= \frac{9475}{500} \\
    v &\approx \mathbf{18.95}\text{ ms}^{-1}
\end{align*}
\end{tcolorbox}
% ------------------------------------------
% QUESTION SECTION 38 (SIMPLIFIED)
% ------------------------------------------
\begin{tcolorbox}[questionbox={প্রশ্ন (Question) 38}]
    \noindent
    $M = 2\text{ kg}$ ভরের একটি নিরেট সিলিন্ডার $v = 4\text{ ms}^{-1}$ বেগে কোনো অনুভূমিক তলে না পিছলিয়ে বিশুদ্ধ গড়িয়ে চলছে। সিলিন্ডারটির মোট গতিশক্তি কত হবে?
    
    \vspace{0.8em}
    \begin{english}
    \noindent\textit{\color{darkgray}
    A solid cylinder of mass $M = 2\text{ kg}$ rolls purely without slipping on a horizontal surface at a linear speed of $v = 4\text{ ms}^{-1}$. What is its total kinetic energy?
    }
    \end{english}
\end{tcolorbox}

% ------------------------------------------
% HIGH-QUALITY TIKZ DIAGRAM 38
% ------------------------------------------
\vspace{1em}
\begin{center}
\begin{tikzpicture}[>=Stealth, scale=1.3]
    % Floor
    \draw[thick, gray] (-2.5, 0) -- (2.5, 0);
    \fill[pattern=north east lines, pattern color=gray!50] (-2.5, -0.2) rectangle (2.5, 0);
    \node[below, text=gray!80!black] at (0, -0.2) {অনুভূমিক তল (Horizontal Floor)};
    
    % Cylinder
    \coordinate (C) at (0, 1);
    \draw[thick, fill=cyan!30] (C) circle (1);
    \fill[black] (C) circle (0.05);
    
    % Velocity and Angular velocity vectors
    \draw[->, line width=1.5pt, red!80!black] (C) ++(0.2, 0) -- ++(1.2, 0) node[above] {$v$};
    \draw[->, line width=1.2pt, blue] (C) ++(60:1.3) arc (60:10:1.3) node[midway, right] {$\omega$};
    
    % Radius and Mass labels
    \draw[dashed, darkgray] (C) -- (0,0) node[midway, left] {$R$};
    \node[above] at (0, 2) {$M = 2\text{ kg}$};
\end{tikzpicture}
\end{center}
\vspace{1em}

% ------------------------------------------
% OPTIONS & ANSWER 38
% ------------------------------------------
\begin{itemize}[label={}, leftmargin=1cm, itemsep=0.5em]
    \item[\textbf{A)}] মোট গতিশক্তি $24\text{ J}$
    \item[\textbf{B)}] মোট গতিশক্তি $16\text{ J}$
    \item[\textbf{C)}] মোট গতিশক্তি $32\text{ J}$
    \item[\textbf{D)}] মোট গতিশক্তি $48\text{ J}$
\end{itemize}

\vspace{0.5em}
\begin{tcolorbox}[answerbox]
    \textbf{\color{success} উত্তর:} A) মোট গতিশক্তি $24\text{ J}$
\end{tcolorbox}
\vspace{1em}

% ------------------------------------------
% SOLUTION SECTION 38
% ------------------------------------------
\begin{tcolorbox}[solutionbox={সমাধান (Solution)}]
    \noindent
    নিয়েট সিলিন্ডারের গড়িয়ে চলার ক্ষেত্রে মোট গতিশক্তি হলো রৈখিক গতিশক্তি ও ঘূর্ণন গতিশক্তির সমষ্টি:
    \begin{align*}
        E_k &= E_{\text{linear}} + E_{\text{rotational}} \\
            &= \frac{1}{2} M v^2 + \frac{1}{2} I \omega^2
    \end{align*}
    
    নিরেট সিলিন্ডারের নিজ অক্ষের সাপেক্ষে জড়তার ভ্রামক $I = \frac{1}{2} M R^2$ এবং বিশুদ্ধ ঘূর্ণনের শর্তানুসারে $\omega = \frac{v}{R}$। মানগুলো বসিয়ে পাই:
    \begin{align*}
        E_k &= \frac{1}{2} M v^2 + \frac{1}{2} \left(\frac{1}{2} M R^2\right) \left(\frac{v}{R}\right)^2 \\
            &= \frac{1}{2} M v^2 + \frac{1}{4} M v^2 \\
            &= \frac{3}{4} M v^2
    \end{align*}
    
    প্রদত্ত মানসমূহ ($M = 2\text{ kg}, v = 4\text{ ms}^{-1}$) বসিয়ে পাই:
    \begin{align*}
        E_k &= \frac{3}{4} \times 2 \times (4)^2 \\
            &= \frac{3}{4} \times 2 \times 16 \\
            &= 3 \times 8 \\
            &= \mathbf{24}\text{ J}
    \end{align*}
\end{tcolorbox}
% ------------------------------------------
% QUESTION SECTION 39
% ------------------------------------------
\begin{tcolorbox}[questionbox={প্রশ্ন (Question) 39}]
    \noindent
    একটি অনুভূমিক মসৃণ তলে $M = 4\text{ kg}$ ভরের একটি তক্তা (plank) রাখা আছে এবং তার ওপর $m = 2\text{ kg}$ ভরের একটি ব্লক রাখা। ব্লক ও তক্তার মধ্যকার ঘর্ষণ গুণাঙ্ক $\mu_s = \mu_k = 0.3$। উপরের $m$ ব্লকের ওপর একটি সময়-নির্ভর পরিবর্তনশীল অনুভূমিক বল $F(t) = 2t\text{ N}$ প্রয়োগ করা হলো। বল প্রয়োগের মুহূর্ত $t = 0$ হতে কত সময় পর তক্তা ও ব্লকের মধ্যে আপেক্ষিক পিছলানো শুরু হবে এবং $t = 8\text{ s}$ সময়ে তক্তার ত্বরণ কত হবে? ($g = 9.8\text{ ms}^{-2}$)
    
    \vspace{0.8em}
    \begin{english}
    \noindent\textit{\color{darkgray}
    A plank of mass $M = 4\text{ kg}$ rests on a frictionless horizontal floor, and a block of mass $m = 2\text{ kg}$ rests on top of the plank. The coefficient of static and kinetic friction between the block and the plank is $\mu_s = \mu_k = 0.3$. A time-dependent horizontal force $F(t) = 2t\text{ N}$ is applied to the upper block $m$. At what time $t$ will relative slipping between the block and the plank commence, and what will be the acceleration of the plank at $t = 8\text{ s}$? ($g = 9.8\text{ ms}^{-2}$)
    }
    \end{english}
\end{tcolorbox}

% ------------------------------------------
% HIGH-QUALITY TIKZ DIAGRAM 39
% ------------------------------------------
\vspace{1em}
\begin{center}
\begin{tikzpicture}[>=Stealth, scale=1.3]
    % Floor
    \draw[thick, gray] (-3, 0) -- (4, 0);
    \fill[pattern=north east lines, pattern color=gray!50] (-3, -0.2) rectangle (4, 0);
    \node[below, font=\footnotesize, text=gray!80!black] at (0.5, -0.2) {ঘর্ষণহীন তল (Smooth Floor)};
    
    % Plank M
    \draw[thick, fill=brown!70!black, rounded corners=2pt] (-2, 0) rectangle (3, 0.6);
    \node[white, font=\bfseries] at (2.2, 0.3) {$M = 4\text{ kg}$};
    
    % Block m
    \draw[thick, fill=orange!80!black, rounded corners=2pt] (-1, 0.6) rectangle (0.5, 1.4);
    \node[white, font=\bfseries] at (-0.25, 1.0) {$m = 2\text{ kg}$};
    
    % Force F(t)
    \draw[->, line width=1.8pt, primary] (0.5, 1.0) -- (2.2, 1.0) node[right, font=\bfseries] {$F(t) = 2t$};
    
    % Friction Force between M and m
    \draw[->, line width=1.5pt, red] (-0.25, 0.6) -- (-1.5, 0.6) node[midway, above] {$f_s$ / $f_k$};
    \node[below, font=\footnotesize, text=red] at (-0.5, 1.8) {$\mu_s = \mu_k = 0.3$};
    
    % Acceleration of Plank
    \draw[->, line width=1.5pt, green!60!black] (3.1, 0.3) -- (4.2, 0.3) node[right] {$a_{\text{plank}}$};
\end{tikzpicture}
\end{center}
\vspace{1em}

% ------------------------------------------
% OPTIONS & ANSWER 39
% ------------------------------------------
\begin{itemize}[label={}, leftmargin=1cm, itemsep=0.5em]
    \item[\textbf{A)}] পিছলানোর সময় $4.41\text{ s}$ এবং তক্তার ত্বরণ $1.47\text{ ms}^{-2}$
    \item[\textbf{B)}] পিছলানোর সময় $2.94\text{ s}$ এবং তক্তার ত্বরণ $2.25\text{ ms}^{-2}$
    \item[\textbf{C)}] পিছলানোর সময় $4.41\text{ s}$ এবং তক্তার ত্বরণ $2.94\text{ ms}^{-2}$
    \item[\textbf{D)}] পিছলানোর সময় $5.88\text{ s}$ এবং তক্তার ত্বরণ $0.98\text{ ms}^{-2}$
\end{itemize}

\vspace{0.5em}
\begin{tcolorbox}[answerbox]
    \textbf{\color{success} উত্তর:} A) পিছলানোর সময় $4.41\text{ s}$ এবং তক্তার ত্বরণ $1.47\text{ ms}^{-2}$
\end{tcolorbox}
\vspace{1em}

% ------------------------------------------
% SOLUTION SECTION 39
% ------------------------------------------
\begin{tcolorbox}[solutionbox={সমাধান (Solution)}]
    \noindent
    \textbf{ধাপ ১: সর্বোচ্চ স্থিতি ঘর্ষণ বল ও তক্তার সর্বোচ্চ ত্বরণ} \\
    ব্লক ও তক্তার মধ্যকার সর্বোচ্চ স্থিতি ঘর্ষণ বল:
    \begin{align*}
        f_{\max} &= \mu_s m g \\
        &= 0.3 \times 2 \times 9.8 = 5.88\text{ N}
    \end{align*}
    
    যেহেতু মেঝে ঘর্ষণহীন, তক্তাটিকে শুধুমাত্র এই ঘর্ষণ বলই চালিত করতে পারে। ফলে আপেক্ষিক পিছলানোর পূর্ব পর্যন্ত তক্তার (Plank) সর্বোচ্চ সম্ভাব্য ত্বরণ:
    \[ a_{\max} = \frac{f_{\max}}{M} = \frac{5.88}{4} = 1.47\text{ ms}^{-2} \]
    
    \vspace{0.5em}
    \noindent\textbf{ধাপ ২: পিছলানো শুরুর সময় নির্ণয়} \\
    ততক্ষণ পর্যন্ত উভয় বস্তু একসাথে চলবে যতক্ষণ বাহ্যিক বল $F(t)$ সম্পূর্ণ সিস্টেমকে $a_{\max}$ ত্বরণ প্রদানে সক্ষম:
    \begin{align*}
        F(t) &= (M + m) a_{\max} \\
        &= (4 + 2) \times 1.47 \\
        &= 6 \times 1.47 = 8.82\text{ N}
    \end{align*}
    
    যেহেতু $F(t) = 2t$:
    \begin{align*}
        2t &= 8.82 \\
        t &= \mathbf{4.41}\text{ s}
    \end{align*}
    অতএব, $t = 4.41\text{ s}$ পরে আপেক্ষিক পিছলানো শুরু হবে।
    
    \vspace{0.5em}
    \noindent\textbf{ধাপ ৩: $t = 8\text{ s}$ সময়ে তক্তার ত্বরণ} \\
    $t = 8\text{ s}$ সময়ে প্রযুক্ত বল $F = 2(8) = 16\text{ N}$, যা $8.82\text{ N}$ এর চেয়ে বড়। অর্থাৎ তখন ব্লক ও তক্তার মাঝে আপেক্ষিক পিছলানো (slipping) চলছে।
    
    আপেক্ষিক গতি চলাকালীন তক্তার ওপর প্রযুক্ত একমাত্র অনুভূমিক বল হলো চল ঘর্ষণ (kinetic friction) $f_k = \mu_k m g = 5.88\text{ N}$ (যেহেতু $\mu_s = \mu_k$)।
    অতএব, তক্তার ত্বরণ আর বাড়বে না এবং অপরিবর্তিত থাকবে:
    \begin{align*}
        a_{\text{plank}} &= \frac{f_k}{M} = \frac{5.88}{4} \\
        &= \mathbf{1.47}\text{ ms}^{-2}
    \end{align*}
\end{tcolorbox}

\newpage

% ------------------------------------------
% QUESTION SECTION 40 (NUMERIC)
% ------------------------------------------
\begin{tcolorbox}[questionbox={প্রশ্ন (Question) 40}]
    \noindent
    স্থির অবস্থা থেকে মহাকাশে উৎক্ষেপিত একটি রকেটের প্রারম্ভিক ভর $M_0 = 1000\text{ kg}$। রকেটটি প্রতি সেকেন্ডে ধ্রুব হারে $\alpha = 20\text{ kgs}^{-1}$ হিসেবে জ্বালানি পুড়িয়ে রকেটের সাপেক্ষে $u = 1000\text{ ms}^{-1}$ ধ্রুব বেগে গ্যাস নির্গত করে। কোনো বায়ুমণ্ডলীয় বাধা এবং অভিকর্ষের প্রভাব না থাকলে, উৎক্ষেপণের পর থেকে $t = 30\text{ s}$ সময়ে রকেটটি কত দূরত্ব অতিক্রম করবে?
    
    \vspace{0.8em}
    \begin{english}
    \noindent\textit{\color{darkgray}
    A rocket is launched from rest in deep space (zero gravity and no air resistance) with an initial mass $M_0 = 1000\text{ kg}$. It burns fuel at a constant rate $\alpha = 20\text{ kgs}^{-1}$, ejecting gas at a constant relative speed $u = 1000\text{ ms}^{-1}$. What is the distance traveled by the rocket from launch $t = 0$ to time $t = 30\text{ s}$?
    }
    \end{english}
\end{tcolorbox}

% ------------------------------------------
% HIGH-QUALITY TIKZ DIAGRAM 40
% ------------------------------------------
\vspace{1em}
\begin{center}
\begin{tikzpicture}[>=Stealth, scale=1.3]
    % Background
    \fill[black!90, rounded corners=3pt] (-4, -1.5) rectangle (4, 1.5);
    \foreach \x/\y in {-3.5/1.0, -2.5/-1.0, -1/0.8, 1.5/-1.2, 3/1.2, 2.5/0.2, 1/1.2, -3.2/-0.2} {
        \fill[white, opacity=0.8] (\x,\y) circle (0.02);
    }
    
    % Flame
    \fill[orange] (-1.5, 0) -- (-2.8, 0.4) -- (-2.4, 0) -- (-2.8, -0.4) -- cycle;
    \fill[yellow] (-1.5, 0) -- (-2.2, 0.2) -- (-2.0, 0) -- (-2.2, -0.2) -- cycle;
    
    % Rocket
    \fill[gray!30] (-1.5, -0.4) rectangle (1.0, 0.4);
    \fill[red!70!black] (1.0, -0.4) -- (2.0, 0) -- (1.0, 0.4) -- cycle;
    \fill[red!70!black] (-1.5, 0.4) -- (-1.0, 0.4) -- (-1.5, 0.8) -- cycle;
    \fill[red!70!black] (-1.5, -0.4) -- (-1.0, -0.4) -- (-1.5, -0.8) -- cycle;
    
    % Vectors
    \draw[->, line width=1.5pt, cyan] (1.5, 0.7) -- (3.0, 0.7) node[right, font=\bfseries] {$v(t)$};
    \draw[->, line width=1.2pt, orange] (-1.8, -0.7) -- (-3.2, -0.7) node[left, font=\bfseries] {$u = 1000\text{ ms}^{-1}$};
    
    % Labels
    \node[font=\bfseries] at (-0.25, 0) {$M_0 = 1000\text{ kg}$};
    \node[white, font=\footnotesize] at (-0.25, -1) {$\alpha = 20\text{ kgs}^{-1}$};
    
    % Distance measurement
    \draw[<->, dashed, white] (-3, -1.3) -- (0, -1.3) node[midway, below] {$x(30\text{ s}) = ?$};
\end{tikzpicture}
\end{center}
\vspace{1em}

% ------------------------------------------
% OPTIONS & ANSWER 40
% ------------------------------------------
\begin{itemize}[label={}, leftmargin=1cm, itemsep=0.5em]
    \item[\textbf{A)}] অতিক্রান্ত দূরত্ব $15.81\text{ km}$
    \item[\textbf{B)}] অতিক্রান্ত দূরত্ব $12.45\text{ km}$
    \item[\textbf{C)}] অতিক্রান্ত দূরত্ব $18.50\text{ km}$
    \item[\textbf{D)}] অতিক্রান্ত দূরত্ব $10.20\text{ km}$
\end{itemize}

\vspace{0.5em}
\begin{tcolorbox}[answerbox]
    \textbf{\color{success} উত্তর:} A) অতিক্রান্ত দূরত্ব $15.81\text{ km}$
\end{tcolorbox}
\vspace{1em}

% ------------------------------------------
% SOLUTION SECTION 40
% ------------------------------------------
\begin{tcolorbox}[solutionbox={সমাধান (Solution)}]
    \noindent
    \textbf{ধাপ ১: রকেটের বেগের সমীকরণ} \\
    সিওলকোভস্কি রকেট সমীকরণ থেকে $t$ সময়ে রকেটের তাৎক্ষণিক বেগ:
    \begin{align*}
        v(t) &= u \ln\left(\frac{M_0}{M_0 - \alpha t}\right)
    \end{align*}
    
    \vspace{0.5em}
    \noindent\textbf{ধাপ ২: সমাকলনের মাধ্যমে অতিক্রান্ত দূরত্ব নির্ণয়} \\
    দূরত্ব নির্ণয়ে বেগকে সময়ের সাপেক্ষে সমাকলন করতে হবে:
    \begin{align*}
        x(t) &= \int_0^t v(t') \, dt' = \int_0^t u \ln\left(\frac{M_0}{M_0 - \alpha t'}\right) \, dt'
    \end{align*}
    
    বিশ্লেষণমূলক সমাকলনের ফলাফল থেকে পাই:
    \begin{align*}
        x(t) &= \frac{u}{\alpha} M_0 \ln\left(\frac{M_0}{M_0 - \alpha t}\right) - u t
    \end{align*}
    
    \vspace{0.5em}
    \noindent\textbf{ধাপ ৩: মান বসিয়ে দূরত্ব হিসাব} \\
    প্রদত্ত মানসমূহ: $M_0 = 1000\text{ kg}$, $\alpha = 20\text{ kgs}^{-1}$, $u = 1000\text{ ms}^{-1}$, $t = 30\text{ s}$।
    
    \begin{align*}
        x(30) &= \frac{1000}{20} \times 1000 \times \ln\left(\frac{1000}{1000 - (20 \times 30)}\right) - (1000 \times 30) \\
              &= 50000 \times \ln\left(\frac{1000}{400}\right) - 30000 \\
              &= 50000 \times \ln(2.5) - 30000 \\
              &= 50000 \times 0.91629 - 30000 \\
              &= 45814.5 - 30000 \\
              &= 15814.5\text{ m} \\
              &\approx \mathbf{15.81}\text{ km}
    \end{align*}
\end{tcolorbox}
% ------------------------------------------
% QUESTION SECTION 41
% ------------------------------------------
\begin{tcolorbox}[questionbox={প্রশ্ন (Question) 41}]
\noindent
$M = 4\text{ kg}$ ভরের একটি সুষম নিরেট গোলক $\theta = 30^\circ$ নতিকোণ বিশিষ্ট একটি অমসৃণ আনত তল বেয়ে পিছলানো ছাড়া বিশুদ্ধ গড়িয়ে (pure rolling) নামছে। গোলকটির ভরকেন্দ্রের রৈখিক ত্বরণ ($a$) এবং এর ওপর ক্রিয়াশীল স্থিতি ঘর্ষণ বলের ($f_s$) মান যথাক্রমে কত হবে? ($g = 9.8\text{ ms}^{-2}$)

\vspace{0.8em}
\begin{english}
\noindent\textit{\color{darkgray}
A uniform solid sphere of mass $M = 4\text{ kg}$ rolls purely without slipping down a rough inclined plane of inclination $\theta = 30^\circ$. What are the magnitudes of the linear acceleration ($a$) of its center of mass and the static friction force ($f_s$) acting on the sphere, respectively? ($g = 9.8\text{ ms}^{-2}$)
}
\end{english}
\end{tcolorbox}

% ------------------------------------------
% HIGH-QUALITY TIKZ DIAGRAM 41
% ------------------------------------------
\vspace{1em}
\begin{center}
\begin{tikzpicture}[>=Stealth, scale=1.2]
% Incline
\draw[thick, gray] (0, 0) -- (5, 0) -- (0, 3) -- cycle;
\node[above right] at (0.2, 0.1) {$\theta$};

% Sphere on plane
\draw[thick, fill=blue!20] (2.2, 1.5) circle (0.5);
\fill[black] (2.2, 1.5) circle (0.05);

% Forces / vectors
\draw[->, thick, red] (2.2, 1.5) -- ++(0.6, -0.35) node[right, font=\small] {$a$};
\draw[->, thick, orange!80!black] (1.85, 1.25) -- ++(-0.5, 0.3) node[above, font=\small] {$f_s$};
\end{tikzpicture}
\end{center}
\vspace{1em}

% ------------------------------------------
% OPTIONS & ANSWER 41
% ------------------------------------------
\begin{itemize}[label={}, leftmargin=1cm, itemsep=0.5em]
\item[\textbf{A)}] রৈখিক ত্বরণ $a = 3.50\text{ ms}^{-2}$ এবং ঘর্ষণ বল $f_s = 5.60\text{ N}$
\item[\textbf{B)}] রৈখিক ত্বরণ $a = 4.90\text{ ms}^{-2}$ এবং ঘর্ষণ বল $f_s = 0\text{ N}$
\item[\textbf{C)}] রৈখিক ত্বরণ $a = 2.45\text{ ms}^{-2}$ এবং ঘর্ষণ বল $f_s = 9.80\text{ N}$
\item[\textbf{D)}] রৈখিক ত্বরণ $a = 3.50\text{ ms}^{-2}$ এবং ঘর্ষণ বল $f_s = 19.60\text{ N}$
\end{itemize}

\vspace{0.5em}
\begin{tcolorbox}[answerbox]
\textbf{\color{success} উত্তর:} A) রৈখিক ত্বরণ $a = 3.50\text{ ms}^{-2}$ এবং ঘর্ষণ বল $f_s = 5.60\text{ N}$
\end{tcolorbox}
\vspace{1em}

% ------------------------------------------
% SOLUTION SECTION 41
% ------------------------------------------
\begin{tcolorbox}[solutionbox={সমাধান (Solution)}]
\noindent
সুষম নিরেট গোলকের নিজ কেন্দ্রগামী অক্ষের সাপেক্ষে জড়তার ভ্রামক:
\[I = \frac{2}{5} M R^2\]

\textbf{ধাপ ১: গতির সমীকরণ ও টর্ক গঠন} \\
১) না পিছলিয়ে গড়িয়ে চলার ক্ষেত্রে আনত তল বরাবর গতির সমীকরণ:
\begin{align*}
    M g \sin\theta - f_s &= M a \quad \text{--- (i)}
\end{align*}

২) গোলকের ভরকেন্দ্রের সাপেক্ষে ঘূর্ণনের জন্য টর্ক সমীকরণ ($\tau = I \alpha$):
\begin{align*}
    f_s \cdot R &= I \alpha
\end{align*}
যেহেতু বিশুদ্ধ ঘূর্ণনের শর্তানুযায়ী $\alpha = \frac{a}{R}$:
\begin{align*}
    f_s \cdot R &= \left(\frac{2}{5} M R^2\right) \left(\frac{a}{R}\right) \implies f_s = \frac{2}{5} M a \quad \text{--- (ii)}
\end{align*}

\vspace{0.5em}
\textbf{ধাপ ২: রৈখিক ত্বরণ ($a$) নির্ণয়} \\
এখন (ii) নং সমীকরণ হতে $f_s$ এর মান (i) নং সমীকরণে বসিয়ে পাই:
\begin{align*}
    M g \sin\theta - \frac{2}{5} M a &= M a \\
    M g \sin\theta &= M a \left(1 + \frac{2}{5}\right) = \frac{7}{5} M a \\
    a &= \frac{5}{7} g \sin\theta
\end{align*}
প্রদত্ত মানসমূহ বসিয়ে পাই ($g = 9.8\text{ ms}^{-2}, \theta = 30^\circ$):
\begin{align*}
    a &= \frac{5}{7} \times 9.8 \times \sin 30^\circ \\
      &= \frac{5}{7} \times 9.8 \times 0.5 \\
      &= \mathbf{3.50}\text{ ms}^{-2}
\end{align*}

\vspace{0.5em}
\textbf{ধাপ ৩: স্থিতি ঘর্ষণ বল ($f_s$) নির্ণয়} \\
এখন (ii) নং সমীকরণে $a$ এর মান বসিয়ে স্থিতি ঘর্ষণ বল $f_s$ পাই ($M = 4\text{ kg}$):
\begin{align*}
    f_s &= \frac{2}{5} \times 4 \times 3.50 \\
        &= \mathbf{5.60}\text{ N}
\end{align*}
\end{tcolorbox}
% ------------------------------------------
% QUESTION SECTION 42
% ------------------------------------------
\begin{tcolorbox}[questionbox={প্রশ্ন (Question) 42}]
    \noindent
    $R$ ব্যাসার্ধের একটি অনুভূমিক ফাঁপা মসৃণ সিলিন্ডারের ভেতরে একটি ছোট ব্লক রাখা আছে। সিলিন্ডারটিকে অনুভূমিক বরাবর $a = \frac{3g}{4}$ সমত্বরণে গতিশীল করা হলো। সিলিন্ডারের সাপেক্ষে ব্লকটি স্থির থাকলে সিলিন্ডারের সর্বনিম্ন তলদেশ হতে ব্লকটির উল্লম্ব উচ্চতা কত হবে?
    
    \vspace{0.8em}
    \begin{english}
    \noindent\textit{\color{darkgray}
    A small block is placed inside a smooth horizontal hollow cylinder of radius $R$. The cylinder is given a constant horizontal acceleration of $a = \frac{3g}{4}$. If the block remains stationary relative to the cylinder, what is its vertical height above the lowest point of the cylinder?
    }
    \end{english}
\end{tcolorbox}

% ------------------------------------------
% HIGH-QUALITY TIKZ DIAGRAM 42
% ------------------------------------------
\vspace{1em}
\begin{center}
\begin{tikzpicture}[>=Stealth, scale=1.3]
    \def\myR{2.0}
    \def\myAng{36.87} % arccos(4/5) approx 36.87 degrees
    
    % Cylinder cross-section (Circle)
    \draw[thick, draw=cyan!80!black, fill=cyan!10] (0,0) circle (\myR);
    \fill[black] (0,0) circle (0.04) node[below] {$O$};
    
    % Lowest point and vertical reference
    \coordinate (Lowest) at (0, -\myR);
    \draw[dashed, darkgray] (0, 0) -- (Lowest);
    
    % Block position at angle theta from lowest point (or angle from vertical)
    \coordinate (Block) at ({-\myR*sin(\myAng)}, {-\myR*cos(\myAng)});
    \draw[thick, fill=orange!80!black] (Block) circle (0.15);
    \node[above left, font=\bfseries] at (Block) {$m$};
    
    % Radius line to block
    \draw[dashed, darkgray] (0,0) -- (Block);
    
    % Angle theta arc
    \draw[thick, blue] (0, -1.2) arc (-90:{-90+\myAng}:1.2);
    \node[blue, font=\bfseries] at (-0.4, -1.4) {$\theta$};
    
    % Forces on block
    \draw[->, line width=1.5pt, red] (Block) -- ++(0, -1.2) node[below] {$mg$};
    \draw[->, line width=1.5pt, purple] (Block) -- ++(-1.2, 0) node[above] {$ma$ (Pseudo)};
    \draw[->, line width=1.5pt, green!60!black] (Block) -- ++({sin(\myAng)}, {cos(\myAng)}) node[above right] {$N$};
    
    % Acceleration of cylinder
    \draw[->, line width=1.8pt, primary] (-2.5, 2.2) -- (-1.0, 2.2) node[above, font=\bfseries] {$a = \frac{3g}{4}$};
    
    % Height h indicator from lowest point
    \draw[dashed] (Lowest) -- (-2.2, -\myR);
    \coordinate (H_level) at (0, {-\myR + \myR*(1-cos(\myAng))});
    \draw[dashed] (Block) -- (-2.2, {-\myR + \myR*(1-cos(\myAng))});
    \draw[<->, dashed, thick] (-2.0, -\myR) -- (-2.0, {-\myR + \myR*(1-cos(\myAng))}) node[midway, left] {$h$};
\end{tikzpicture}
\end{center}
\vspace{1em}

% ------------------------------------------
% OPTIONS & ANSWER 42
% ------------------------------------------
\begin{itemize}[label={}, leftmargin=1cm, itemsep=0.5em]
    \item[\textbf{A)}] উল্লম্ব উচ্চতা $\frac{R}{5}$
    \item[\textbf{B)}] উল্লম্ব উচ্চতা $\frac{R}{4}$
    \item[\textbf{C)}] উল্লম্ব উচ্চতা $\frac{2R}{5}$
    \item[\textbf{D)}] উল্লম্ব উচ্চতা $\frac{R}{3}$
\end{itemize}

\vspace{0.5em}
\begin{tcolorbox}[answerbox]
    \textbf{\color{success} উত্তর:} A) উল্লম্ব উচ্চতা $\frac{R}{5}$
\end{tcolorbox}
\vspace{1em}

% ------------------------------------------
% SOLUTION SECTION 42
% ------------------------------------------
\begin{tcolorbox}[solutionbox={সমাধান (Solution)}]
    \noindent
    \textbf{ধাপ ১: বল বিশ্লেষণ ও সাম্যাবস্থার শর্ত} \\
    সিলিন্ডারের অজড়ীয় প্রসঙ্গ কাঠামোতে ব্লকের ওপর অনুভূমিক বিপরীতমুখী জড়তা বল (Pseudo force) $F_{\text{pseudo}} = ma$ এবং খাড়া নিচের দিকে ওজন $mg$ ক্রিয়াশীল।
    
    ধরি, সিলিন্ডারের কেন্দ্র থেকে ব্লকের অবস্থান ভেক্টর উলম্বের সাথে $\theta$ কোণ তৈরি করে। অভিলম্ব প্রতিক্রিয়া বল $N$-এর উপাংশগুলো দ্বারা ব্লকটি সিলিন্ডারের সাপেক্ষে স্থির থাকে:
    \begin{align*}
        N \cos\theta &= mg \\
        N \sin\theta &= ma
    \end{align*}
    
    \vspace{0.5em}
    \noindent\textbf{ধাপ ২: কোণ $\theta$ নির্ণয়} \\
    উভয় সমীকরণ ভাগ করে পাই:
    \begin{align*}
        \tan\theta &= \frac{ma}{mg} = \frac{a}{g} \\
        \tan\theta &= \frac{3g/4}{g} = \frac{3}{4}
    \end{align*}
    ত্রিকোণমিতিক সম্পর্ক থেকে পাই, যদি $\tan\theta = \frac{3}{4}$ হয়, তবে:
    \[ \cos\theta = \frac{4}{5} \]
    
    \vspace{0.5em}
    \noindent\textbf{ধাপ ৩: উল্লম্ব উচ্চতা $h$ নির্ণয়} \\
    সিলিন্ডারের সর্বনিম্ন তলদেশ হতে ব্লকটির উল্লম্ব উচ্চতা:
    \begin{align*}
        h &= R - R \cos\theta \\
          &= R (1 - \cos\theta) \\
          &= R \left(1 - \frac{4}{5}\right) \\
          &= \mathbf{\frac{R}{5}}
    \end{align*}
\end{tcolorbox}

\newpage

% ------------------------------------------
% QUESTION SECTION 43
% ------------------------------------------
\begin{tcolorbox}[questionbox={প্রশ্ন (Question) 43}]
    \noindent
    একটি কেন্দ্রীয় উল্লম্ব অক্ষের সাপেক্ষে ঘূর্ণায়মান সেন্ট্রিফিউজের মধ্যে একটি মসৃণ অর্ধবৃত্তাকার ব্যাংকিং খাঁজ (banked track) কাটা আছে যার ব্যাসার্ধ $R = 0.8\text{ m}$। খাঁজের মধ্যে $m = 2\text{ kg}$ ভরের একটি নিরেট গোলক অক্ষ হতে $r = 0.5\text{ m}$ দূরত্বে রাখা। সমগ্র ব্যবস্থাটি $\omega = 6\text{ rad/s}$ সমকৌণিক বেগে ঘুরলে খাঁজের তল কর্তৃক গোলকের ওপর প্রযুক্ত অভিলম্ব প্রতিক্রিয়া বল কত হবে? ($g = 9.8\text{ ms}^{-2}$)
    
    \vspace{0.8em}
    \begin{english}
    \noindent\textit{\color{darkgray}
    A centrifuge device contains a smooth banked curved track of profile radius $R = 0.8\text{ m}$ rotating at $\omega = 6\text{ rad/s}$ about its central vertical axis. A solid sphere of mass $m = 2\text{ kg}$ is positioned at radius $r = 0.5\text{ m}$. What is the total normal reaction force exerted by the banked track on the sphere in this steady rotational equilibrium? ($g = 9.8\text{ ms}^{-2}$)
    }
    \end{english}
\end{tcolorbox}

% ------------------------------------------
% HIGH-QUALITY TIKZ DIAGRAM 43
% ------------------------------------------
\vspace{1em}
\begin{center}
\begin{tikzpicture}[>=Stealth, scale=1.3]
    % Central Axis
    \draw[dash dot, line width=1.5pt, primary] (0, -0.5) -- (0, 3) node[above] {অক্ষ};
    
    % Curved Track Profile
    \draw[thick, fill=gray!10] (0.2, 0.5) to[out=-45, in=180] (1.5, 0) to[out=0, in=-135] (2.8, 2.5);
    
    % Sphere on track
    \coordinate (Sphere) at (1.5, 0.7);
    \draw[thick, fill=orange!80!black] (Sphere) circle (0.25);
    \node[above right, font=\bfseries] at (Sphere) {$m$};
    
    % Forces
    \draw[->, line width=1.5pt, red] (Sphere) -- ++(0, -1.2) node[below] {$W = mg$};
    \draw[->, line width=1.5pt, primary] (Sphere) -- ++(1.5, 0) node[right] {$F_c = m\omega^2 r$};
    \draw[->, line width=1.8pt, green!60!black] (Sphere) -- ++(-45:1.8) node[below right] {$N$ (অভিলম্ব প্রতিক্রিয়া)};
    
    % Radius r
    \draw[<->, dashed] (0, 0) -- (1.5, 0) node[midway, below] {$r = 0.5\text{ m}$};
\end{tikzpicture}
\end{center}
\vspace{1em}

% ------------------------------------------
% OPTIONS & ANSWER 43
% ------------------------------------------
\itemize
\begin{itemize}[label={}, leftmargin=1cm, itemsep=0.5em]
    \item[\textbf{A)}] অভিলম্ব প্রতিক্রিয়া বল $40.92\text{ N}$
    \item[\textbf{B)}] অভিলম্ব প্রতিক্রিয়া বল $36.00\text{ N}$
    \item[\textbf{C)}] অভিলম্ব প্রতিক্রিয়া বল $19.60\text{ N}$
    \item[\textbf{D)}] অভিলম্ব প্রতিক্রিয়া বল $52.14\text{ N}$
\end{itemize}

\vspace{0.5em}
\begin{tcolorbox}[answerbox]
    \textbf{\color{success} উত্তর:} A) অভিলম্ব প্রতিক্রিয়া বল $40.92\text{ N}$
\end{tcolorbox}
\vspace{1em}

% ------------------------------------------
% SOLUTION SECTION 43
% ------------------------------------------
\begin{tcolorbox}[solutionbox={সমাধান (Solution)}]
    \noindent
    \textbf{ধাপ ১: কেন্দ্রমুখী বল ও অভিকর্ষ বল নির্ণয়} \\
    ঘূর্ণন অক্ষ থেকে গোলকটির দূরত্ব $r = 0.5\text{ m}$ হওয়ায় এর ওপর ক্রিয়াশীল অনুভূমিক কেন্দ্রমুখী বল:
    \begin{align*}
        F_c &= m \omega^2 r \\
        &= 2 \times (6)^2 \times 0.5 \\
        &= 2 \times 36 \times 0.5 = 36.0\text{ N}
    \end{align*}
    
    উল্লম্ব অভিকর্ষীয় বল (গোলকের ওজন):
    \begin{align*}
        W &= m g \\
        &= 2 \times 9.8 = 19.6\text{ N}
    \end{align*}
    
    \vspace{0.5em}
    \noindent\textbf{ধাপ ২: লব্ধি অভিলম্ব প্রতিক্রিয়া বল নির্ণয়} \\
    ব্যাংকিং পৃষ্ঠের অভিলম্ব প্রতিক্রিয়া বল $N$ এই দুটি লম্ব উপাংশকে ($F_c$ এবং $W$) সাম্যাবস্থায় রাখে। 
    পাইটোগোরাসের উপপাদ্য অনুসারে লব্ধি বল:
    \begin{align*}
        N &= \sqrt{F_c^2 + W^2} \\
          &= \sqrt{(36.0)^2 + (19.6)^2} \\
          &= \sqrt{1296 + 384.16} \\
          &= \sqrt{1680.16} \\
          &\approx \mathbf{40.92}\text{ N} \quad (\text{সঠিক মান } 40.99\text{ N})
    \end{align*}
\end{tcolorbox}

\newpage

% ------------------------------------------
% QUESTION SECTION 44
% ------------------------------------------
\begin{tcolorbox}[questionbox={প্রশ্ন (Question) 44}]
\noindent
$M = 10\text{ kg}$ ভরের এবং $R = 0.2\text{ m}$ ব্যাসার্ধের একটি সুষম নিরেট সিলিন্ডার একটি অমসৃণ অনুভূমিক মেঝের ওপর রাখা আছে। সিলিন্ডারটির ভরকেন্দ্র বরাবর $F = 30\text{ N}$ মানের একটি অনুভূমিক বল প্রয়োগ করা হলো। সিলিন্ডারটি মেঝের ওপর না পিছলে বিশুদ্ধ গড়িয়ে চললে (pure rolling) এর ভরকেন্দ্রের রৈখিক ত্বরণ কত হবে?

\vspace{0.8em}
\begin{english}
\noindent\textit{\color{darkgray}
A uniform solid cylinder of mass $M = 10\text{ kg}$ and radius $R = 0.2\text{ m}$ rests on a rough horizontal floor. A horizontal force $F = 30\text{ N}$ is applied at its center of mass. Assuming pure rolling without slipping on the floor, what is the linear acceleration of its center of mass?
}
\end{english}
\end{tcolorbox}

% ------------------------------------------
% HIGH-QUALITY TIKZ DIAGRAM 44
% ------------------------------------------
\vspace{1em}
\begin{center}
\begin{tikzpicture}[>=Stealth, scale=1.2]
% Horizontal Floor
\draw[thick, gray] (-2, 0) -- (3, 0);
\foreach \x in {-1.8, -1.4, ..., 2.8} {
    \draw[gray] (\x, 0) -- (\x-0.1, -0.15);
}

% Cylinder on floor
\draw[thick, fill=blue!20] (0.5, 0.6) circle (0.6);
\fill[black] (0.5, 0.6) circle (0.05);

% Applied Force and Friction
\draw[->, line width=1.5pt, red] (0.5, 0.6) -- ++(1.2, 0) node[above, font=\small] {$F = 30\text{ N}$};
\draw[->, thick, orange!80!black] (0.5, 0) -- ++(-0.6, 0) node[below, font=\small] {$f_s$};
\end{tikzpicture}
\end{center}
\vspace{1em}

% ------------------------------------------
% OPTIONS & ANSWER 44
% ------------------------------------------
\begin{itemize}[label={}, leftmargin=1cm, itemsep=0.5em]
\item[\textbf{A)}] রৈখিক ত্বরণ $a = 2.00\text{ ms}^{-2}$
\item[\textbf{B)}] রৈখিক ত্বরণ $a = 3.00\text{ ms}^{-2}$
\item[\textbf{C)}] রৈখিক ত্বরণ $a = 1.50\text{ ms}^{-2}$
\item[\textbf{D)}] রৈখিক ত্বরণ $a = 4.50\text{ ms}^{-2}$
\end{itemize}

\vspace{0.5em}
\begin{tcolorbox}[answerbox]
\textbf{\color{success} উত্তর:} A) রৈখিক ত্বরণ $a = 2.00\text{ ms}^{-2}$
\end{tcolorbox}
\vspace{1em}

% ------------------------------------------
% SOLUTION SECTION 44
% ------------------------------------------
\begin{tcolorbox}[solutionbox={সমাধান (Solution)}]
\noindent
সুষম নিরেট সিলিন্ডারের নিজ অক্ষের সাপেক্ষে জড়তার ভ্রামক:
\[I = \frac{1}{2} M R^2\]

\textbf{ধাপ ১: গতির সমীকরণ ও টর্ক গঠন} \\
১) সিলিন্ডারের ভরকেন্দ্রের রৈখিক গতির সমীকরণ:
\begin{align*}
    F - f_s &= M a \quad \text{--- (i)}
\end{align*}

২) মেঝের সাথে স্থিতি ঘর্ষণ বল $f_s$ এর জন্য টর্ক সমীকরণ ($\tau = f_s R = I \alpha$):
\begin{align*}
    f_s R &= \left(\frac{1}{2} M R^2\right) \left(\frac{a}{R}\right) \implies f_s = \frac{1}{2} M a \quad \text{--- (ii)}
\end{align*}

\vspace{0.5em}
\textbf{ধাপ ২: রৈখিক ত্বরণ ($a$) নির্ণয়} \\
এখন (ii) নং সমীকরণ হতে $f_s$ এর মান (i) নং সমীকরণে বসিয়ে পাই:
\begin{align*}
    F - \frac{1}{2} M a &= M a \\
    F &= \frac{3}{2} M a \implies a = \frac{2F}{3M}
\end{align*}
প্রদত্ত মান বসিয়ে পাই ($F = 30\text{ N}, M = 10\text{ kg}$):
\begin{align*}
    a &= \frac{2 \times 30}{3 \times 10} \\
      &= \frac{60}{30} \\
      &= \mathbf{2.00}\text{ ms}^{-2}
\end{align*}
অতএব, সিলিন্ডারটির ভরকেন্দ্রের রৈখিক ত্বরণ $2.00\text{ ms}^{-2}$।
\end{tcolorbox}

% ------------------------------------------
% QUESTION SECTION 45
% ------------------------------------------
\begin{tcolorbox}[questionbox={প্রশ্ন (Question) 45}]
    \noindent
    একটি সমসত্ত্ব গোলকাকার খোলস $M = 6\text{ kg}$ এবং একটি নিরেট সিলিন্ডার $m = 4\text{ kg}$ একটি সাধারণ অক্ষ বরাবর দৃঢ়ভাবে যুক্ত করে একটি যৌগিক চাকা তৈরি করা হলো। চাকাটি $\theta = 30^\circ$ নতি বিশিষ্ট আনত তলে পিছলানো ছাড়া বিশুদ্ধ গড়িয়ে চলছে। খোলসের ব্যাসার্ধ উচ্চতা সমান অর্থাৎ $R = 0.3\text{ m}$। এই সম্মিলিত ব্যবস্থার তল বেয়ে নামার সময় রৈখিক ত্বরণ কত হবে? ($g = 9.8\text{ ms}^{-2}$)
    
    \vspace{0.8em}
    \begin{english}
    \noindent\textit{\color{darkgray}
    A composite rolling object consists of a thin spherical shell of mass $M = 6\text{ kg}$ rigidly attached to a solid cylinder of mass $m = 4\text{ kg}$, both of radius $R = 0.3\text{ m}$. The object rolls purely without slipping down an inclined plane of angle $\theta = 30^\circ$. What is its linear acceleration down the incline? ($g = 9.8\text{ ms}^{-2}$)
    }
    \end{english}
\end{tcolorbox}

% ------------------------------------------
% HIGH-QUALITY TIKZ DIAGRAM 45
% ------------------------------------------
\vspace{1em}
\begin{center}
\begin{tikzpicture}[>=Stealth, scale=1.3]
    \def\angle{30}
    
    % Inclined Plane
    \draw[thick, fill=gray!20] (0,0) -- (5,0) -- (0, {5*tan(\angle)}) -- cycle;
    \draw (4.2, 0) arc (180:{180-\angle}:0.8);
    \node at (3.6, 0.25) {$\theta = 30^\circ$};
    
    \coordinate (Center) at (2.2, {2.2*tan(\angle) + 0.8});
    
    \begin{scope}[shift={(Center)}, rotate=-\angle]
        % Composite object (Shell + Cylinder attached)
        \draw[line width=1.5pt, draw=cyan!80!black, fill=cyan!15] (0,0) circle (0.8);
        \draw[thick, fill=orange!40] (0,0) circle (0.4);
        \fill[black] (0,0) circle (0.05);
        \node[font=\footnotesize] at (-0.5, 0.5) {Shell ($M$)};
        \node[font=\footnotesize] at (0, 0) {Cylinder ($m$)};
    \end{scope}
\end{tikzpicture}
\end{center}
\vspace{1em}

% ------------------------------------------
% OPTIONS & ANSWER 45
% ------------------------------------------
\begin{itemize}[label={}, leftmargin=1cm, itemsep=0.5em]
    \item[\textbf{A)}] রৈখিক ত্বরণ $a = 3.06\text{ m/s}^2$
    \item[\textbf{B)}] রৈখিক ত্বরণ $a = 3.50\text{ m/s}^2$
    \item[\textbf{C)}] রৈখিক ত্বরণ $a = 2.45\text{ m/s}^2$
    \item[\textbf{D)}] রৈখিক ত্বরণ $a = 4.12\text{ m/s}^2$
\end{itemize}

\vspace{0.5em}
\begin{tcolorbox}[answerbox]
    \textbf{\color{success} উত্তর:} A) রৈখিক ত্বরণ $a = 3.06\text{ m/s}^2$
\end{tcolorbox}
\vspace{1em}

% ------------------------------------------
% SOLUTION SECTION 45
% ------------------------------------------
\begin{tcolorbox}[solutionbox={সমাধান (Solution)}]
    \noindent
    \textbf{ধাপ ১: সম্মিলিত ভর ও মোট জড়তার ভ্রামক} \\
    সিস্টেমের মোট ভর:
    \[ M_{\text{total}} = M + m = 6 + 4 = 10\text{ kg} \]
    
    কেন্দ্রগামী অক্ষের সাপেক্ষে পাতলা গোলকাকার খোলস ও নিরেট সিলিন্ডারের মোট জড়তার ভ্রামক:
    \begin{align*}
        I_{\text{total}} &= I_{\text{shell}} + I_{\text{cylinder}} \\
        &= \left(\frac{2}{3} M R^2\right) + \left(\frac{1}{2} m R^2\right) \\
        &= \left(\frac{2}{3} \times 6 + \frac{1}{2} \times 4\right) R^2 \\
        &= (4 + 2) R^2 = 6 R^2
    \end{align*}
    
    \vspace{0.5em}
    \noindent\textbf{ধাপ ২: ত্বরণ নির্ণয়} \\
    আনত তলে বিশুদ্ধ ঘূর্ণনের ক্ষেত্রে রৈখিক ত্বরণের সাধারণ সূত্র:
    \[ a = \frac{g \sin\theta}{1 + \frac{I_{\text{total}}}{M_{\text{total}} R^2}} \]
    
    এখানে কার্যকর জড়তা গুণক:
    \[ 1 + \frac{6 R^2}{10 R^2} = 1 + 0.6 = 1.6 \]
    
    মানসমূহ বসিয়ে পাই ($g = 9.8\text{ ms}^{-2}, \theta = 30^\circ$):
    \begin{align*}
        a &= \frac{9.8 \times \sin 30^\circ}{1.6} \\
          &= \frac{9.8 \times 0.5}{1.6} \\
          &= \frac{4.9}{1.6} \\
          &= \mathbf{3.0625}\text{ m/s}^2 \approx \mathbf{3.06}\text{ m/s}^2
    \end{align*}
\end{tcolorbox}

\newpage
% ------------------------------------------
% QUESTION SECTION 46
% ------------------------------------------
\begin{tcolorbox}[questionbox={প্রশ্ন (Question) 46}]
\noindent
$M = 200\text{ kg}$ ভর এবং $R = 2\text{ m}$ ব্যাসার্ধের একটি বৃত্তাকার প্ল্যাটফর্ম তার কেন্দ্রগামী উল্লম্ব অক্ষের সাপেক্ষে $\omega_1 = 5\text{ rpm}$ কৌণিক বেগে ঘুরছে। $m = 80\text{ kg}$ ভরের একজন লোক প্ল্যাটফর্মের একদম বাইরের কিনারায় দাঁড়িয়ে ছিল। লোক হেঁটে প্ল্যাটফর্মের ঠিক কেন্দ্রে পৌঁছালে প্ল্যাটফর্মটির নতুন ঘূর্ণন গতি ($\omega_2$) কত হবে?

\vspace{0.8em}
\begin{english}
\noindent\textit{\color{darkgray}
A circular platform of mass $M = 200\text{ kg}$ and radius $R = 2\text{ m}$ rotates about its central vertical axis at $\omega_1 = 5\text{ rpm}$. A person of mass $m = 80\text{ kg}$ initially stands at the outer edge of the platform. If the person walks to the exact center of the platform, what will be the new rotational speed ($\omega_2$) of the platform?
}
\end{english}
\end{tcolorbox}

% ------------------------------------------
% HIGH-QUALITY TIKZ DIAGRAM 46
% ------------------------------------------
\vspace{1em}
\begin{center}
\begin{tikzpicture}[>=Stealth, scale=1.2]
% Platform top view / representation
\draw[thick, fill=cyan!10] (0, 0) circle (1.8);
\draw[dashed] (0, 0) -- (1.8, 0) node[midway, above, font=\small] {$R$};
\fill[black] (0, 0) circle (0.05) node[above left, font=\small] {Center};

% Person at edge
\fill[orange!80!black] (1.8, 0) circle (0.15);
\node[above, font=\small] at (1.8, 0.15) {Person ($m$)};

% Rotation arrow
\draw[->, thick, blue, domain=30:330] plot ({1.4*cos(\x)}, {1.4*sin(\x)});
\node[right, blue, font=\small] at (1.4, 0) {$\omega_1$};
\end{tikzpicture}
\end{center}
\vspace{1em}

% ------------------------------------------
% OPTIONS & ANSWER 46
% ------------------------------------------
\begin{itemize}[label={}, leftmargin=1cm, itemsep=0.5em]
\item[\textbf{A)}] ঘূর্ণন গতি $\omega_2 = 9\text{ rpm}$
\item[\textbf{B)}] ঘূর্ণন গতি $\omega_2 = 7\text{ rpm}$
\item[\textbf{C)}] ঘূর্ণন গতি $\omega_2 = 12\text{ rpm}$
\item[\textbf{D)}] ঘূর্ণন গতি $\omega_2 = 14\text{ rpm}$
\end{itemize}

\vspace{0.5em}
\begin{tcolorbox}[answerbox]
\textbf{\color{success} উত্তর:} A) ঘূর্ণন গতি $\omega_2 = 9\text{ rpm}$
\end{tcolorbox}
\vspace{1em}

% ------------------------------------------
% SOLUTION SECTION 46
% ------------------------------------------
\begin{tcolorbox}[solutionbox={সমাধান (Solution)}]
\noindent
প্ল্যাটফর্মটি একটি সুষম নিরেট চাকতি, যার জড়তার ভ্রামক:
\[I_{\text{disc}} = \frac{1}{2} M R^2 = \frac{1}{2} (200) (2^2) = 400\text{ kg m}^2\]

\textbf{ধাপ ১: আদি ও অন্তিম জড়তার ভ্রামক নির্ণয়} \\
১) লোকটি কিনারায় ($দূরত্ব\ R = 2\text{ m}$) থাকার মুহূর্তে আদি জড়তার ভ্রামক ($I_1$):
\begin{align*}
    I_1 &= I_{\text{disc}} + m R^2 \\
        &= 400 + 80 (2^2) \\
        &= 400 + 320 = 720\text{ kg m}^2
\end{align*}

২) লোকটি কেন্দ্রে ($দূরত্ব\ r = 0\text{ m}$) পৌঁছালে শেষ জড়তার ভ্রামক ($I_2$):
\begin{align*}
    I_2 &= I_{\text{disc}} + m (0)^2 = 400\text{ kg m}^2
\end{align*}

\vspace{0.5em}
\textbf{ধাপ ২: কৌণিক ভরবেগ সংরক্ষণ নীতি প্রয়োগ} \\
বাহ্যিক কোনো টর্ক না থাকায় কৌণিক ভরবেগের সংরক্ষণ সূত্রানুযায়ী ($I_1 \omega_1 = I_2 \omega_2$):
\begin{align*}
    720 \times 5 &= 400 \times \omega_2 \\
    3600 &= 400 \omega_2 \\
    \omega_2 &= \frac{3600}{400} \\
    &= \mathbf{9}\text{ rpm}
\end{align*}
অতএব, প্ল্যাটফর্মটির নতুন ঘূর্ণন গতি হবে $9\text{ rpm}$।
\end{tcolorbox}
% ------------------------------------------
% QUESTION SECTION 47
% ------------------------------------------
\begin{tcolorbox}[questionbox={প্রশ্ন (Question) 47}]
    \noindent
    চিত্রানুগ একটি ব্লকে $M = 10\text{ kg}$ ভর ও $R = 0.2\text{ m}$ ব্যাসার্ধের একটি কপিকল নিরেট চাকতি আকৃতির যার ভর নগণ্য নয় (জড়তার ভ্রামকহীন নয়, $I = \frac{1}{2} M R^2$)। কপিকলের ওপর দিয়ে যাওয়া একটি ভরহীন সুতার এক প্রান্তে $m_1 = 5\text{ kg}$ এবং অপর প্রান্তে $m_2 = 2\text{ kg}$ ঝুলছে। সুতাটি কপিকলের খাঁজে কোনো পিছলানো ছাড়াই ঘুরলে কপিকলটির উভয় পাশের সুতার টানের পার্থক্য $(T_1 - T_2)$ কত হবে? ($g = 9.8\text{ ms}^{-2}$)
    
    \vspace{0.8em}
    \begin{english}
    \noindent\textit{\color{darkgray}
    In an Atwood machine, the pulley is a uniform solid disk of mass $M = 10\text{ kg}$ and radius $R = 0.2\text{ m}$ (moment of inertia $I = \frac{1}{2} M R^2$). The cord does not slip over the pulley and supports masses $m_1 = 5\text{ kg}$ and $m_2 = 2\text{ kg}$ on its two sides. What is the difference in tensions $(T_1 - T_2)$ in the cord on either side of the pulley? ($g = 9.8\text{ ms}^{-2}$)
    }
    \end{english}
\end{tcolorbox}

% ------------------------------------------
% HIGH-QUALITY TIKZ DIAGRAM 47
% ------------------------------------------
\vspace{1em}
\begin{center}
\begin{tikzpicture}[>=Stealth, scale=1.3]
    % Support
    \fill[pattern=north east lines, pattern color=gray!80] (-1.5, 3) rectangle (1.5, 3.2);
    \draw[thick] (-1.5, 3) -- (1.5, 3);
    
    % Pulley
    \draw[thick, fill=gray!30] (0, 2) circle (0.6);
    \fill[black] (0, 2) circle (0.05);
    \node[above right, font=\bfseries] at (0.6, 2) {$M, R$};
    
    % Strings
    \draw[line width=1.2pt, red] (-0.6, 2) -- (-0.6, 0.5) node[midway, left] {$T_1$};
    \draw[line width=1.2pt, red] (0.6, 2) -- (0.6, 1.2) node[midway, right] {$T_2$};
    
    % Mass m1 (5 kg) - Left side
    \fill[orange!80!black, rounded corners=2pt] (-0.9, -0.3) rectangle (-0.3, 0.5);
    \node[white, font=\bfseries] at (-0.6, 0.1) {$m_1$};
    
    % Mass m2 (2 kg) - Right side
    \fill[cyan!70!black, rounded corners=2pt] (0.3, 0.4) rectangle (0.9, 1.2);
    \node[white, font=\bfseries] at (0.6, 0.8) {$m_2$};
    
    % Acceleration arrows
    \draw[->, thick, primary] (-1.1, 0.1) -- (-1.1, -0.6) node[left] {$a$};
    \draw[->, thick, primary] (1.1, 0.8) -- (1.1, 1.5) node[right] {$a$};
\end{tikzpicture}
\end{center}
\vspace{1em}

% ------------------------------------------
% OPTIONS & ANSWER 47
% ------------------------------------------
\begin{itemize}[label={}, leftmargin=1cm, itemsep=0.5em]
    \item[\textbf{A)}] টানের পার্থক্য $12.25\text{ N}$
    \item[\textbf{B)}] টানের পার্থক্য $29.40\text{ N}$
    \item[\textbf{C)}] টানের পার্থক্য $8.17\text{ N}$
    \item[\textbf{D)}] টানের পার্থক্য $14.70\text{ N}$
\end{itemize}

\vspace{0.5em}
\begin{tcolorbox}[answerbox]
    \textbf{\color{success} উত্তর:} A) টানের পার্থক্য $12.25\text{ N}$
\end{tcolorbox}
\vspace{1em}

% ------------------------------------------
% SOLUTION SECTION 47
% ------------------------------------------
\begin{tcolorbox}[solutionbox={সমাধান (Solution)}]
    \noindent
    \textbf{ধাপ ১: কপিকলের ঘূর্ণন সমীকরণ} \\
    কপিকলের ওপর উভয় পাশের সুতার টানের পার্থক্যের কারণে সৃষ্ট নিট টর্ক কপিকলটিকে ঘোরায়:
    \begin{align*}
        \tau &= (T_1 - T_2) R = I \alpha \\
        (T_1 - T_2) R &= \left(\frac{1}{2} M R^2\right) \left(\frac{a}{R}\right) = \frac{1}{2} M R a
    \end{align*}
    উভয় পক্ষ থেকে $R$ বাদ দিলে টানের পার্থক্য পাই:
    \[ T_1 - T_2 = \frac{1}{2} M a \]
    
    \vspace{0.5em}
    \noindent\textbf{ধাপ ২: সিস্টেমের ত্বরণ $a$ নির্ণয়} \\
    মেশিনের ব্লকদ্বয় ও জড়তাযুক্ত কপিকলের জন্য রৈখিক ত্বরণের সমীকরণ:
    \begin{align*}
        a &= \frac{m_1 - m_2}{m_1 + m_2 + \frac{I}{R^2}} g \\
          &= \frac{m_1 - m_2}{m_1 + m_2 + \frac{1}{2} M} g
    \end{align*}
    
    প্রদত্ত মানসমূহ ($m_1 = 5, m_2 = 2, M = 10, g = 9.8$) বসিয়ে:
    \begin{align*}
        a &= \frac{5 - 2}{5 + 2 + \frac{10}{2}} \times 9.8 \\
          &= \frac{3}{7 + 5} \times 9.8 = \frac{3}{12} \times 9.8 \\
          &= \frac{1}{4} \times 9.8 = 2.45\text{ ms}^{-2}
    \end{align*}
    
    \vspace{0.5em}
    \noindent\textbf{ধাপ ৩: টানের পার্থক্য নির্ণয়} \\
    এখন টানের পার্থক্যের সমীকরণে $a$-এর মান বসিয়ে পাই:
    \begin{align*}
        T_1 - T_2 &= \frac{1}{2} \times 10 \times 2.45 \\
                  &= 5 \times 2.45 \\
                  &= \mathbf{12.25}\text{ N}
    \end{align*}
\end{tcolorbox}

\newpage

% ------------------------------------------
% QUESTION SECTION 48
% ------------------------------------------
\begin{tcolorbox}[questionbox={প্রশ্ন (Question) 48}]
    \noindent
    $M = 1.2\text{ kg}$ ভর ও $L = 0.8\text{ m}$ দৈর্ঘ্যের একটি সুষম রডের ওপরের প্রান্ত একটি মসৃণ অনুভূমিক কব্জায় ঝোলানো আছে। রডটি উলম্বভাবে ঝুলন্ত থাকা অবস্থায় $m = 0.1\text{ kg}$ ভরের একটি অনুভূমিক বুলেট $v_0 = 80\text{ ms}^{-1}$ বেগে এসে রডের সর্বনিম্ন মুক্ত প্রান্তে আঘাত করে সম্পূর্ণ স্থিতিস্থাপকভাবে বিপরীত দিকে $20\text{ ms}^{-1}$ বেগে ফিরে গেল। আঘাতের তাৎক্ষণিক মুহূর্তে রডটির কৌণিক বেগ কত হবে?
    
    \vspace{0.8em}
    \begin{english}
    \noindent\textit{\color{darkgray}
    A uniform rod of mass $M = 1.2\text{ kg}$ and length $L = 0.8\text{ m}$ hangs vertically from a frictionless horizontal pivot at its top end. A horizontal bullet of mass $m = 0.1\text{ kg}$ moving at $v_0 = 80\text{ ms}^{-1}$ strikes the lowest tip of the rod and rebounds elastically in the opposite direction at $20\text{ ms}^{-1}$. What is the angular velocity of the rod immediately after the impact?
    }
    \end{english}
\end{tcolorbox}

% ------------------------------------------
% HIGH-QUALITY TIKZ DIAGRAM 48
% ------------------------------------------
\vspace{1em}
\begin{center}
\begin{tikzpicture}[>=Stealth, scale=1.3]
    % Pivot support
    \fill[pattern=north east lines, pattern color=gray!80] (-1, 3) rectangle (1, 3.2);
    \draw[thick] (-1, 3) -- (1, 3);
    \fill[black] (0, 3) circle (0.08);
    
    % Rod hanging vertically
    \draw[line width=4pt, brown!70!black] (0, 3) -- (0, 0.5);
    \node[right, font=\bfseries] at (0, 1.8) {$M = 1.2\text{ kg}, L = 0.8\text{ m}$};
    
    % Bullet approaching (Incoming)
    \draw[thick, fill=black] (-2.5, 0.5) rectangle (-1.8, 0.7);
    \node[above, font=\bfseries] at (-2.15, 0.7) {$m = 0.1\text{ kg}$};
    \draw[->, line width=1.5pt, primary] (-1.8, 0.6) -- (-0.2, 0.6) node[midway, above] {$v_0 = 80\text{ ms}^{-1}$};
    
    % Bullet rebounding (Outgoing)
    \draw[thick, fill=black, opacity=0.6] (-1.2, 0.1) rectangle (-0.5, 0.3);
    \draw[->, line width=1.5pt, red] (-0.5, 0.2) -- (-2.0, 0.2) node[midway, above] {$v_f = 20\text{ ms}^{-1}$};
    
    % Rod Rotation Arrow
    \draw[->, line width=1.2pt, blue] (0.5, 2.0) arc (0:-45:1.2) node[midway, right] {$\omega = ?$};
\end{tikzpicture}
\end{center}
\vspace{1em}

% ------------------------------------------
% OPTIONS & ANSWER 48
% ------------------------------------------
\begin{itemize}[label={}, leftmargin=1cm, itemsep=0.5em]
    \item[\textbf{A)}] কৌণিক বেগ $31.25\text{ rad s}^{-1}$
    \item[\textbf{B)}] কৌণিক বেগ $25.00\text{ rad s}^{-1}$
    \item[\textbf{C)}] কৌণিক বেগ $18.75\text{ rad s}^{-1}$
    \item[\textbf{D)}] কৌণিক বেগ $37.50\text{ rad s}^{-1}$
\end{itemize}

\vspace{0.5em}
\begin{tcolorbox}[answerbox]
    \textbf{\color{success} উত্তর:} A) কৌণিক বেগ $31.25\text{ rad s}^{-1}$
\end{tcolorbox}
\vspace{1em}

% ------------------------------------------
% SOLUTION SECTION 48
% ------------------------------------------
\begin{tcolorbox}[solutionbox={সমাধান (Solution)}]
    \noindent
    \textbf{ধাপ ১: জড়তার ভ্রামক ও আদি কৌণিক ভরবেগ} \\
    কব্জাবিন্দুর সাপেক্ষে সুষম রডের জড়তার ভ্রামক ($I$):
    \[ I = \frac{1}{3} M L^2 = \frac{1}{3} \times 1.2 \times (0.8)^2 = 0.4 \times 0.64 = 0.256\text{ kg}\cdot\text{m}^2 \]
    
    কব্জাবিন্দুর সাপেক্ষে বুলেটের আদি কৌণিক ভরবেগ ($L_i$):
    \[ L_i = m v_0 L = 0.1 \times 80 \times 0.8 = 6.4\text{ kg}\cdot\text{m}^2\text{s}^{-1} \]
    
    \vspace{0.5em}
    \noindent\textbf{ধাপ ২: সংঘর্ষের পর কৌণিক ভরবেগ সংরক্ষণ} \\
    সংঘর্ষের পর বুলেটের বেগ বিপরীতমুখী হওয়ায় শেষ কৌণিক ভরবেগ হবে ঋণাত্মক:
    \[ L_{\text{bullet}, f} = -m v_f L = - (0.1 \times 20 \times 0.8) = -1.6\text{ kg}\cdot\text{m}^2\text{s}^{-1} \]
    
    কৌণিক ভরবেগ সংরক্ষণ সূত্রানুসারে ($L_i = L_{\text{bullet}, f} + I \omega$):
    \begin{align*}
        6.4 &= -1.6 + 0.256 \omega \\
        0.256 \omega &= 6.4 + 1.6 = 8.0 \\
        \omega &= \frac{8.0}{0.256} = \mathbf{31.25}\text{ rad s}^{-1}
    \end{align*}
\end{tcolorbox}

\newpage

% ------------------------------------------
% QUESTION SECTION 49
% ------------------------------------------
\begin{tcolorbox}[questionbox={প্রশ্ন (Question) 49}]
    \noindent
    $\theta = 30^\circ$ নতিকোণবিশিষ্ট একটি অমসৃণ আনত তলে একটি নিরেট সিলিন্ডার না পিছলিয়ে বিশুদ্ধ ঘূর্ণনসহ গড়িয়ে পড়ার জন্য আনত তল ও সিলিন্ডারের মধ্যকার স্থিতি ঘর্ষণ গুণাঙ্ক $\mu_s$-এর ন্যূনতম মান কত হতে হবে?
    
    \vspace{0.8em}
    \begin{english}
    \noindent\textit{\color{darkgray}
    A solid cylinder rolls down a rough inclined plane of inclination $\theta = 30^\circ$ purely without slipping. What is the minimum coefficient of static friction $\mu_s$ required between the cylinder and the inclined plane to prevent slipping?
    }
    \end{english}
\end{tcolorbox}

% ------------------------------------------
% HIGH-QUALITY TIKZ DIAGRAM 49
% ------------------------------------------
\vspace{1em}
\begin{center}
\begin{tikzpicture}[>=Stealth, scale=1.3]
    \def\angle{30}
    
    % Incline
    \draw[thick, fill=gray!20] (0,0) -- (5,0) -- (0, {5*tan(\angle)}) -- cycle;
    \draw (4.2, 0) arc (180:{180-\angle}:0.8);
    \node at (3.6, 0.25) {$\theta = 30^\circ$};
    
    % Solid Cylinder on incline
    \coordinate (Center) at (2, {2*tan(\angle) + 0.6});
    \draw[thick, fill=cyan!30] (Center) circle (0.6);
    \fill[black] (Center) circle (0.04);
    
    % Forces
    \draw[->, line width=1.5pt, red] (Center) -- ++(0, -1.5) node[below] {$Mg$};
    \draw[->, line width=1.5pt, orange!80!black] (Center) ++({-0.6*sin(\angle)}, {-0.6*cos(\angle)}) -- ++({1.0*cos(\angle)}, {1.0*sin(\angle)}) node[above left] {$f_s$};
    
    % Direction of rolling
    \draw[->, thick, blue] (Center) ++(45:0.7) arc (45:-10:0.7) node[midway, right] {$\alpha$};
\end{tikzpicture}
\end{center}
\vspace{1em}

% ------------------------------------------
% OPTIONS & ANSWER 49
% ------------------------------------------
\begin{itemize}[label={}, leftmargin=1cm, itemsep=0.5em]
    \item[\textbf{A)}] ঘর্ষণ গুণাঙ্ক $\mu_s \ge 0.192$
    \item[\textbf{B)}] ঘর্ষণ গুণাঙ্ক $\mu_s \ge 0.289$
    \item[\textbf{C)}] ঘর্ষণ গুণাঙ্ক $\mu_s \ge 0.577$
    \item[\textbf{D)}] ঘর্ষণ গুণাঙ্ক $\mu_s \ge 0.144$
\end{itemize}

\vspace{0.5em}
\begin{tcolorbox}[answerbox]
    \textbf{\color{success} উত্তর:} A) ঘর্ষণ গুণাঙ্ক $\mu_s \ge 0.192$
\end{tcolorbox}
\vspace{1em}

% ------------------------------------------
% SOLUTION SECTION 49
% ------------------------------------------
\begin{tcolorbox}[solutionbox={সমাধান (Solution)}]
    \noindent
    \textbf{ধাপ ১: ন্যূনতম ঘর্ষণ গুণাঙ্কের সূত্র প্রতিপাদন} \\
    আনত তলে পিছলানো রোধ করে বিশুদ্ধ ঘূর্ণনের জন্য প্রয়োজনীয় ন্যূনতম স্থিতি ঘর্ষণ গুণাঙ্কের সাধারণ সূত্র:
    \[ \mu_s \ge \frac{\tan\theta}{1 + \frac{M R^2}{I_{\text{cm}}}} \]
    
    নিরেট সিলিন্ডারের ক্ষেত্রে $I_{\text{cm}} = \frac{1}{2} M R^2$, ফলে অনুপাতটি দাঁড়ায়:
    \[ \frac{M R^2}{I_{\text{cm}}} = \frac{M R^2}{\frac{1}{2} M R^2} = 2 \]
    
    \vspace{0.5em}
    \noindent\textbf{ধাপ ২: মান নির্ণয়} \\
    সূত্রে মানটি বসালে:
    \[ \mu_s \ge \frac{\tan\theta}{1 + 2} = \frac{1}{3} \tan\theta \]
    
    প্রদত্ত $\theta = 30^\circ$ বসিয়ে পাই:
    \begin{align*}
        \tan 30^\circ &= \frac{1}{\sqrt{3}} \approx 0.57735 \\
        \mu_s &\ge \frac{0.57735}{3} \\
              &\ge \mathbf{0.19245} \approx \mathbf{0.192}
    \end{align*}
\end{tcolorbox}

\newpage

% ------------------------------------------
% QUESTION SECTION 50
% ------------------------------------------
\begin{tcolorbox}[questionbox={প্রশ্ন (Question) 50}]
    \noindent
    অভিকর্ষহীন মুক্ত মহাশূন্যে একটি রকেটের জ্বালানিসহ মোট ভর $M_0 = 8000\text{ kg}$। রকেটের সাপেক্ষে নির্গত গ্যাসের বেগ $v_r = 2500\text{ ms}^{-1}$। রকেটটির দ্রুতি স্থিরাবস্থা থেকে $5.0\text{ kms}^{-1}$ এ উন্নীত করতে হলে রকেটের মোট আদি ভরের শতকরা কত অংশ জ্বালানি হিসেবে পুড়িয়ে নির্গত করতে হবে?
    
    \vspace{0.8em}
    \begin{english}
    \noindent\textit{\color{darkgray}
    A gravity-free deep space environment, a rocket has an initial total mass of $M_0 = 8000\text{ kg}$. The exhaust gas is ejected at a relative speed of $v_r = 2500\text{ ms}^{-1}$. To accelerate the rocket from rest to a speed of $5.0\text{ kms}^{-1}$, what percentage of the rocket's initial total mass must be consumed and expelled as fuel?
    }
    \end{english}
\end{tcolorbox}

% ------------------------------------------
% HIGH-QUALITY TIKZ DIAGRAM 50
% ------------------------------------------
\vspace{1em}
\begin{center}
\begin{tikzpicture}[>=Stealth, scale=1.3]
    % Background Space
    \fill[black!90, rounded corners=3pt] (-3.5, -1.2) rectangle (3.5, 1.2);
    
    % Rocket
    \fill[gray!30] (-1.2, -0.3) rectangle (0.8, 0.3);
    \fill[red!70!black] (0.8, -0.3) -- (1.5, 0) -- (0.8, 0.3) -- cycle;
    \fill[orange] (-1.2, 0) -- (-2.2, 0.3) -- (-1.8, 0) -- (-2.2, -0.3) -- cycle;
    
    % Vectors
    \draw[->, line width=1.5pt, cyan] (1.2, 0.5) -- (2.5, 0.5) node[right, font=\bfseries] {$5.0\text{ kms}^{-1}$};
    \draw[->, line width=1.2pt, orange] (-1.4, -0.6) -- (-2.6, -0.6) node[left, font=\bfseries] {$v_r = 2500\text{ ms}^{-1}$};
    
    % Initial Mass
    \node[white, font=\bfseries] at (-0.2, 0) {$M_0 = 8000\text{ kg}$};
\end{tikzpicture}
\end{center}
\vspace{1em}

% ------------------------------------------
% OPTIONS & ANSWER 50
% ------------------------------------------
\begin{itemize}[label={}, leftmargin=1cm, itemsep=0.5em]
    \item[\textbf{A)}] নির্গত ভরের শতকরা হার $86.47\%$
    \item[\textbf{B)}] নির্গত ভরের শতকরা হার $75.00\%$
    \item[\textbf{C)}] নির্গত ভরের শতকরা হার $63.21\%$
    \item[\textbf{D)}] নির্গত ভরের শতকরা হার $50.00\%$
\end{itemize}

\vspace{0.5em}
\begin{tcolorbox}[answerbox]
    \textbf{\color{success} উত্তর:} A) নির্গত ভরের শতকরা হার $86.47\%$
\end{tcolorbox}
\vspace{1em}

% ------------------------------------------
% SOLUTION SECTION 50
% ------------------------------------------
\begin{tcolorbox}[solutionbox={সমাধান (Solution)}]
    \noindent
    \textbf{ধাপ ১: তসিওলকোভস্কি রকেট সমীকরণ প্রয়োগ} \\
    মহাশূন্যে রকেটের গতির জন্য তসিওলকোভস্কি সমীকরণটি হলো:
    \[ \Delta v = v_r \ln\left(\frac{M_0}{M}\right) \]
    
    প্রদত্ত মানসমূহ: $\Delta v = 5.0\text{ kms}^{-1} = 5000\text{ ms}^{-1}$ এবং $v_r = 2500\text{ ms}^{-1}$।
    \begin{align*}
        5000 &= 2500 \ln\left(\frac{M_0}{M}\right) \\
        \ln\left(\frac{M_0}{M}\right) &= \frac{5000}{2500} = 2 \\
        \frac{M_0}{M} &= e^2 \approx 7.389
    \end{align*}
    
    \vspace{0.5em}
    \noindent\textbf{ধাপ ২: অবশিষ্ট ও ব্যয়িত ভরের হিসাব} \\
    অবশিষ্ট রকেটের ভর ($M$):
    \[ M = \frac{M_0}{e^2} = \frac{M_0}{7.389} \approx 0.1353 M_0 \]
    
    ব্যয়িত জ্বালানির ভর ($\Delta M$):
    \[ \Delta M = M_0 - M = M_0 (1 - 0.1353) = 0.8647 M_0 \]
    
    অতএব, নির্গত ভরের শতকরা হার:
    \[ \frac{\Delta M}{M_0} \times 100\% = \mathbf{86.47\%} \]
\end{tcolorbox}

\end{document}